\documentclass[11pt]{article}
\usepackage[T1]{fontenc}
\usepackage[utf8]{inputenc}
\usepackage{lmodern,microtype,amsmath,amssymb,amsthm,mathtools}
\usepackage[margin=1in]{geometry}
\usepackage{booktabs,longtable,array,enumitem}
\usepackage[colorlinks=true,linkcolor=blue,citecolor=blue,urlcolor=blue]{hyperref}
\usepackage{fancyhdr}
\setlength{\headheight}{14pt}
\pagestyle{fancy}\fancyhf{}
\fancyhead[L]{\small BR\_PENCIL}
\fancyhead[R]{\small Residual subspaces and border rank}
\fancyfoot[C]{\thepage}
\setlist{nosep}\setlength{\emergencystretch}{2em}
\newcommand{\R}{\mathbb R}\newcommand{\C}{\mathbb C}
\newcommand{\Z}{\mathbb Z}\newcommand{\Q}{\mathbb Q}
\newcommand{\F}{\mathbb F}\newcommand{\A}{\mathbb A}
\newcommand{\BR}{\underline{\mathrm R}}\newcommand{\TR}{\mathrm R}
\newcommand{\BSR}{\underline{\mathrm R}_{\mathrm{sym}}}
\newcommand{\GL}{\mathrm{GL}}\newcommand{\SL}{\mathrm{SL}}
\newcommand{\Sym}{\operatorname{Sym}}\newcommand{\Spec}{\operatorname{Spec}}
\newcommand{\rank}{\operatorname{rank}}\newcommand{\sgn}{\operatorname{sgn}}
\newcommand{\diag}{\operatorname{diag}}\newcommand{\tr}{\operatorname{Tr}}
\newcommand{\id}{\operatorname{id}}\newcommand{\Span}{\operatorname{span}}
\newcommand{\Gr}{\operatorname{Gr}}\newcommand{\Hilb}{\operatorname{Hilb}}
\newcommand{\coker}{\operatorname{coker}}\newcommand{\calO}{\mathcal O}
\newcommand{\CP}{\mathbb P}\newcommand{\eps}{\varepsilon}
\newcommand{\ceil}[1]{\left\lceil#1\right\rceil}
\newcommand{\abs}[1]{\lvert#1\rvert}
\newenvironment{finding}[2]{\par\medskip\noindent\textbf{Finding #1 (#2).}\ \itshape}{\par\medskip}
\newenvironment{background}[2]{\par\medskip\noindent\textbf{Background #1 (#2).}\ }{\par\medskip}
\title{\textbf{Border rank of the $4\times4$ determinant and permanent\\by multidimensional residual substitution}\\[5pt]
\large Exact Grassmannian minima, Koszul and prolongation certificates,\\
all-field models, and the Grothendieck dictionary}
\author{Research report for BR\_PENCIL\\\small Run BR\_PENCIL-20261121-R1}
\date{21 November 2026\\\small Local date: Europe/Berlin}
\begin{document}
\maketitle
\begin{abstract}
We study the border rank of tensors obtained by retaining several
 dimensions of a determinant or permanent flattening image. An explicit
 torus degeneration reduces every Grassmannian minimum to a finite
 minimum on coordinate subspaces. For $n=4,b=2,j=2$, the exact minimum
 is $3$ for each row family and $6$ for each polynomial family. The
 polynomial calculation yields border rank at least $40$. Going further,
 the polynomial residual minima for $j=3,4,5,6$ are respectively $8,9,11,12$,
 giving the all-field bounds $\BR(\mathcal D_4)\geq42$ and
 $\BR(\mathcal P_4)\geq42$ in the integral coefficient-tensor model.
 In characteristic zero, complete finite covers with Koszul and
 prolongation certificates prove that the seventh, eighth and ninth
 minima are $14,15,16$ for both polynomial families and every $n\geq4$. Consequently
 $\BR(\mathcal D_4)\geq43$ and $\BR(\mathcal P_4)\geq43$;
 these CP bounds also lower-bound the usual complex symmetric
 border ranks. Explicit degenerations attain the stated residual
 minima. We distinguish the integral divided-power model from the
 ordinary-power convention in small characteristic. Separately,
 we give a rational, zero-gap DPS certificate for a residual image
 outside the four families and a quantitative hierarchy error bound.
 The proposed maximization of the projector overlap over a full
 general-linear orbit is refuted: its supremum is one, generally
 unattained. A separate part translates each construction into schemes,
 locally free sheaves and the appropriate Grothendieck foundations,
 without attributing the modern rank results to those texts.
\end{abstract}
\clearpage
{\small\tableofcontents}
\newpage

\section*{Conventions, provenance, and scope}
\addcontentsline{toc}{section}{Conventions, provenance, and scope}
The ordinary CP border rank allows independent vectors in every
tensor factor. Symmetric border rank allows only powers. Thus a
CP lower bound for a symmetric tensor also bounds its symmetric
border rank; the reverse implication for lower bounds is not used.
Row tensors $D_n,P_n\in(k^n)^{\otimes n}$ are different from
polynomial coefficient tensors in $(k^{n^2})^{\otimes n}$.
The displayed bounds $43$ for both polynomial families in characteristic zero
and $42$ for both over every algebraically closed field
are lower bounds for the latter tensors. They are not assertions
that the full polynomial border ranks equal the displayed numbers.

All-field statements use integer ordered-matching coefficients.
In characteristic zero these are nonzero scalar multiples of the
normalized coefficient tensors of the determinant and permanent.
In small characteristic one must use the divided-power model;
the ordinary polynomial power map is a different construction.
The precise formulas and the factorial obstruction are given in
the fields section. Throughout, $\mathcal D_n=S_{\det,n}$ and
$\mathcal P_n=S_{\mathrm{per},n}$ denote these integral polynomial
coefficient tensors. Over characteristic zero the nonzero scalar
normalization does not affect either border rank.

A finding is a deduction established in this run, not a claim of
historical priority. Known methods and comparison results retain
their source locators. The predecessor is GCT\_GINEQ,
run GCT\_GINEQ-2026-11-20-R1, findings F-13, F-18, F-19,
F-20, F-21, F-22 and F-23; the attached intake of
21 November 2026, \S5.1(b), is the source of the main question.
No result of Alp\"oge or Masullo is used, so no irrelevant
step, Lean-file or statement locator is invented.

The model is GPT-6 Astra Pro, with inherited
reasoning configuration; no separate numerical effort setting is
exposed. Whole-proof reads, independent re-proofs and exact
computations are distinguished in \texttt{CHECK\_EVENTS.tsv}.
The certificates are ordinary mathematical and exact-arithmetic
certificates, not formal theorem-prover proofs or external refereeing.

\part{Checked main findings and certificates}
\section{The all-field multidimensional substitution statement}

For a finite-dimensional space $W=V_1\otimes\cdots\otimes V_b$
and a $j$-plane $L\subset W$, write
\[
 \mathsf P_L=\sum_{i=1}^j a_i\otimes u_i
       \in k^j\otimes V_1\otimes\cdots\otimes V_b,
\]
where $(u_i)$ is any basis of $L$. Its CP border rank is independent
of the basis. All the $b$ residual factors remain separate.

\begin{finding}{R-01}{proper incidence and substitution of $j$ dimensions}
Let $k$ be algebraically closed, let
$T\in A_0\otimes V_1\otimes\cdots\otimes V_b$ be nonzero,
and let $U$ be its $A_0$-flattening image, of dimension $q$.
For $1\leq j\leq q$,
\begin{equation}\label{eq:substitution}
 \BR_k(T)\ \geq\ q-j+
       \min_{L\in\Gr(j,U)}\BR_k(\mathsf P_L).
\end{equation}
The same inequality lower-bounds an original tensor from which $T$
was obtained by grouping factors into $A_0$.
\end{finding}
\begin{proof}
For $s\geq1$ let $Z_{j,s}\subset\Gr(j,W)$ be the locus
$\BR(\mathsf P_L)\leq s$. It is closed. Indeed, on each
standard Grassmannian chart choose the regular echelon frame of
the tautological bundle. The associated tensor is a regular
map into $k^j\otimes W$. The inverse image of its closed
Segre-secant cone is closed in that chart. Changing the frame
acts through $\GL_j$, which preserves the cone, so these local
conditions agree and define the claimed closed locus.

In $\Gr(j,W)\times\Gr(q,W)$ the flag incidence $L\subset U'$
is closed: it is the vanishing of the map from the rank-$j$
tautological subbundle to the rank-$(\dim W-q)$ universal
quotient. Therefore
\[
 E_{q,j,s}=\{U':\text{some }L\subset U',\ \dim L=j,
                         \BR(\mathsf P_L)\leq s\}
\]
is closed. It is the image of this incidence intersected with
$Z_{j,s}\times\Gr(q,W)$ under the projective projection.
The precise properness inputs are EGA~II, Definition~5.4.1,
Proposition~5.4.2 and Theorem~5.5.3(i); see the dictionary.

Replace $A_0$ by the $q$-dimensional support of $T$ in that
factor. An inclusion and a left-inverse projection show that
this replacement preserves its border rank. On the open set
$O$ of tensors whose first flattening has full rank $q$, the
image map $f:O\to\Gr(q,W)$ is regular: its coordinates on
a nonzero-minor chart are the usual echelon-frame coordinates.
Put $r=\BR(T)$; an ordinary flattening gives $r\geq q$.

For an actual rank-at-most-$r$ tensor $T'\in O$, let $S$ be
the span of the pure residual atoms of a length-at-most-$r$
decomposition. Then $U'=f(T')\subset S$ and
$q\leq d:=\dim S\leq r$. Choose a basis of $S$ from those
pure atoms and let $M$ be the span of any $d-q+j$ of them.
The dimension inequality gives $\dim(U'\cap M)\geq j$.
Choose a $j$-plane $L\subset U'\cap M$. Expanding its
basis in the selected atoms expresses $\mathsf P_L$ as a
sum of at most $d-q+j\leq r-q+j$ pure tensors, including
its new first factor. Thus $f(T')\in E_{q,j,r-q+j}$.

The set $f^{-1}(E_{q,j,r-q+j})$ is closed in $O$. Since
$T$ is in the closure in $O$ of the actual rank-at-most-$r$
tensors in $O$, it also belongs to that inverse image. Hence
some $L\subset U$ has $\BR(\mathsf P_L)\leq r-q+j$.
Rearranging proves \eqref{eq:substitution}. Grouping original
factors maps every original pure tensor to a pure tensor in
the grouped format and hence cannot increase border rank;
the last assertion follows.
\end{proof}

The multidimensional three-factor principle is already
Landsberg--Micha{\l}ek, Proposition~2.3\cite{LM2018}.
The proof above records the precise arbitrary residual-order,
arbitrary-algebraically-closed-field form needed here. Its
$j=1$ case is the proper-incidence version of predecessor
GCT\_GINEQ F-13/F-18; no new attribution of the substitution
principle is intended. The new numerical work is the residual
minimum and its certificates.

\paragraph{Geometric form.}
Each $L$ is a point of the Grassmannian with its tautological
rank-$j$ locally free sheaf. Its inclusion into the tensor
bundle is the tautological residual tensor. Border rank forces
the image point $[U]$ into the proper image $E_{q,j,s}$ of a
closed flag incidence. The proof is carried out in each
algebraically closed fiber. It does not assume that an
integral secant ideal commutes with every reduction modulo a
prime, and it uses no Euclidean compactness or Hermitian norm.

% Checked mathematical fragments from /root/pencil.
% Dependencies: finding environment with two arguments; amsmath, amsthm,
% amssymb; macros \BR and \Span. No division by a factorial is used.
% The final P-08 here is the finite-difference upper certificate:
% its source-note locator is PENCIL_PROOFS.md, P-09 (upper-bound part).

\section{Exact two-dimensional residual minima}

Throughout this section the field is algebraically closed and the tensors
have integral ordered-coordinate coefficients \(1\) or \(-1\). For the
polynomial families these are the divided-power coefficient tensors.
In characteristic zero they differ from the usual symmetric coefficient
tensors only by a nonzero scalar. In positive characteristic the integral
tensor, rather than evaluation of a symmetric tensor as a polynomial, is
the convention.

For a \(j\)-plane \(W\subset U_{n,b}\), choose a basis \(u_1,\ldots,u_j\)
and put \(P_W=\sum_{h=1}^j a_h\otimes u_h\). Its border rank is independent
of the basis. Write
\[
 \mu_j(U_{n,b})=\min_{W\in\operatorname{Gr}(j,U_{n,b})}\BR(P_W).
\]
The residual factors of \(P_W\) remain ungrouped.

\begin{finding}{P-01}{the Grassmannian minimum is a coordinate minimum}
In each of the four families and for every allowed \(n,b,j\),
\[
 \mu_j(U_{n,b})=
 \min_{\substack{\mathcal S\subset\mathcal B_{n,b}\\|\mathcal S|=j}}
 \BR\left(\sum_{Q\in\mathcal S}a_Q\otimes Q\right),
\]
where \(\mathcal B_{n,b}\) is the coordinate residual basis: \(Q_I\)
in a row family and \(Q_{I,J}\) in a polynomial family.
\end{finding}
\begin{proof}
For a row family, act on \(e_i\), \(0\leq i<n\), by
\(g(t)e_i=t^{2^i}e_i\). The basis tensor \(Q_I\) is an eigenvector
of weight \(\sum_{i\in I}2^i\), and these weights are distinct.
For a polynomial family act on \(e_{ij}\) by
\(t^{2^i+2^{n+j}}\). Then \(Q_{I,J}\) has weight
\(\sum_{i\in I}2^i+\sum_{j\in J}2^{n+j}\), again with distinct
weights. Integer exponents define distinct characters in every
characteristic.

Order the residual basis \(q_1,\ldots,q_N\) by increasing weights
\(w_1<\cdots<w_N\). Row reduction gives a basis of any chosen plane
of the form
\[
 u_h=q_{p_h}+\sum_{i>p_h}c_{hi}q_i,\qquad h=1,\ldots,j,
\]
with distinct pivot indices. Consequently
\[
 u_h(t):=t^{-w_{p_h}}g(t)^{\otimes b}u_h
 =q_{p_h}+\sum_{i>p_h}c_{hi}t^{w_i-w_{p_h}}q_i
\]
is polynomial in \(t\), and \(u_h(0)=q_{p_h}\).
For \(t\ne0\), the tensor \(\sum_h a_h\otimes u_h(t)\) is
obtained from \(P_W\) by invertible maps in its factors: \(g(t)\)
in each residual factor and a diagonal map in the first factor.
Its border rank therefore equals \(\BR(P_W)\).
The border-rank-at-most-\(r\) locus is closed, so specialization
at zero cannot increase border rank. This gives a coordinate plane
of no larger border rank. The converse inequality follows because
every coordinate plane is an allowed plane.
\end{proof}

\begin{finding}{P-02}{all two-dimensional minima reduce to an adjacent pair}
For \(n\geq b+1\), the minimum \(\mu_2(U_{n,b})\) is the border rank
of the following single tensor, independently of \(n\):
\[
 \begin{cases}
 a_1\otimes Q_{[b]}+
 a_2\otimes Q_{[b-1]\cup\{b+1\}},&\text{row families},\\[2mm]
 a_1\otimes Q_{[b],[b]}+
 a_2\otimes Q_{[b],[b-1]\cup\{b+1\}},&\text{polynomial families}.
 \end{cases}
\]
This is an exact smaller-tensor reduction; the displayed border
rank for arbitrary \(b>2\) is not evaluated here.
\end{finding}
\begin{proof}
By P-01 it suffices to consider coordinate pairs.
For two distinct \(b\)-subsets \(I,K\), write their common part as
\(C\), with \(|C|=c\), and their equally sized differences as
\(\{i_1,\ldots,i_d\}\) and \(\{k_1,\ldots,k_d\}\), where
\(d=b-c\geq1\). Map the common coordinates bijectively onto
\(e_1,\ldots,e_c\); map both \(e_{i_h}\) and \(e_{k_h}\) to
\(e_{c+h}\) for \(h<d\); and map \(e_{i_d},e_{k_d}\) to
\(e_b,e_{b+1}\), respectively. Send unused coordinates to zero.
The map is injective on each of the two selected coordinate sets.
Its tensor power therefore takes the two selected row tensors to
the adjacent pair, up to signs in the determinant case. Those
signs are absorbed in the first factor.

For polynomial tensors, suppose first that the column sets \(J,L\)
are distinct. Map the two row sets \(I,K\) injectively onto the
same \(b\)-element row set, by matching their differences and
retaining their intersection. Map the column sets \(J,L\) to the
adjacent pair just constructed. The induced linear map on matrix
coordinates, \(e_{ij}\mapsto e_{f(i),g(j)}\), is injective on the
rows and columns within each selected term. It therefore takes
each selected minor or permanent to the desired adjacent minor
or permanent, up to a determinant sign. If \(J=L\), then
\(I\ne K\); interchange rows and columns first, and transpose
afterwards.

Products of linear maps cannot increase border rank. Thus every
coordinate pair has border rank at least that of the adjacent
pair. The adjacent pair occurs in \(U_{n,b}\), proving equality.
\end{proof}

\begin{finding}{P-03}{a uniform star degeneration}
Let \(k\geq1\), let \(a_1,\ldots,a_k\) and \(e_0,\ldots,e_k\)
be independent, and let \(\delta\in\{1,-1\}\). Then
\[
 S_{k,\delta}:=\sum_{i=1}^k
 a_i\otimes(e_0\otimes e_i+\delta e_i\otimes e_0)
 \qquad\text{satisfies}\qquad
 \BR(S_{k,\delta})=k+1
\]
over every algebraically closed field.
\end{finding}
\begin{proof}
Contraction in the last factor gives the \(k+1\) independent
vectors
\(\delta\sum_i a_i\otimes e_i\) and \(a_i\otimes e_0\),
\(1\leq i\leq k\). Thus the corresponding flattening has rank
\(k+1\), proving the lower bound.
For \(t\ne0\), the following is a sum of \(k+1\) pure tensors:
\[
 \frac1t\left[
 \sum_{i=1}^k a_i\otimes(e_0+\delta t e_i)\otimes(e_0+t e_i)
 -\left(\sum_{i=1}^k a_i\right)\otimes e_0\otimes e_0
 \right].
\]
Its exact expansion is
\[
 S_{k,\delta}+\delta t\sum_{i=1}^k a_i\otimes e_i\otimes e_i.
\]
It therefore specializes to \(S_{k,\delta}\). There is no
division by an integer; the identity remains valid in
characteristic two.
\end{proof}

\begin{finding}{P-04}{the requested \(n=4,b=2,j=2\) computation}
For every \(n\geq3\), the exact two-dimensional residual minima are
\[
 \begin{array}{c|cccc}
 \text{family}&D_n&P_n&\mathcal D_n&\mathcal P_n\\ \hline
 \mu_2(U_{n,2})&3&3&6&6 .
 \end{array}
\]
In particular, at \(n=4\) the two-dimensional substitution
bound is \(7\) in the row families and \(40\) in the polynomial
families. These assertions hold for the integral coefficient
tensors over every algebraically closed field.
\end{finding}
\begin{proof}
The adjacent row pair in P-02 is \(S_{2,-1}\) for the
determinant and \(S_{2,1}\) for the permanent, up to labels and
signs. P-03 gives its exact border rank three.

For the polynomial pair use the six matrix-coordinate variables
\[
 a=e_{11},\ b=e_{12},\ c=e_{13},\
 d=e_{21},\ e=e_{22},\ f=e_{23}.
\]
Its two slices are the ordered coefficient tensors of
\(ae+\sigma bd\) and \(af+\sigma cd\), where \(\sigma=-1\)
for the determinant and \(+1\) for the permanent.
The associated tensor is the sum of a symmetric star
\(S_{2,1}\) on \((a,e,f)\) and \(\sigma S_{2,1}\) on
\((d,b,c)\), with the same two-dimensional first factor.
Adding the two explicit curves of P-03 gives a six-term
degeneration.

For the matching lower bound, each coordinate quadratic is a
nonsingular signed matching on its four-variable coordinate
support. Its matrix kernel is precisely the coordinate
complement of that support. For a distinct pair
\((I,J),(K,L)\), the common slice kernel is therefore the
complement of
\((I\times J)\cup(K\times L)\). The pencil flattening has rank
\[
 |(I\times J)\cup(K\times L)|
   =8-|I\cap K|\,|J\cap L|\geq6.
\]
The last inequality uses distinctness of the pairs; the only
excluded product larger than two is \(2\cdot2\), which would
make the pairs identical. The adjacent pair attains six.
P-01 proves uniformity over the entire Grassmannian, not only
coordinate planes.

Finally the generalized substitution theorem R-01 gives
\(\BR(T_n)\geq\binom n2+1\) for the row families and
\(\BR(T_n)\geq\binom n2^2+4\) for the polynomial families.
At \(n=4\) these are the stated numbers.
\end{proof}

\begin{finding}{P-05}{exact Koszul certificates for the pencils}
For a pencil \(P=a_1\otimes Q_1+a_2\otimes Q_2\) in
\(A\otimes B\otimes C\), with \(\dim A=2\), the \(p=1\)
Koszul map
\[
 K_P:A\otimes B^*\longrightarrow\bigwedge^2A\otimes C,
 \qquad
 a\otimes\beta\longmapsto
 \sum_{i=1}^2a_i\wedge a\otimes Q_i(\beta)
\]
has matrix \([-Q_2\mid Q_1]\), up to an overall sign convention.
Its pure-tensor cost is \(\binom11=1\).
For the coordinate row pencils its rank is the union size of
the two edges; for the coordinate polynomial pencils its rank
is \(8-|I\cap K||J\cap L|\). Thus P-04 is certified by
Koszul flattenings with exact costs.
\end{finding}
\begin{proof}
For a pure tensor the displayed map has rank at most one.
Rank subadditivity and closed determinantal rank loci give
\(\BR(P)\geq\operatorname{rank}K_P\).
The image of each slice is its coordinate support, proving
the stated concatenated ranks. More explicitly, for each
variable in the support union choose a slice containing it
and the unique input coordinate paired with it in that
signed matching. These are distinct columns, each with one
nonzero entry in the chosen rows, and give a square minor
with determinant \(1\) or \(-1\).

At \(n=4\), exhaustive integer certificates give \(15\) row
pairs, with rank distribution \(12\) of rank three and \(3\)
of rank four, and \(630\) polynomial pairs, with \(144\) of
rank six, \(288\) of rank seven, and \(198\) of rank eight.
The checker \path{check_pencils.py} constructs the minors
and verifies both signs, as well as the star degenerations.
The two-dimensional Koszul map is an ordinary flattening
after a basis change; no stronger Koszul estimate is needed
for this exact minimum.
\end{proof}

\section{Larger residual subspaces}

\begin{finding}{P-06}{anchored collections have explicit exact border rank}
Let \(H\subset\{1,\ldots,r\}\times\{1,\ldots,s\}\) use every
row and column vertex, let \(h=|H|>0\), and let
\(a,x_i,y_j,z_{ij}\), \((i,j)\in H\), be independent.
For \(\sigma=\pm1\), define
\[
 T_H=\sum_{(i,j)\in H}c_{ij}\otimes
 \bigl(a\otimes z_{ij}+z_{ij}\otimes a
       +\sigma(x_i\otimes y_j+y_j\otimes x_i)\bigr).
\]
Then \(\BR(T_H)=h+r+s+1\), in every characteristic.
\end{finding}
\begin{proof}
Every slice is a nonsingular signed matching on its
four-variable support. The union contains all
\(h+r+s+1\) variables, so the last-factor flattening has
that rank.

For the upper bound put
\[
 \begin{split}
 v_{00}&=a,\qquad v_{i0}=a+t x_i,\qquad v_{0j}=a+t y_j,\\
 v_{ij}&=a+t x_i+t y_j+\sigma t^2z_{ij}.
 \end{split}
\]
The curve
\[
 T(t)=\sigma t^{-2}\sum_{(i,j)\in H}c_{ij}\otimes
 \bigl(v_{ij}^{\otimes2}-v_{i0}^{\otimes2}
       -v_{0j}^{\otimes2}+v_{00}^{\otimes2}\bigr)
\]
is polynomial after cancellation and has value \(T_H\) at
zero. Indeed, its constant term after division by \(t^2\)
is exactly the displayed slice. Grouping by the \(h+r+s+1\)
vectors \(v_{ij},v_{i0},v_{0j},v_{00}\) expresses it as at
most that many pure tensors for \(t\ne0\).
The identity is over \(\mathbb Z\).
\end{proof}

\begin{finding}{P-07}{determinant coordinate minima reduce to order ideals}
For either determinant family, P-01 can be restricted to
coordinate subsets that are lower order ideals of cardinality
\(j\). The order is componentwise order on increasing
\(b\)-subsets in the row family, and product componentwise
order on pairs of increasing \(b\)-subsets in the polynomial
family. The assertion holds in every characteristic.
\end{finding}
\begin{proof}
Fix \(v>u\) and apply the transvection
\(e_v\mapsto e_v+t e_u\). In the exterior-power coordinate
basis, it sends
\[
 Q_I\longmapsto Q_I+\varepsilon_I t
 Q_{(I\setminus\{v\})\cup\{u\}}
\]
when \(v\in I\), \(u\notin I\), and fixes the vector
otherwise. If both indices occur, the added exterior
product has a repeated coordinate and vanishes. Thus the
nontrivial actions form disjoint two-element chains
with fixed lower partner. For polynomial minors apply
the same construction separately to row or column
exterior powers.

If both partners of a chain are selected, their span is
unchanged. If only the lower partner is selected, it is
fixed. If only the upper partner is selected, substitute
\(t=s^{-1}\) and normalize its frame vector to
\(sQ_{\rm upper}+\varepsilon Q_{\rm lower}\); it tends to
the lower partner. Different chains have different
partners, so the limiting frame remains independent.
For \(s\ne0\) the residual tensors differ only by product
general linear transformations and changes of first-factor
basis. Border rank cannot increase at the limit.

Whenever a replacement occurs, the sum of all selected
integer coordinate indices strictly decreases. Iterating
must terminate. At termination the selected subset is
closed under replacement of any index by a smaller absent
index. Such replacements generate componentwise order:
to pass from \(I\) to \(J\leq I\), successively replace the
first unequal coordinate by that of \(J\). The polynomial
case gives the product order by doing rows and columns
separately. The resulting coordinate plane is therefore
an order ideal and has no larger border rank.
\end{proof}

The exterior-power cancellation in P-07 is essential.
The analogous permanent image is not invariant under
these transvections, so P-07 is not asserted for it.

\begin{finding}{P-08}{finite differences for the exceptional permanent subspaces}
Let \(T_5\) have the five ordered coefficient slices
\[
 ae+bd,\quad af+cd,\quad bf+ce,\quad ah+bg,\quad ai+cg
\]
in the row-major variables
\(\left(\begin{smallmatrix}a&b&c\\d&e&f\\g&h&i\end{smallmatrix}\right)\).
Then \(\BR(T_5)\leq11\).
Adding the slice \(bi+ch\) gives \(T_6\) with
\(\BR(T_6)\leq12\).
Both upper bounds have integral, explicit finite-difference
certificates. In conjunction with R-04, \(\BR(T_5)=11\).
\end{finding}
\begin{proof}
Put
\[
 v(u,w,y,z)=(r_0+u r_1+w r_2)\otimes(s_0+y s_1+z s_2),
 \qquad F=v\otimes v.
\]
The coefficients of \(uy,uz,uyz,wy,wz\) in \(F\) are
exactly the five stated ordered symmetric quadrics.
For example, the coefficient of \(uy\) is
\(a\otimes e+e\otimes a+b\otimes d+d\otimes b\);
that of \(uyz\) is
\(b\otimes f+f\otimes b+c\otimes e+e\otimes c\).

For a set \(S\) of parameters define
\[
 D_S(t)=t^{-|S|}
 \sum_{R\subset S}(-1)^{|S|-|R|}F(t\,1_R).
\]
Inclusion-exclusion kills every monomial with support
properly contained in \(S\); parameters outside \(S\)
are zero. After cancellation \(D_S(t)\) is polynomial
and its value at zero is the coefficient of the
squarefree monomial supported on \(S\).

Take the first-factor-labelled sum of these differences
for \(S=uy,uz,uyz,wy,wz\).
The distinct evaluation sets are the eight subsets of
\(\{u,y,z\}\) and the three additional sets
\(\{w\},\{w,y\},\{w,z\}\). Grouping their first-factor
coefficients yields eleven pure tensors, with the desired
limit. For \(T_6\), the additional set \(S=wyz\) adds only
the evaluation set \(\{w,y,z\}\), giving twelve terms.
The checker \path{check_finite_differences.py} expands
the Laurent expressions over \(\mathbb Z\), verifies
cancellation of all negative powers, and checks every
constant tensor coefficient.
\end{proof}

\subsection*{Scheme and representation dictionary for P-01--P-08}
The row and column diagonal actions are split tori inside
general linear group schemes over \(\mathbb Z\). In P-01 the
echelon frame is a rank-\(j\) subbundle of the trivial residual
bundle on \(\mathbb A^1\), because its pivot coefficient
submatrix remains the identity. Its tautological tensor
specializes in the closed Segre secant cone.
P-02 uses morphisms of free modules, hence of their tensor
bundles. P-03, P-06 and P-08 supply integral curves in
secant cones. P-05 is an exterior-power morphism of locally
free sheaves; its unit minors give the asserted
determinantal inequalities after every base change.
In P-07 the elementary lowering root subgroups and the
diagonal torus generate a split Borel; the terminal
coordinate planes are its invariant subspaces.
This is an explicit form of the projective Borel-fixed
reduction, requiring no semisimplicity in positive
characteristic. The exact EGA, SGA and FGA locators for
these constructions are supplied in the dictionary
section of the main paper.

% P-10/P-11, all-field finite differences and the cubic adjacent pair.
% Dependencies: finding environment, \BR, P-02, R-01, integral convention.
\section{Common-core permanent subspaces and cubic row pencils}

\begin{finding}{P-10}{an integral common-core permanent degeneration}
Let $b\geq2$, $j\geq1$, and let
$e_0,e_1,\ldots,e_{b-2},f_1,\ldots,f_j$ be independent.
Write $Q_i$ for the ordered permanent tensor on
$e_0,e_1,\ldots,e_{b-2},f_i$, and let $a_1,\ldots,a_j$
be independent first-factor vectors. Over every algebraically
closed field,
\[
 \BR\left(\sum_{i=1}^j a_i\otimes Q_i\right)
 \leq (j+1)2^{b-2}.
\]
In particular $\BR(P_b)\leq2^{b-1}$ in every characteristic.
\end{finding}
\begin{proof}
Let $C=\{e_1,\ldots,e_{b-2}\}$, put $S_i=C\cup\{f_i\}$,
and define the tensor-valued polynomial
\[
 F_R(t)=\left(e_0+t\sum_{v\in R}v\right)^{\otimes b},
 \qquad
 D_i(t)=t^{-(b-1)}\sum_{R\subseteq S_i}
                 (-1)^{b-1-|R|}F_R(t).
\]
Inclusion--exclusion kills every monomial missing one
of the $b-1$ indicated directions. Every surviving
monomial has parameter degree at least $b-1$, so the
Laurent expression is polynomial. At degree exactly
$b-1$, every direction occurs once and exactly one
slot contains $e_0$. Its coefficient is therefore
the ordered tensor $Q_i$, with each tensor-basis word
occurring once. Thus $D_i(0)=Q_i$ over $\mathbb Z$.

In $\sum_i a_i\otimes D_i(t)$ group equal evaluation
sets $R$. They are the $2^{b-2}$ subsets of $C$ and,
for each $i$, the $2^{b-2}$ subsets of $C$ with $f_i$
adjoined. Hence there are $(j+1)2^{b-2}$ pure tensor
terms for $t\ne0$, and the limit is the claimed tensor.
No integer is divided out, so the proof survives every
characteristic. The final assertion is the case $j=1$.
\end{proof}

\begin{finding}{P-11}{exact cubic row pencil minima}
For independent $x,y,z,w$ and first-factor basis $a_0,a_1$,
let
\[
 T_+=a_0\otimes\operatorname{per}(x,y,z)
       +a_1\otimes\operatorname{per}(x,y,w),\qquad
 T_-=a_0\otimes\operatorname{alt}(x,y,z)
       +a_1\otimes\operatorname{alt}(x,y,w),
\]
where $\operatorname{per}$ is the sum of all six ordered
words and $\operatorname{alt}$ is their alternating sum.
Then
\[
 \BR(T_+)=6\quad\hbox{in every characteristic},
 \qquad
 \BR(T_-)=8\quad\hbox{if }\operatorname{char}k\ne2.
\]
In characteristic two $T_-=T_+$ and its border rank is six.
Consequently, for every $n\geq4$ the row permanent has
$\mu_2(U_{n,3})=6$; the row determinant has the same
minimum in characteristic two, and minimum eight
in other characteristics. Moreover,
\[
 \BR(P_4)=8\quad\hbox{over every algebraically closed field},
 \qquad \BR(D_4)=8\quad\hbox{in characteristic two}.
\]
\end{finding}
\begin{proof}
The permanent upper bound is P-10 with $b=3,j=2$.
Concretely, the two mixed differences use only the six
evaluation vectors
\[
 x,\quad x+ty,\quad x+tz,\quad x+tw,\quad
 x+ty+tz,\quad x+ty+tw.
\]
Divide the $\{y,z\}$ and $\{y,w\}$ differences of
their third tensor powers by $t^2$, and label them by
$a_0,a_1$. Their constant tensor coefficients are
exactly $T_+$.

For the lower bounds put $A=\langle a_0,a_1\rangle$,
$E=\langle x,y,z,w\rangle$, and use the double-Koszul map
\[
 K_T:A\otimes E\otimes E^*
       \longrightarrow\Lambda^2A\otimes\Lambda^2E\otimes E.
\]
For $T=\sum t_{ijkl}a_i\otimes e_j\otimes e_k\otimes e_l$,
its value on $\alpha\otimes u\otimes\beta$ is
\[
 \sum t_{ijkl}(a_i\wedge\alpha)\otimes(e_j\wedge u)
                   \otimes\beta(e_k)e_l.
\]
Its pure-tensor cost is exactly three: wedge by a nonzero
vector has rank one on the two-dimensional factor and
rank three on the four-dimensional factor, while evaluation
and the remaining fixed vector each have rank one.
Closed determinantal loci therefore give
$\BR(T)\geq\lceil\operatorname{rank}K_T/3\rceil$.

Use zero-based indices, the ordered basis $(x,y,z,w)$,
lexicographically ordered exterior pairs, and product
bases in the displayed factor order. Thus the matrix
has $24$ rows and $32$ columns. For $T_+$ the rows
\[
 (3,8,0,17,9,1,15,7,5,19,11,12,4,13,14,6)
\]
and columns
\[
 (0,1,3,4,5,7,8,9,11,12,13,16,17,20,24,25)
\]
give an order-$16$ minor of determinant $1$.
For $T_-$ the rows
\[
 (3,8,0,17,9,1,15,7,5,19,11,16,12,4,13,2,14,6,18)
\]
and columns
\[
 (0,1,3,4,5,7,8,9,11,12,13,15,16,17,20,21,24,25,28)
\]
give an order-$19$ minor of determinant $4$.
The exact integer matrices and these determinants are
reconstructed from the displayed map by
\path{verify_adjacent_cubic.py}; their index data are
also in \path{adjacent_cubic_certificates.json}.
Thus the permanent lower bound six holds in every
characteristic, and the determinant lower bound seven
holds away from characteristic two.

We now strengthen the determinant lower bound to eight.
The following proper incidence is a necessary condition
for border rank at most seven. Write the three copies of
$E$ as $B,C,D$. There must exist three nine-dimensional
subspaces
\[
 I_{BC}\subset\operatorname{Ann}_{BC}(T_-),\quad
 I_{BD}\subset\operatorname{Ann}_{BD}(T_-),\quad
 I_{CD}\subset\operatorname{Ann}_{CD}(T_-)
\]
of $B^*\otimes C^*$, $B^*\otimes D^*$, and
$C^*\otimes D^*$ respectively, such that
\[
 \dim\bigl(I_{BC}D^*+I_{BD}C^*+I_{CD}B^*\bigr)\leq57.
 \tag{P11.1}
\]
Here products mean insertion into the three tensor
factors, and $\operatorname{Ann}_{BC}(T_-)$ is the
kernel of contraction to $A\otimes D$. To prove the
condition, start with seven generic Segre atoms. Their
pairwise $BC,BD,CD$ projections each span a seven-plane.
The corresponding annihilators are nine-planes, and
their products annihilate all seven independent
$BCD$ projections. Thus their sum has dimension at
most $64-7=57$. The three Grassmannians are projective;
containment in the annihilators and (P11.1) are closed
conditions, the latter being a determinantal condition
on a map from three tautological bundles. Proper
projection therefore extends the condition from a
dense open set of seven-atom sums to its entire secant
closure. This proves the asserted necessary condition
over every algebraically closed field.

Let $H$ be the upper triangular subgroup of
$\operatorname{GL}(E)$ in the basis $(x,y,z,w)$.
It preserves the plane
$U=\langle\operatorname{alt}(x,y,z),
\operatorname{alt}(x,y,w)\rangle$. Give $A=U^*$ its
contragredient action. Then $H$ fixes $T_-$ and
preserves the projective incidence fiber just constructed.
This fiber has an $H$-fixed point, as follows directly
from triangular matrices. Embed the product of the
three Grassmannians by its Segre--Pl\"ucker embedding.
The ambient representation is a tensor product of
exterior powers of triangular representations, hence
has an $H$-stable complete flag $0=V_0\subset\cdots
\subset V_N$. Choose the least $s$ for which the fiber
meets $\mathbb P(V_s)$. This intersection is projective
and lies in the affine chart
$\mathbb P(V_s)\setminus\mathbb P(V_{s-1})$, so it is
finite. Since $H$ is connected, each of its points is
fixed. Thus all three nine-planes may simultaneously
be assumed $H$-stable. This argument uses triangularity
and properness in every characteristic; it makes no
semisimplicity assumption.

Let $\epsilon_0,\epsilon_1,\epsilon_2,\epsilon_3$ be
the dual basis. For any pair of the three $E$ factors put
\[
 S_{ii}=\epsilon_i\otimes\epsilon_i,
 \qquad
 S_{ij}=\epsilon_i\otimes\epsilon_j+
         \epsilon_j\otimes\epsilon_i\quad(i<j),
 \qquad
 Z=\epsilon_2\otimes\epsilon_3-
       \epsilon_3\otimes\epsilon_2.
\]
When the characteristic is not two, direct contraction
gives the eleven-dimensional space
\[
 \operatorname{Ann}_{BC}(T_-)
   =\langle S_{ij}:0\leq i\leq j\leq3\rangle\oplus kZ,
 \tag{P11.2}
\]
and the same description for the other two pairs.
Indeed the pair contraction has the five independent
images $a_0\otimes x,a_0\otimes y,a_1\otimes x,
a_1\otimes y,a_0\otimes z+a_1\otimes w$.

There are exactly two $H$-stable nine-planes in (P11.2):
\[
 \begin{split}
 J_0&=\langle S_{ij}:(i,j)\ne(0,0)\rangle,\\
 J_1&=\langle S_{ij}:(i,j)\notin\{(0,0),(0,1)\}\rangle
          \oplus kZ.
 \end{split}
 \tag{P11.3}
\]
Here is the full classification. The diagonal torus
has a one-dimensional weight space at each $S_{ij}$,
except that the $(2,3)$ weight also contains $Z$.
An invariant nine-plane cannot contain $S_{00}$:
its upper-triangular orbit spans all ten symmetric
vectors. If it contains $S_{01}$, that orbit spans
all nine symmetric vectors except $S_{00}$, so the
plane equals $J_0$. To check the latter statement,
the elementary substitution
$\epsilon_0\mapsto\epsilon_0+t\epsilon_1$ gives
$S_{01}+2tS_{11}$, and substitutions with larger
indices generate the other eight vectors. Two is
invertible here. If the plane contains neither
$S_{00}$ nor $S_{01}$, torus invariance excludes both
of those one-dimensional weight spaces. The remaining
space has dimension nine, forcing the plane to be
$J_1$. Both displayed spaces are stable, which also
verifies that the list is exact.

It remains to test the eight possible triples from
(P11.3). The multiplication map in (P11.1) has $64$
rows and $108$ columns. Order the rows lexicographically
in $B^*,C^*,D^*$; order the columns by pair $BC,BD,CD$,
then the basis in (P11.3) with the $S_{ij}$ in
lexicographic order and $Z$ last, then the free index
$0,1,2,3$. The complete certificate supplies an
order-$58$ minor for each triple, as follows:
\[
\begin{array}{c|rrrrrrrr}
 (J_{BC},J_{BD},J_{CD})
 &(0,0,0)&(0,0,1)&(0,1,0)&(0,1,1)
 &(1,0,0)&(1,0,1)&(1,1,0)&(1,1,1)\\ \hline
 \det&-4&4&-4&8&16&32&-64&-64
\end{array}
\]
The rows and columns are recorded in
\path{adjacent_cubic_borel_certificate.json}.
The standalone \path{verify_adjacent_cubic.py}
reconstructs every entry directly from (P11.3) and
evaluates these integer determinants by fraction-free
elimination, without importing the generator or using
a rank oracle. Every displayed minor is nonzero when
two is invertible, giving rank at least $58$ in all
eight cases and contradicting (P11.1). Hence
$\BR(T_-)\geq8$ in every odd characteristic and in
characteristic zero.

For the determinant upper bound put $z_0=z,z_1=w$ and
consider the two simultaneous tangent tensors
\[
 H_x=\sum_{i=0}^1a_i\otimes
 (z_i\otimes x\otimes y-x\otimes z_i\otimes y
                         +x\otimes x\otimes z_i),
\]
\[
 H_y=\sum_{i=0}^1a_i\otimes
 (-z_i\otimes y\otimes x+y\otimes z_i\otimes x
                         -y\otimes y\otimes z_i).
\]
Each has border rank at most three. For $H_x$, the
three-term curve is
\[
 \frac{1}{t}\left[
 \sum_{i=0}^1a_i\otimes(x+tz_i)\otimes(x-tz_i)\otimes(y+tz_i)
 -(a_0+a_1)\otimes x\otimes x\otimes y\right].
\]
For $H_y$, it is
\[
 \frac{1}{t}\left[
 \sum_{i=0}^1a_i\otimes(y-tz_i)\otimes(y+tz_i)\otimes(x-tz_i)
 -(a_0+a_1)\otimes y\otimes y\otimes x\right].
\]
Both numerators have zero constant coefficient, so
the displayed expressions extend polynomially over
$t=0$, with limits $H_x,H_y$. Direct expansion gives
the integral identity
\[
 T_-=H_x+H_y-
       \sum_{i=0}^1a_i\otimes(x+y)\otimes(x-y)\otimes z_i.
\]
The last sum consists of two pure tensors. Adding it
to the two curves produces an eight-term degeneration
of $T_-$. This upper bound in fact holds in every
characteristic: only the parameter $t$ is inverted,
and no integer is divided out. The standalone checker
reconstructs all eight terms, verifies cancellation of
every negative power, and checks all $128$ constant
tensor coefficients over $\mathbb Z$.

P-02 identifies every two-dimensional row minimum with
this adjacent pair. For $P_4$, R-01 with $b=3,j=2$
gives $\BR(P_4)\geq4-2+6=8$; P-10 with $b=4,j=1$
gives the matching upper bound eight. Finally the two
row tensors agree in characteristic two.
\end{proof}

\paragraph{Geometric dictionary.}
P-10 is an integral family of Segre sections over
$\mathbb G_m$ whose inclusion--exclusion sum extends
over $\mathbb A^1$. P-11 uses exterior-power maps of
locally free sheaves and their actual integer minors;
both upper degenerations are integral polynomial
families after their parameter poles cancel. The exact
foundations are EGA~I \S9.5,
SGA~3 Expos\'e~I \S\S4.5--4.6, and SGA~6 Expos\'e~V,
Definition~2.1, Examples~2.2(b) and Lemma~2.2.1, with
the printed-page and facsimile locators in the main
dictionary. The rank-eight row corollary further uses
R-01's EGA~II proper incidence. The determinant lower
bound uses a projective product of Grassmannians,
tautological multiplication maps, and a triangular
Segre--Pl\"ucker representation; the finite incidence
projection is proper by EGA~II, Proposition~5.4.2 and
Theorem~5.5.3(i). Its eight power-of-two minors work
over $\mathbb Z[1/2]$. No archive cote or
cohomological rank theorem is inferred from these
constructions.

\section{Bounds for the polynomial cubic adjacent pencil}

\begin{finding}{P-18}{a twenty-seven-term border upper bound and exact lower certificates for the cubic polynomial seed}
Let $M_-$ and $M_+$ be the first-factor-labelled ordered
coefficient tensors of the two $3\times3$ determinants,
respectively permanents, in the columns $012$ and $013$
of a $3\times4$ matrix. Their format is $2\times12\times12\times12$.
Over every algebraically closed field,
\[
 \BR(M_-)\leq27,\qquad 21\leq\BR(M_+)\leq24.
\]
If the characteristic is not three, also $\BR(M_-)\geq21$.
In characteristic two the two tensors coincide, so their
common border rank is between $21$ and $24$.
In particular, in characteristic zero the bounds are
$21\leq\BR(M_-)\leq27$ and $21\leq\BR(M_+)\leq24$.
No exact value for either polynomial pencil is asserted.
\end{finding}
\begin{proof}
Write $V=k^3$ and let $E$ have basis $e_{r,c}$ with
$0\leq r\leq2$ and $0\leq c\leq3$. Let
$X,Y,Z_0,Z_1:V\to E$ be the split embeddings into
columns $0,1,2,3$. Define the linear map
\[
 \Phi:V^{\otimes3}\longrightarrow k^2\otimes E^{\otimes3}
\]
by
\[
 \Phi(u\otimes v\otimes w)
 =\sum_{i=0}^1a_i\otimes
   \operatorname{per}\bigl(X(u),Y(v),Z_i(w)\bigr),
\]
where $\operatorname{per}$ means the six ordered tensor
words, including their multiplicities. This expression
is trilinear in $u,v,w$, hence defines the indicated
linear map. Expanding in the row basis gives
\[
 \Phi(D_3)=M_-,\qquad \Phi(P_3)=M_+.
\]
This is an identity of integral coefficient tensors;
no division by $3!$ or polynomial polarization is used.

For a pure input tensor, P-10 with $b=3,j=2$ supplies
a six-term border decomposition of its image. Therefore
\[
 \BR(\Phi(T))\leq6\BR(T)
 \tag{P18.1}
\]
for every $T$: the inverse image under $\Phi$ of the
closed $6r$-secant cone contains every sum of $r$
pure tensors and hence its closure. The one-label
specialization of P-11's determinant tangent formula
has five terms, giving $\BR(D_3)\leq5$ over
$\mathbb Z$ after parameter cancellation. P-10 gives
$\BR(P_3)\leq4$ in every characteristic. Equation
(P18.1) proves auxiliary upper bounds $30$ and $24$.

We improve the determinant upper bound to $27$ by a
partially symmetric intermediate space. First, if
$\operatorname{char}k\ne2$, every element of
$F\otimes\Gamma^2H$, where $\dim F=3$ and $\dim H=6$,
is a limit of sums of nine tensors $f\otimes v^{\otimes2}$.
Here $\Gamma^2H=(H^{\otimes2})^{\mathfrak S_2}$ is the
space of invariant tensors; its coordinates are the
coefficients at $(i,j)$ for $0\leq i\leq j\leq5$.
To prove this assertion consider the polynomial map
\[
 \Theta:(F\times H)^9\longrightarrow F\otimes\Gamma^2H,
 \qquad ((c_h,v_h))\longmapsto\sum_{h=0}^8c_h\otimes v_h^{\otimes2}.
 \tag{P18.3}
\]
Let $f_0,f_1,f_2$ denote the standard basis of $F$,
and use the following explicit point of the source:
\[
\begin{split}
 &c_h=f_{\lfloor h/3\rfloor}\quad(0\leq h\leq8),
 \qquad v_h=e_h\quad(0\leq h\leq5),\\
 &v_6=(1,1,0,1,1,0),\quad
 v_7=(1,0,1,2,0,1),\quad
 v_8=(0,1,1,0,2,2).
\end{split}
\]
Thus the nine labels are three copies of each basis
vector. Order the target
coordinates by $(a,i,j)$ with $a$ first and $i\leq j$
lexicographically; order the source coordinates by
the summand $h$, its three label coordinates, and
then its six vector coordinates. The Jacobian entries
are, explicitly,
\[
 \frac{\partial\Theta_{aij}}{\partial c_{hb}}
 =\delta_{ab}v_{hi}v_{hj},\qquad
 \frac{\partial\Theta_{aij}}{\partial v_{h\ell}}
 =c_{ha}(\delta_{i\ell}v_{hj}+\delta_{j\ell}v_{hi}).
\]
The certificate \path{nine_square_jacobian_certificate.json}
specifies $63$ columns of this $63\times81$ matrix;
their exact integer determinant is $-16$.
The standard-library checker
\path{verify_nine_square_jacobian.py} reconstructs the
point and every Jacobian entry from these formulas
and computes that determinant by fraction-free elimination.

For completeness, this Jacobian certificate proves
dominance without a characteristic-zero assumption.
Fix the other eighteen source coordinates at the
displayed point. Translate that point and its image
to zero. The resulting polynomial map from $\mathbb A^{63}$
to itself has an invertible linear term whenever
$\operatorname{char}k\ne2$. Solving one homogeneous
degree at a time gives a formal power series inverse.
A polynomial relation vanishing on its image would,
after composition with that inverse, be the zero
formal polynomial, and hence the zero polynomial.
Thus (P18.3) has dense image in these characteristics,
which proves the nine-term border assertion.

Apply this fact with $H=X(V)\oplus Y(V)$, the common
two columns of the two determinants. In the cofactor
expansion
\[
 \det[X,Y,Z_i]=\sum_{r=0}^2 (Z_i)_r\,q_r(X,Y),
\]
let $Q_r\in\Gamma^2H$ be the ordered coefficient tensor
of the signed quadratic cofactor $q_r$. In particular,
an off-diagonal monomial with coefficient $c$ contributes
$c(u\otimes v+v\otimes u)$; no coefficient is divided
by two. Put $Q=\sum_r f_r\otimes Q_r$ and identify
$Z_i(f_r)=e_{r,2+i}$. Define the linear map
\[
 \Psi:F\otimes\Gamma^2H\longrightarrow k^2\otimes E^{\otimes3}
\]
by inserting $Z_i(f)$ in each of the three positions
beside the two ordered factors of a quadratic tensor,
and summing with label $a_i$. In particular,
\[
 \Psi(f\otimes v^{\otimes2})=
 \sum_{i=0}^1a_i\otimes
 \bigl(v\otimes v\otimes Z_i(f)+v\otimes Z_i(f)\otimes v
                 +Z_i(f)\otimes v\otimes v\bigr).
\]
The cofactor expansion gives $\Psi(Q)=M_-$. Each
displayed image has border rank at most three: it
is the limit at zero of
\[
 \frac1t\left[\sum_{i=0}^1a_i\otimes
       (v+tZ_i(f))^{\otimes3}
       -(a_0+a_1)\otimes v^{\otimes3}\right].
 \tag{P18.4}
\]
There are three pure summands for $t\ne0$; their sum
is regular at zero. Consequently the inverse image
under $\Psi$ of the closed border-rank-$27$ cone
contains every sum of nine labelled squares and
therefore their entire closure. This proves
$\BR(M_-)\leq27$ in characteristic different from two.
In characteristic two $M_-=M_+$, and the already
proved permanent upper bound $24$ is stronger. Thus
the asserted determinant upper bound holds in every
characteristic.

There are also explicit one-parameter certificates
for the auxiliary determinant bound $30$ and the
permanent bound $24$.
The five-term row determinant curve has scalar poles
of order at most one. Apply the six-point difference
to each pure row term, with its difference parameter
$s=t^2$. After the difference is divided by $s^2$,
its error beyond the desired ordered cubic is divisible
by $s$. Multiplication by a row coefficient with a pole
of order at most one leaves an error divisible by $t$.
Thus the resulting thirty-term expression is regular
at zero with limit $M_-$. For the four-term row permanent
curve the scalar poles have order at most two, and
$s=t^3$ gives the corresponding twenty-four-term curve.
The standalone program
\path{verify_polynomial_cubic_upper.py} reconstructs both
curves and checks all $2\cdot12^3=3456$ constant
coefficients and every negative-power cancellation over
$\mathbb Z$. Its explicit atom lists are recorded in
\path{polynomial_cubic_upper_certificate.json}.

For the lower bounds use the double-Koszul map
\[
 K_M:A\otimes\Lambda^4E\otimes E^*
 \longrightarrow\Lambda^2A\otimes\Lambda^5E\otimes E,
 \qquad \dim A=2,
\]
given, for $M=\sum m_{a b c d}a_a\otimes e_b\otimes e_c\otimes e_d$, by
\[
 \alpha\otimes\eta\otimes\beta
 \longmapsto\sum m_{a b c d}
 (a_a\wedge\alpha)\otimes(e_b\wedge\eta)
                       \otimes\beta(e_c)e_d.
\]
Its matrix has $9504$ rows and $11880$ columns. Its rank
on a nonzero pure tensor is
$\binom{11}{4}=330$. Consequently
\[
 \BR(M)\geq\left\lceil\frac{\rank K_M}{330}\right\rceil.
 \tag{P18.2}
\]

The following table summarizes exact integer minors
in $905$ different weight blocks for each tensor,
with $1810$ certified minors in total:
\[
\begin{array}{c|r|l}
\text{tensor and certificate}&\text{sum of minor orders}
 &\text{nonvanishing condition}\\\hline
M_-&6680&\text{each determinant is a unit times }2^u3^v\\
M_+&6626&\text{every determinant is a unit}
\end{array}
\]
Here $u,v\geq0$, and a unit means $1$ or $-1$.
To explain why the block orders add, give $e_{r,c}$
weight $R_r+C_c$ and $a_i$ weight
\[
 -R_0-R_1-R_2-C_0-C_1-C_{2+i}.
\]
Every coefficient term of $M$ has total weight zero.
The map $K_M$ is equivariant. Its source weight at
$a_\alpha\otimes e_I\otimes e_k^*$ is
$\operatorname{wt}(a_\alpha)+\sum_{i\in I}\operatorname{wt}(e_i)
-\operatorname{wt}(e_k)$; its target weight at
$(a_0\wedge a_1)\otimes e_J\otimes e_l$ is
$\operatorname{wt}(a_0)+\operatorname{wt}(a_1)
+\sum_{j\in J}\operatorname{wt}(e_j)+\operatorname{wt}(e_l)$.
Different weights therefore contribute disjoint matrix
blocks, in every characteristic.

The complete row and column indices and integer
determinants are in
\path{polynomial_cubic_lower_certificate.json}.
The independent standard-library verifier
\path{verify_polynomial_cubic_lower.py} reconstructs
every selected entry from the displayed map, verifies
the weights and their distinctness, and evaluates all
minors by exact fraction-free elimination. It also
checks the prime divisors stated in the table. Thus
the first row gives rank at least $6680$ when the
characteristic is neither two nor three; the second
gives $6626$ in every characteristic. Both exceed
$20\cdot330=6600$, so (P18.2) gives the claimed lower
bounds. In characteristic two $M_-=M_+$, which also
covers the determinant there. The displayed minors
do not establish the determinant lower bound twenty-one
in characteristic three, and no such assertion is made.
\end{proof}

\paragraph{What this leaves open.}
P-02 identifies these tensors with the polynomial
$b=3,j=2$ minima for every $n\geq4$. At $n=4$ the
flattening image has dimension $16$. The permanent
upper bound therefore caps the result obtainable from
this particular substitution at $16-2+24=38$.
The determinant upper bound caps it at
$16-2+27=41$. Thus this particular two-plane cubic
substitution cannot improve the proved determinant
lower bound $43$, regardless of the unresolved exact
value of the cubic pencil. The present lower
certificates give $21$ for that pencil in characteristic
different from three.

\paragraph{Geometric dictionary for P-18.}
The maps $\Phi$ and $\Psi$ are maps of free integral
modules and their associated tensor bundles. The
Jacobian minor proves density for a map of affine
schemes after inverting two. The upper curves
are integral Segre sections on $\mathbb G_m$ whose
grouped sum extends over $\mathbb A^1$. The Koszul
map is a morphism involving exterior powers of locally
free sheaves. The split-torus characters give its
weight-block decomposition, and the actual integer
minors carry its rank information. These are the
EGA~I \S9.5 image/closure construction, the SGA~3
Expos\'e~I \S\S4.5--4.6 module/group schemes, and
the SGA~6 Expos\'e~V Definition~2.1, Examples~2.2(b)
and Lemma~2.2.1 exterior operations, with the exact
printed and facsimile locators in the common dictionary.
Neither a $K_0$ class nor an unproved smoothability
claim replaces these coefficient and minor checks.

\section{A characteristic-free obstruction for the exceptional permanent plane}
\label{sec:permanent-apolar}

\paragraph{Background for R-04 (necessary prolongation conditions).}
Let $T\in A\otimes B\otimes C$ over an algebraically closed field,
and let $U=T(C^*)$. If $r\leq\dim(A\otimes B)$ and
$\underline{\operatorname R}(T)\leq r$, there is an $r$-plane
$E\subseteq A\otimes B$ containing $U$ such that
\[
 \operatorname{rank}\Phi_A(E)\leq(\dim A-1)r,\qquad
 \operatorname{rank}\Phi_B(E)\leq(\dim B-1)r,
\]
where
\begin{align*}
 \Phi_A(E)&:A\otimes E\longrightarrow\Lambda^2A\otimes B,
 &v\otimes(a\otimes b)&\longmapsto(v\wedge a)\otimes b,\\
 \Phi_B(E)&:B\otimes E\longrightarrow A\otimes\Lambda^2B,
 &w\otimes(a\otimes b)&\longmapsto a\otimes(w\wedge b).
\end{align*}
These are the degree-$(210)$ and degree-$(120)$ necessary
border-apolarity tests. The following short proof records the
rectangular, arbitrary-characteristic form needed here.

\begin{proof}
On $\operatorname{Gr}(r,A\otimes B)$ these are bundle maps
from a constant bundle tensored with the tautological bundle;
the stated rank conditions are closed determinantal conditions.
If $E$ is spanned by $r$ independent pure elements
$a_\ell\otimes b_\ell$, the image of the first map is the sum
of the $r$ spaces $(A\wedge a_\ell)\otimes b_\ell$, each of
dimension $\dim A-1$. This proves its rank bound, and the
second follows in the same way.

The incidence $T(C^*)\subseteq E$ is closed, being the
vanishing of the map into the tautological quotient. Intersect
it with both determinantal loci in tensor space times the
Grassmannian. Its projection to tensor space is proper and
therefore has closed image. Tuples of $r$ independent pure
$A\otimes B$ elements form a nonempty open subset of the
irreducible product of $r$ affine Segre cones; $r$ distinct
matrix units show nonemptiness. With arbitrary $C$ factors,
their sums lie in the projected incidence. The closure of
these sums contains all rank-at-most-$r$ tensors, by density
of that open set in the sum parametrization, and hence
contains their border-rank closure. The required $E$ exists.
The argument applies in every characteristic.
\end{proof}

This is a known necessary part of border apolarity, not a new
general method. See Conner--Huang--Landsberg,
\emph{Bad and good news for Strassen's laser method},
author manuscript \texttt{CHLlaser5-6-21.pdf}, Proposition~2.1,
equation~(4), printed p.~7, and its version with $A,B$ interchanged.
The all-characteristic extension used here is supplied by the proof above.

\begin{finding}{R-04}{the exceptional permanent five-plane}
Let $A=k^5$ have basis $c_0,\ldots,c_4$, and let $B,C$ be
copies of $k^9$ with bases $a,b,c,d,e,f,g,h,i$. In the following
display $xy+yx$ abbreviates $x_B\otimes y_C+y_B\otimes x_C$:
\begin{align*}
 T={}&c_0\otimes(ae+ea+bd+db)
     +c_1\otimes(af+fa+cd+dc)\\
   &+c_2\otimes(bf+fb+ce+ec)
     +c_3\otimes(ah+ha+bg+gb)
     +c_4\otimes(ai+ia+cg+gc).
\end{align*}
Then $\underline{\operatorname R}(T)\geq11$ over every
algebraically closed field. This is the permanent coordinate
orbit $[0,1,3,6,7]$ in the rectangle enumeration.
\end{finding}

\begin{proof}
Let $U$ be the $45\times9$ matrix of $T(C^*)$, with $A\otimes B$
coordinates ordered first by $A$. Let
$Z_s(X)=c_s\wedge X\in\Lambda^2A\otimes B$ and set
\[
 K=[Z_0U\mid\cdots\mid Z_4U],\qquad
 E_0=[U\mid c_2\otimes a].
\]
The exact certificate specifies an integral $50\times90$
matrix $L$ with $82$ nonzero entries. Put
\[
 H=\begin{pmatrix}LZ_0\\\vdots\\LZ_4\end{pmatrix},
 \qquad R=\Phi_B(\operatorname{span}E_0).
\]
The included certificate and standalone verifier give the
following integer identities and unit minors:
\[
 LK=0,\qquad HE_0=0,
\]
\begin{center}
\begin{tabular}{c|c|c|c}
Matrix&Shape&Minor size&Exact determinant\\\hline
$K$&$90\times45$&40&$1$\\
$L$&$50\times90$&50&$-1$\\
$H$&$250\times45$&35&$-1$\\
$E_0$&$45\times10$&10&$1$\\
$R$&$180\times90$&82&$-1$
\end{tabular}
\end{center}
All row/column selections and the nonzero entries of $L$ are
in \path{permanent_apolar_certificate.json}. Every other
matrix is reconstructed from the five displayed quadrics.
The standard-library-only program
\path{check_permanent_apolar_certificate.py} checks the
zero products and computes the five determinants by exact
fraction-free elimination over $\mathbb Z$. Thus these facts
hold after every field extension and in every characteristic.

The ten columns of $E_0$ are independent, so $\dim U=9$.
The full row rank of $L$ and $LK=0$ imply
$\operatorname{rank}K\leq40$; its unit minor gives equality.
Therefore $\operatorname{im}K=\ker L$. Similarly, $H$ has
rank at least $35$, and its kernel contains the ten
independent columns of $E_0$. Hence
\[
 \ker H=U+k(c_2\otimes a).
\]
If $\underline{\operatorname R}(T)\leq10$, the preceding necessary conditions supply a
ten-plane $E\supset U$ with
$\operatorname{rank}\Phi_A(E)\leq40$ and
$\operatorname{rank}\Phi_B(E)\leq80$.
The first map already has rank $40$ on $A\otimes U$, so
its image on $A\otimes E$ must equal $\operatorname{im}K$.
For every $X\in E$ and every $s$ this forces
$LZ_s(X)=0$. Consequently $E\subseteq\ker H$; equality of
dimensions gives $E=\operatorname{span}E_0$. But the unit
$82$-minor of $R$ gives $\operatorname{rank}\Phi_B(E)\geq82$,
contradicting its upper bound $80$. This proves the claim.
\end{proof}

The first test forces a unique extension, so this proof
covers every possible ten-plane and needs no torus-fixed
assumption. Combined with the other coordinate-orbit
certificates, coordinate degeneration and the anchored
five-edge upper witness, it gives over every algebraically closed field, using the integral coefficient-tensor convention,
\[
 \min_{W\in\operatorname{Gr}(5,U_{4,2}^{\mathrm{per}})}
     \underline{\operatorname R}(P_W)=11,\qquad
 \underline{\operatorname R}(\mathrm{per}_4)\geq36-5+11=42.
\]
The global field scope uses the separate unit-minor checks for every other coordinate orbit as well as R-04.

\paragraph{Geometric dictionary.}
The necessary-condition lemma uses the tautological bundle on a relative Grassmannian,
exterior-power bundle maps, determinantal rank loci and a
proper incidence projection. This is the framework of
EGA~II, Definition~5.4.1, Proposition~5.4.2 and
Theorem~5.5.3(i), as in predecessor F-18.
R-04 is a certificate by morphisms of free $\mathbb Z$-modules
and unit minors, so its kernel/image implications survive
arbitrary field base change. In $K_0$ the first target has
class $\lambda^2[A]\,[B]$, but the actual minors, not this
class, supply the obstruction. There is no metric, norm,
$\ell$-adic weight or motivic input. No manuscript cote or
scan is claimed to contain this new calculation.

% Exact uniform j=3,4,5 residual certificates, independently verified by
% /root/pencil. Dependencies: finding/proof, amsmath/amssymb, longtable,
% booktabs, \BR. P-01, P-06, R-01 and R-04 must appear elsewhere.

\subsection{How the finite certificates cover the Grassmannian}

For this subsection \(n=4,b=2\) and \(U\) is either polynomial
residual image. Use the ordered list
\[
 \mathcal P=((1,2),(1,3),(1,4),(2,3),(2,4),(3,4))
\]
for both row and column pairs. Label the rectangle
\((\mathcal P_a,\mathcal P_b)\) by \(6a+b\), with
\(0\leq a,b\leq5\). Thus certificate subset labels range from
zero through \(35\); matrix-entry indices remain \(1,\ldots,4\).

For a coordinate subset \(S\) of cardinality \(j\), let
\(E_S\) be the span of its used matrix variables and put
\(m= \dim E_S\). Every slice is a nonsingular signed matching
on its four-coordinate support, so the last-factor flattening
has rank \(m\). For the residual tensor
\(T_S\in A_S\otimes E_S\otimes E_S\), \(\dim A_S=j\), use
\[
 K_{S,p}:\bigwedge^p A_S\otimes E_S^*
 \longrightarrow \bigwedge^{p+1}A_S\otimes E_S.
\]
Its pure-tensor cost is \(c_{j,p}=\binom{j-1}{p}\).
An integer minor of order \(d\) and nonzero determinant
therefore proves
\[
 \BR(T_S)\geq \left\lceil d/c_{j,p}\right\rceil
\]
over every field in which that determinant remains nonzero.
The proof of the cost is direct: for a pure tensor
\(a\otimes b\otimes c\), the map is the product of evaluation
by \(b\), the wedge map \(a\wedge-\), and the fixed output
vector \(c\). The wedge map has rank \(\binom{j-1}{p}\).
Rank subadditivity and closed determinantal loci prove the
border-rank bound.

Row permutations, column permutations, and transposition form
a group of order \(24^2\cdot2=1152\) preserving border rank
of coordinate residual tensors. In the determinant case they
may change individual slice signs, which are absorbed by
an invertible diagonal map on \(A_S\).
The complete orbit counts are
\[
 \begin{array}{c|r|r}
 j&\binom{36}{j}&\text{coordinate orbits}\\ \hline
 3&7140&21\\
 4&58905&99\\
 5&376992&454.
 \end{array}
\]
The source program \path{coordinate_koszul.py} constructs
all group images and all orbit representatives. Its modular
elimination is used only to locate candidate minors.
The independent certificate package
\path{global_j345_certificate.json} records the actual
row and column indices and integer determinant of every
minor needed below.

The independent, standard-library-only checker
\path{verify_global_bounds.py} rebuilds every minor
from the defining wedge map and verifies its determinant by
fraction-free integer elimination. It also enumerates
\emph{every} \(j\)-subset of the \(36\) rectangles and
checks that every subset whose ordinary support is below
the claimed threshold lies in one of the listed low-support
orbits. Thus no generic-rank sampling or unverified orbit
representative assumption is used. Finally P-01 transfers
the coordinate lower bound to all planes in the Grassmannian.

\begin{finding}{R-02}{the three-dimensional polynomial residual minimum}
For both polynomial families over every algebraically closed
field,
\[
 \mu_3(U_{4,2})=8.
\]
Consequently R-01 gives
\(\BR(\mathcal D_4),\BR(\mathcal P_4)\geq41\).
\end{finding}
\begin{proof}
Every three-rectangle subset with at least eight variables
is settled by its ordinary flattening. All remaining subsets
form a single orbit, represented by \(S=\{0,1,3\}\); it
contains \(48\) coordinate subsets and uses six variables.
For \(p=1\), where the pure cost is two, the determinant
tensor has an order-\(18\) minor of determinant \(-4\);
the permanent tensor has an order-\(16\) minor of determinant
\(1\). Hence the required lower bound eight holds for the
permanent in every characteristic and for the determinant
away from characteristic two. In characteristic two the
two tensors coincide, so the permanent certificate also
proves the determinant assertion.

For the upper bound use the anchored shape with one row
vertex, three column vertices, and its three edges.
P-06 gives border rank \(3+1+3+1=8\). This shape occurs
at \(S=\{0,1,2\}\). Hence the minimum is exactly eight.
Substitution gives \(36-3+8=41\).
\end{proof}

\begin{finding}{R-03}{the four-dimensional polynomial residual minimum}
For both polynomial families over every algebraically closed
field,
\[
 \mu_4(U_{4,2})=9.
\]
The resulting substitution bound is again \(41\).
\end{finding}
\begin{proof}
Subsets with at least nine variables are settled by the
ordinary flattening. Exactly \(900\) coordinate subsets
remain; their four orbits have the following \(p=1\)
certificates. The pure cost is three; each pair in the
table is (minor order, integer determinant).
\[
 \begin{array}{c|r r|r r}
 S&\multicolumn{2}{c|}{\text{determinant}}&
    \multicolumn{2}{c}{\text{permanent}}\\ \hline
 \{0,1,2,3\}&30&-4&30&1\\
 \{0,1,3,6\}&29&4&28&-1\\
 \{0,1,4,5\}&32&-1&32&-1\\
 \{0,1,9,21\}&31&1&31&-1
 \end{array}
\]
All permanent minors are units; the determinant minors
are nonzero away from characteristic two. In characteristic
two use equality of the two coefficient tensors. These
bounds exceed nine on the exceptional orbits and establish
the uniform lower bound.
The complete anchored \(2\times2\) shape
\(S=\{0,1,6,7\}\) has border rank \(4+2+2+1=9\) by P-06.
Thus the minimum is nine, and substitution gives
\(36-4+9=41\).
\end{proof}

\begin{finding}{R-05}{the five-dimensional minimum and the bound forty-two}
For both polynomial families over every algebraically closed
field, in the integral coefficient-tensor convention,
\[
 \boxed{\mu_5(U_{4,2})=11},\qquad
 \boxed{\BR(\mathcal D_4)\geq42,\quad
        \BR(\mathcal P_4)\geq42}.
\]
Over characteristic zero these are ordinary CP border-rank
lower bounds for the usual determinant and permanent
polynomial tensors, and therefore also symmetric
border-rank lower bounds.
\end{finding}
\begin{proof}
Ordinary flattenings settle every coordinate subset with
at least eleven variables. The remaining subsets have
exactly \(26\) orbits. Table~\ref{tab:five-minors}
gives one sufficient certificate for each family on every
one of these orbits. Its entries are triples
\((p;d;\Delta)\), where \(p\) is the Koszul exterior degree,
\(d\) the minor order, and \(\Delta\) its exact integer
determinant. For \(p=1\) the pure cost is four, and every
listed \(d\geq41\) certifies at least eleven. For \(p=2\)
the pure cost is six, and every listed \(d\geq61\)
does the same.

All listed permanent determinants are \(1\) or \(-1\).
The single permanent exception \(\{0,1,3,6,7\}\) is
precisely \(T_5\) of P-08; the unit-minor prolongation
obstruction R-04 gives its lower bound eleven in every
characteristic. All determinant minors in the table are
units times powers of two. They therefore prove the
determinant bound in every characteristic other than two,
and in characteristic two the permanent certificates apply
because the two integral tensors coincide. The
order-\(61\), determinant-\(-1\) certificate in the first
row is included specifically to remove the characteristic
three exceptional prime of a larger, unnecessary minor.

This proves the uniform lower bound. An anchored five-edge
shape in a \(2\times3\) grid, for example
\(S=\{0,1,2,6,7\}\), has border rank
\(5+2+3+1=11\) by P-06. Hence the minimum is attained and
equals eleven for both signs. R-01 now gives
\(36-5+11=42\).
\end{proof}

{\small
\begin{longtable}{@{}lcc@{}}
\caption{Complete low-support five-plane certificates.
Entries are \((p;d;\Delta)\).}
\label{tab:five-minors}\\
\toprule
Rectangle subset \(S\)&Determinant&Permanent\\
\midrule
\endfirsthead
\toprule
Rectangle subset \(S\)&Determinant&Permanent\\
\midrule
\endhead
\(\{0,1,2,3,4\}\) & \(2;61;-1\) & \(2;72;1\) \\
\(\{0,1,2,3,6\}\) & \(1;44;4\) & \(1;44;-1\) \\
\(\{0,1,2,3,8\}\) & \(1;46;-4\) & \(1;46;-1\) \\
\(\{0,1,2,3,9\}\) & \(1;46;2\) & \(1;46;-1\) \\
\(\{0,1,2,3,10\}\) & \(1;47;4\) & \(1;47;1\) \\
\(\{0,1,2,6,18\}\) & \(1;42;-2\) & \(1;41;1\) \\
\(\{0,1,2,9,21\}\) & \(1;48;1\) & \(1;48;-1\) \\
\(\{0,1,3,6,7\}\) & \(1;41;2\) & \(\text{R-04}\) \\
\(\{0,1,3,6,8\}\) & \(1;46;4\) & \(1;45;1\) \\
\(\{0,1,3,6,13\}\) & \(1;43;-2\) & \(1;42;-1\) \\
\(\{0,1,3,6,18\}\) & \(2;71;64\) & \(2;66;1\) \\
\(\{0,1,3,6,19\}\) & \(1;42;-2\) & \(1;42;-1\) \\
\(\{0,1,3,6,20\}\) & \(1;47;-4\) & \(1;47;-1\) \\
\(\{0,1,3,6,24\}\) & \(1;44;4\) & \(1;44;1\) \\
\(\{0,1,3,6,25\}\) & \(1;46;2\) & \(1;46;-1\) \\
\(\{0,1,3,6,30\}\) & \(1;48;-4\) & \(1;47;1\) \\
\(\{0,1,3,8,10\}\) & \(1;47;4\) & \(1;46;1\) \\
\(\{0,1,3,8,20\}\) & \(1;47;-4\) & \(1;47;-1\) \\
\(\{0,1,4,5,6\}\) & \(1;48;-1\) & \(1;48;-1\) \\
\(\{0,1,4,5,8\}\) & \(1;46;-1\) & \(1;46;-1\) \\
\(\{0,1,4,8,20\}\) & \(1;46;1\) & \(1;46;1\) \\
\(\{0,1,4,11,23\}\) & \(1;50;-1\) & \(1;50;1\) \\
\(\{0,1,6,7,21\}\) & \(1;44;1\) & \(1;44;-1\) \\
\(\{0,1,6,9,19\}\) & \(1;43;1\) & \(1;43;1\) \\
\(\{0,1,6,15,27\}\) & \(1;48;-1\) & \(1;48;1\) \\
\(\{0,1,9,22,23\}\) & \(1;50;1\) & \(1;50;1\) \\
\bottomrule
\end{longtable}
}

The table is not a claim of optimal Koszul rank for every
orbit; it records enough exact minors for the uniform
bound. R-04 is needed because the direct Koszul certificates
alone leave the exceptional permanent five-plane at ten.
No priority assertion is made by R-02, R-03 or R-05;
the literature comparison is a separate part of the paper.

\paragraph{Algebraic-geometric form.}
The Grassmannian minima are avoided closed relative Segre
secant loci for the tautological residual tensor.
The Koszul maps are morphisms of locally free sheaves
formed with exterior powers, and their integer minors
define determinantal ideals. The units in the certificate
give open rank conditions after every base change.
The only residual characteristic issue is removed by the
identity of the two signs in characteristic two.
R-04 adds a proper Grassmannian incidence with the
first-prolongation determinantal loci, the low-degree
Hilbert-function test of border apolarity.

% R-06 authored by /root; complete argument checked and written by
% /root/fields_literature. Dependencies: finding environment, \BR,
% R-01--R-05, P-01, P-03, P-06 and the integral convention.
\section{The same minima for larger matrix sizes}

\begin{finding}{R-06}{uniform residual minima as the matrix size grows}
For both polynomial families, every algebraically closed field,
and every $n\geq4$,
\[
 \mu_3(U_{n,2})=8,\qquad
 \mu_4(U_{n,2})=9,\qquad
 \mu_5(U_{n,2})=11.
\]
In particular,
\[
 \BR(\mathcal D_n),\ \BR(\mathcal P_n)
 \ \geq\ \binom n2^{\!2}+6\qquad(n\geq4).
\]
For the two row families, every $n\geq2$ and
$1\leq j\leq n-1$ satisfy
\[
 \mu_j(U_{n,2})=j+1.
\]
All statements use the integral coefficient-tensor convention.
\end{finding}
\begin{proof}
P-01 reduces the polynomial assertion to sets of $j$
coordinate $2\times2$ rectangles. Let $E$ be the union of
their matrix-entry supports, and write $m=|E|$. As in the
finite certificates, the last-factor flattening has rank
$m$. Moreover, every used row and column contains at least
two elements of $E$, since each belongs to a selected
rectangle.

For $j=3$, the cases $m\geq8$ already have the required
lower bound. If $m\leq7$, at most three rows and three
columns occur, so row and column relabeling embeds the
configuration in a $4\times4$ matrix. R-02 applies there.
For $j=4$, the same argument with $m\leq8$ gives at most
four rows and four columns, so R-03 applies. Passing
between a tensor and its larger coordinate ambient space
does not change border rank: inclusion and coordinate
projection give the two inequalities.

For $j=5$ the cases $m\geq11$ are settled. If $m\leq10$,
at most five rows and five columns occur. A configuration
using at most four of each is handled by R-05. Suppose
that five rows occur. Then $m=10$, and every row has
exactly two column neighbors. Partition the five rows
according to their unordered pair of column neighbors.
Every part has at least two rows: a used row belongs to
a selected rectangle, and its other row has the same
two neighbors. A selected rectangle lies entirely in
one part. If the partition is $3+2$, at most
$\binom32+\binom22=4$ distinct rectangles exist, which
cannot contain the selected five. The only remaining
partition is $5$. Consequently the support is a
$5\times2$ grid. The case of five columns follows by
transposition.

The five selected rectangles in a $5\times2$ grid are
five edges on its five row vertices, with no isolated
vertex. There are exactly $222$ such labeled graphs:
from the $\binom{10}{5}=252$ five-edge subsets remove
the $5\binom65=30$ subsets with an isolated vertex.
Two isolated vertices are impossible for five edges.
The five permutation orbits and sufficient Koszul
certificates are as follows. Vertices in this table
are numbered $1,\ldots,5$; an entry $ij$ denotes an edge.
\[
\begin{array}{c|r|rr}
 \text{edge set}&\text{orbit size}&\Delta_{\det}&\Delta_{\mathrm{per}}\\ \hline
 12,13,14,15,23&30&-4&1\\
 12,13,14,23,25&60&-4&1\\
 12,13,14,23,45&60&4&1\\
 12,13,14,25,35&60&-1&1\\
 12,13,24,35,45&12&1&1
\end{array}
\]
Each entry is the determinant of an order-$41$ minor of
the $p=1$ Koszul map for the five-dimensional first
factor and the ten-dimensional support space. Its
pure-tensor cost is $\binom41=4$, so every nonzero
minor proves border rank at least
$\lceil41/4\rceil=11$. The certificate file
\path{five_by_two_extra.json} specifies all row and
column indices. The separate checker
\path{verify_five_by_two.py} reconstructs the maps
directly from the five ordered quadrics, computes these
ten determinants by integer elimination, and checks
that the disjoint permutation orbits cover all $222$
graphs. The permanent minors are units, and the
determinant minors are nonzero away from characteristic
two. In characteristic two the two tensors coincide.
This proves the lower bound in every characteristic.

The anchored shapes with $(h,r,s)=(3,1,3),(4,2,2),
(5,2,3)$ give upper bounds $8,9,11$ by P-06.
Each embeds in a $4\times4$ matrix, so these upper
bounds hold for all $n\geq4$. R-01 with $j=5$ now
gives $\binom n2^{2}-5+11=\binom n2^{2}+6$.

For a row residual space and any $j$-plane $W$, its
first factor is concise of dimension $j$. The all-field
one-dimensional residual theorem FL-03 gives matrix
rank at least two for every nonzero element of $W$.
Applying R-01 to $P_W$ while retaining one dimension
therefore gives $\BR(P_W)\geq j-1+2=j+1$.
For $j\leq n-1$, the coordinate star on the edges
$\{0,i\}$, $1\leq i\leq j$, is available. P-03 gives
its border rank $j+1$, proving equality.
\end{proof}

\begin{finding}{R-07}{monotonicity of the residual gain and the sixth minimum}
For any fixed residual space $U$ and
$1\leq i\leq j\leq\dim U$,
\[
 \mu_j(U)-j\ \geq\ \mu_i(U)-i.
\]
For both polynomial families in the integral convention,
over every algebraically closed field and every $n\geq4$,
\[
 \mu_6(U_{n,2})=12.
\]
\end{finding}
\begin{proof}
Apply R-01 to the concise first factor of $P_W$, where
$W\subset U$ has dimension $j$, retaining dimension $i$.
The surviving tensors are precisely $P_V$ for
$V\in\operatorname{Gr}(i,W)$. Hence
\[
 \BR(P_W)\geq j-i+
 \min_{V\in\operatorname{Gr}(i,W)}\BR(P_V)
 \geq j-i+\mu_i(U).
\]
Taking the minimum over $W$ proves the inequality.
R-06 then gives $\mu_6\geq6-5+11=12$.
The complete anchored $2\times3$ shape has six edges
and fits in a $3\times4$ matrix. P-06 gives its exact
border rank $6+2+3+1=12$, proving the matching upper
bound. This sixth minimum yields the same substitution
bound as the fifth minimum.
\end{proof}

% Exact characteristic-zero certificates for four determinant seven-planes.
% Depends on the necessary-prolongation lemma preceding R04/R-04.

\section{Four determinant seven-plane obstructions}

\begin{finding}{P-09}{two determinant types and their transposes have border rank at least fourteen}
In the rectangle numbering of the finite certificates, let
\[
 \begin{split}
 S_A&=\{0,1,3,6,7,9,12\},\\
 S_B&=\{0,1,3,6,7,12,13\}.
 \end{split}
\]
The corresponding determinant residual tensors have border
rank at least \(14\) over every algebraically closed field
of characteristic zero. The same holds for their transposes,
whose subsets are
\[
 S_A^{\mathsf t}=\{0,1,2,6,7,18,19\},\qquad
 S_B^{\mathsf t}=\{0,1,2,6,7,8,18\}.
\]
These are individual residual-orbit lower bounds; a uniform
seven-plane minimum also requires the other coordinate or
Borel-fixed cases.
\end{finding}
\begin{proof}
We supply the whole reduction and the finite exact certificate.
Let \(T\in A\otimes B\otimes C\) be one of the two tensors,
where \(\dim A=7\), and \(B,C\) are the spaces of its used
matrix variables. Their common dimension \(m\) is \(11\)
for type \(A\) and \(12\) for type \(B\). Every slice is a
nonsingular signed matching on its four-coordinate support,
and the union uses all \(m\) variables; therefore
\[
 U=T(C^*)\subset A\otimes B,\qquad \dim U=m.
\]

Use the necessary-prolongation condition proved before R-04.
If \(\BR(T)\leq13\), there is a thirteen-plane
\(L\supset U\) such that
\[
 \dim\ker\Psi_A(L)\geq13,\qquad
 \dim\ker\Psi_B(L)\geq13.
 \tag{*}
\]
Here \(\Psi_A:L\otimes A\to B\otimes\bigwedge^2A\) and
\(\Psi_B:L\otimes B\to A\otimes\bigwedge^2B\) are the
exterior maps. They differ from the earlier notation only
by factor permutations and an overall sign. In particular
\[
 \operatorname{rank}\Psi_A(L)\leq6\cdot13=78,\qquad
 \operatorname{rank}\Psi_B(L)\leq(m-1)13.
 \tag{**}
\]

Give a variable \(e_{ij}\) the row/column torus weight
\(R_i+C_j\), and give a slice label indexed by \((I,J)\)
the weight \(-\sum_{i\in I}R_i-\sum_{j\in J}C_j\).
This split torus stabilizes \(T\). It preserves \(U\), the
incidence \(L\supset U\), and both closed rank conditions
in (*). Choose an integral one-parameter subgroup separating the finitely
many ambient torus characters. Its limit in this projective incidence
is a direct sum of subspaces of the corresponding weight blocks and allows \(L\) to be torus-invariant. Hence \(L/U\) is a sum
of weight spaces.

Let \(V_w\) denote a full weight space of \(A\otimes B\).
The quotient \((A\otimes B)/U\) has \(47\) nonzero weight
blocks for type \(A\), and \(50\) for type \(B\).
These numbers and every weight coordinate are explicitly
enumerated in the certificate. For type \(B\),
\(\dim(L/U)=1\); its weight \(w\) satisfies
\(L\subset U+V_w\). For type \(A\),
\(\dim(L/U)=2\); either both dimensions lie in one weight
block of quotient dimension at least two, or they are
lines in two distinct blocks. There are \(14\) choices
of the former kind and \(\binom{47}{2}=1081\) of the
latter, for \(1095\) choices in total.

\paragraph{First stage: full weight-space extensions.}
For a chosen set of weights \(D\), put
\(L_D=U+\sum_{w\in D}V_w\). It contains every possible
\(L\) with those quotient weights. The certificate gives
a minor of one of the two exterior maps on \(L_D\)
of order
\[
 \dim(\text{its domain})-12
\]
with nonzero determinant modulo \(1000003\).
Consequently the corresponding kernel on \(L_D\) has
dimension at most \(12\), as does the kernel on any
subspace \(L\subset L_D\). This contradicts (*).
This stage excludes \(1042\) of the \(1095\) type \(A\)
choices, and \(47\) of the \(50\) type \(B\) choices.
In particular, no type \(A\) choice with both quotient
dimensions in one weight remains.

\paragraph{Second stage: a finite linear test on each weight line.}
We describe the test used on the remaining weights.
Fix \(w\) and write a quotient vector \(X\in V_w\)
in a basis of \(V_w/(U\cap V_w)\).
If \(U\cap V_w\ne0\), its nonzero vector has a unit
coefficient on every coordinate in its support, so one
coordinate may be omitted to choose this quotient basis.
If \(U\cap V_w=0\), use the whole coordinate basis.
Thus \(X\) has a nonzero parameter vector of dimension
at most five in all cases used below.

For an exterior map in a factor \(F\), let
\(\pi\) be the quotient of its target by the image of
that map on \(U\otimes F\). If \(f_1,\ldots,f_d\) are
the coordinate vectors of \(F\), then
\[
 \pi\Psi(X\otimes f_i)
\]
has weight \(w+\operatorname{wt}(f_i)\).
These weights are pairwise distinct. The rank increment
on adjoining \(X\) is therefore exactly the number of
nonzero projected columns.

The certificate supplies integer linear functionals that
annihilate \(\Psi(U\otimes F)\), separately in each weight
block. Their annihilation is checked as an exact integer
zero identity. Applying them to the new columns gives
matrices \(H_i\), linear in the quotient coordinates of
\(X\). If a projected column is zero, then \(H_iX=0\).
To prove a rank increment at least \(\ell\), it is enough
to check that for every set of \(d-\ell+1\) columns,
the stacked matrix of their \(H_i\)'s has full column
rank. This excludes a nonzero parameter vector for which
all those columns vanish. All these tests are finite
linear rank tests, certified modulo the same explicit
prime; no sampled or generic parameter is substituted.

\paragraph{Type \(B\).}
A certified minor gives
\[
 \operatorname{rank}\Psi_B(U)\geq137.
\]
The three remaining quotient weights have dimensions
\(3,5,3\). For each, the certificate verifies full column
rank for every one of the
\(\binom{12}{6}=924\) stacks corresponding to six zero
columns. Thus adjoining any nonzero line of one of these
weights increases rank by at least seven. This gives
\[
 \operatorname{rank}\Psi_B(L)\geq137+7=144
   >11\cdot13=143,
\]
contradicting (**).

\paragraph{Type \(A\).}
A certified minor gives
\[
 \operatorname{rank}\Psi_A(U)\geq72.
\]
The \(53\) remaining choices consist of pairs of distinct
weight lines. They use \(27\) quotient weight blocks.
For each such block \(w\), the finite linear test gives
a uniform increment lower bound \(\ell_w\), one of
\(2,3,4,6\), for the seven \(A\)-columns.
For two weights \(w,v\), let
\[
 F_w=\{w+\operatorname{wt}(a_s):1\leq s\leq7\},
 \qquad k_{w,v}=|F_w\cap F_v|.
\]
The projected column spaces of the two lines can intersect
only in these common weights; each line contributes at
most one vector in each weight. Their intersection
therefore has dimension at most \(k_{w,v}\).
The total rank increment is at least
\(\ell_w+\ell_v-k_{w,v}\).
For every one of the \(53\) remaining pairs, the certificate
checks
\[
 \ell_w+\ell_v-k_{w,v}>6.
\]
Consequently
\(\operatorname{rank}\Psi_A(L)\geq72+7=79>78\),
again contradicting (**).

All possible quotient weight choices have now been excluded.
Thus neither tensor has border rank at most \(13\).
Transposition is an invertible permutation of matrix
coordinates and preserves the determinant slices, proving
the same result for the two transposed subsets.
\end{proof}

\paragraph{Certificate and independent verifier.}
The complete data are in
\path{det7_two_type_certificate.json}.
The independent program
\path{verify_det7_certificate.py} reconstructs the
tensors, weights, and every required exterior-map entry
directly. It does not import the generator or use a
rank/nullspace oracle. It verifies the baseline minors,
all \(1089\) full-extension minors, all integer annihilator
identities, every required linear zero pattern, every
weight-overlap inequality, and complete coverage of the
\(1095+50\) quotient choices.
It uses integer NumPy arrays for modular arithmetic;
the prime is smaller than \(2^{20}\), all modular products
are smaller than \(2^{40}\), and explicit coefficient
bounds protect the unreduced integer zero identities.
The proof is therefore exact finite arithmetic, with no
floating-point tolerance. A nonzero minor modulo the
verified prime certifies a nonzero integer minor and
therefore the stated characteristic-zero result.

\paragraph{Geometric dictionary.}
The ambient construction is the proper Grassmannian
incidence with first-prolongation determinantal loci.
The torus limit gives a weight decomposition of a quotient
of the tautological bundle. Saturating selected weights
replaces a variable subbundle by an explicit larger
coordinate subbundle. The second stage tests linear
sections of projective quotient weight spaces; full-column
rank proves that the corresponding projective common zero
locus is empty. All maps and annihilator identities are
integral. Their nonzero modular minors certify the required
rank conditions in characteristic zero; no assertion about
all positive characteristics is made here.

% Additional fully certified cases, using the complete criterion of P-09.

\begin{finding}{P-12}{four more determinant seven-planes have border rank at least fourteen}
The determinant residual tensors for
\[
 \begin{split}
 S_C&=\{0,1,2,3,6,7,8\},\\
 S_D&=\{0,1,2,3,4,6,12\}
 \end{split}
\]
and their transposes
\[
 S_C^{\mathsf t}=\{0,1,6,7,12,13,18\},\qquad
 S_D^{\mathsf t}=\{0,1,2,6,12,18,24\}
\]
have border rank at least \(14\) over every algebraically
closed field of characteristic zero.
\end{finding}
\begin{proof}
Both tensors have \(A\)-dimension seven and used-variable
dimension twelve in each of the other factors. As in P-09,
their last-factor images \(U\subset A\otimes B\) have
dimension twelve. Border rank at most thirteen would give
a thirteen-plane \(L\supset U\) satisfying both necessary
prolongation conditions. Taking a torus limit makes
\(L/U\) a homogeneous line.

There are \(52\) possible quotient weight blocks for
type \(C\) and \(60\) for type \(D\). The exact certificate
excludes respectively \(50\) and \(58\) of these by
full-weight-space saturation: on the corresponding
larger space \(L_w=U+V_w\), one of the two prolongations
has kernel dimension at most twelve. The proof of this
exclusion and its monotonicity under \(L\subset L_w\)
are exactly those supplied in P-09.

In both cases a certified baseline minor gives
\(\operatorname{rank}\Psi_A(U)\geq74\). The necessary
rank bound on \(A\otimes L\) is \(6\cdot13=78\), so
the added homogeneous line could increase this rank
by at most four. Its seven coordinate columns have
distinct torus weights. The two remaining quotient
blocks have dimensions \(5,3\) for type \(C\), and
\(4,4\) for type \(D\). For each of the four blocks,
the certificate constructs integer annihilators of
\(\Psi_A(U)\) and verifies that every stack of the
linear conditions for three zero columns has full
column rank. There are \(\binom73=35\) such stacks
per block. Thus no nonzero quotient vector can make
three projected columns vanish, and every added line
increases rank by at least five. This gives
\[
 \operatorname{rank}\Psi_A(L)\geq74+5=79>78,
 \]
the required contradiction.
Every weight line has been treated, and transposition
preserves border rank, proving all four assertions.
\end{proof}

The complete data are
\path{det7_extra_m12_certificate.json}.
The independent verifier
\path{verify_det7_certificate.py}, with this file as
its \texttt{--certificate} argument, reconstructs all
matrices and confirms \(108\) saturation minors,
the two baseline minors, all integer annihilator
identities, and \(140\) linear zero-pattern tests.
The exact modular and characteristic-zero interpretation
is the same as in P-09. The algebraic-geometric dictionary
is likewise P-09's proper Grassmannian incidence,
torus weight quotient, and empty projective linear
sections; no additional geometric hypothesis is used.

% Exact support-11 determinant seven-plane obstructions.

\begin{finding}{P-14}{the remaining support-eleven determinant types}
The determinant residual tensors for
\[
 \begin{split}
 S_E&=\{0,1,2,3,6,7,9\},\\
 S_F&=\{0,1,2,3,6,7,18\}
 \end{split}
\]
and their transposes
\[
 S_E^{\mathsf t}=\{0,1,6,7,12,18,19\},\qquad
 S_F^{\mathsf t}=\{0,1,3,6,7,12,18\}
\]
have border rank at least \(14\) over every algebraically
closed field of characteristic zero.
\end{finding}
\begin{proof}
Use the full proper-incidence, torus, saturation, and
linear increment criteria proved in P-09. Here both
tensors have \(A\)-dimension seven and used-variable
dimension eleven in \(B,C\), so \(U=T(C^*)\) has
dimension eleven. Border rank at most thirteen would
give a torus-invariant thirteen-plane \(L\supset U\),
with two-dimensional quotient, satisfying
\[
 \operatorname{rank}\Psi_A(L)\leq78,\qquad
 \operatorname{rank}\Psi_B(L)\leq130.
 \tag{1}
\]

\paragraph{Type \(E\).}
The quotient has \(49\) nonzero weight blocks. There
are \(1189\) possible weight choices for a two-plane:
\(\binom{49}{2}=1176\) distinct-weight pairs and \(13\)
single-weight choices of dimension at least two.
The certificate excludes \(1185\) choices by full
weight-space saturation and a kernel bound at most twelve.
The four remaining choices are distinct-weight pairs.
A baseline minor gives
\(\operatorname{rank}\Psi_A(U)\geq75\).
Their five used weight blocks have uniform individual
increment lower bounds \(5,5,6,5,4\).
The zero-pattern certificates proving these bounds are
the integer-annihilator linear tests of P-09.
For each of the four pairs, subtracting the number of
common output weights from the sum of the two increment
bounds still gives at least four. Hence
\(\operatorname{rank}\Psi_A(L)\geq75+4=79\), contradicting
(1).

\paragraph{Type \(F\): all but one weight pair.}
The quotient has \(48\) nonzero weight blocks and
\(\binom{48}{2}+15=1143\) possible weight choices.
Saturation excludes \(1076\) of them. A baseline minor
gives \(\operatorname{rank}\Psi_A(U)\geq71\).
The individual-line increment and output-weight-overlap
tests exclude \(66\) of the remaining \(67\) distinct
pairs: in each, the total increment is at least eight,
giving rank at least \(79>78\).

\paragraph{The last type \(F\) pair.}
This pair consists of quotient blocks numbered \(11\)
and \(14\) in the certificate, with quotient dimensions
one and two. Let \(X\) be the selected coordinate
representative of the first block. Since that quotient
weight space is one-dimensional, every possible \(L\)
with this pair contains the same intermediate space
\[
 U_0=U+kX,\qquad \dim U_0=12.
 \tag{2}
\]
Changing the representative by an element of \(U\)
does not change (2). The other quotient line has a
nonzero two-component coordinate vector \(Y\), modulo
\(U\); its weight is different from that of \(X\), so
its quotient basis is unchanged after adjoining \(X\).

An additional exact minor gives
\(\operatorname{rank}\Psi_B(U_0)\geq125\).
Condition (1) allows an increment at most five when
adjoining \(Y\). Its eleven \(B\)-columns have pairwise
distinct torus weights. The certificate supplies
integer annihilators of \(\Psi_B(U_0)\), not merely
of \(\Psi_B(U)\), and checks every one of the
\(\binom{11}{6}=462\) stacks of linear conditions for
six zero columns. Each has full column rank two.
Thus no nonzero \(Y\) can make six columns vanish;
the increment is at least six. Therefore
\[
 \operatorname{rank}\Psi_B(L)\geq125+6=131>130,
\]
which excludes the final pair.

All \(1189+1143\) quotient choices have been covered.
Both tensors have border rank greater than thirteen,
and transposition gives the other two assertions.
\end{proof}

The complete certificate is
\path{det7_extra_m11_certificate.json}.
The independent verifier
\path{verify_det7_certificate.py} reconstructs all
maps directly, verifies \(2261\) saturation minors,
the baseline and conditional-base minors, the integer
annihilator identities, every required zero pattern,
and full coverage of the quotient choices. Its
conditional-base check explicitly verifies that the
fixed block has quotient dimension one, that the two
weights differ, and that the new annihilators kill
the enlarged image from (2).
The modular arithmetic certifies characteristic-zero
rank inequalities exactly, with the safeguards and
geometric dictionary already given in P-09.

\section{Five permanent seven-plane obstructions}

\begin{finding}{P-13}{five permanent coordinate orbits have border rank at least fourteen}
In the standard rectangle numbering, the permanent residual tensors
for each of the following subsets have border rank at least fourteen
over every algebraically closed field of characteristic zero:
\[
\begin{aligned}
 C_1&=\{0,1,2,3,4,6,12\},&
 C_2&=\{0,1,2,3,6,7,8\},\\
 C_3&=\{0,1,2,6,7,8,18\},&
 C_4&=\{0,1,2,6,7,18,20\},\\
 C_5&=\{0,1,3,6,7,9,14\}.&&
\end{aligned}
\]
This is an assertion about five specified orbits; it does not alone
evaluate the uniform seven-plane minimum.
\end{finding}
\begin{proof}
Each tensor has format $7\times12\times12$, and its last-factor
image $U\subset A\otimes B$ has dimension twelve. If its border
rank were at most thirteen, the necessary conditions preceding
R-04 would supply a thirteen-plane $L\supset U$ with
\[
 \rank\Psi_A(L)\leq78,\qquad\rank\Psi_B(L)\leq143.
\]
The row/column torus preserves the tensor and these determinantal
incidences. The one-parameter fixed-limit argument of P-09 therefore
allows $L/U$ to be one homogeneous line.

We use exactly the complete saturation and linear-column criterion
proved in P-09, with permanent coefficients $+1$ in all four
ordered entries of each slice. The finite certificate gives:
\begin{center}
\begin{tabular}{c|r|r|c|r|r}
Case & Weights & Saturation & Map
 & Base rank & Increment\\\hline
$C_1$ &60&58&$\Psi_A$&74&5\\
$C_2$ &52&49&$\Psi_A$&74&5\\
$C_3$ &50&45&$\Psi_B$&136&8\\
$C_4$ &50&46&$\Psi_A$&74&5\\
$C_5$ &56&56&none&78&--
\end{tabular}
\end{center}
For every saturation exclusion, the certificate exhibits a nonzero
minor of order $\dim(L_w\otimes F)-12$ in an exterior map on
$L_w=U+V_w$. Its kernel has dimension at most twelve. Inclusion
$L\subset L_w$ then contradicts the necessary kernel dimension
at least thirteen, as in P-09.

For $C_1,C_2,C_4$ the remaining quotient blocks number two,
three and four, respectively. An integer annihilator of each
baseline exterior image turns vanishing of a new projected
column into explicit linear conditions on a nonzero quotient
vector. Every stack for three of the seven $A$ columns has
full column rank. Thus at most two columns vanish, and the
rank increment is at least five. Its distinct nonzero column
weights ensure independence. Hence $\rank\Psi_A(L)\geq79>78$.
For $C_3$ there are five remaining blocks. Every stack for
five of the twelve $B$ columns has full column rank, so at
least eight projected columns are nonzero and independent.
Thus $\rank\Psi_B(L)\geq144>143$.
Case $C_5$ has no remaining block. All possibilities have
therefore been excluded.

The data in \path{per7_m12_certificate.json} specify the
baseline and saturation minors and all integer annihilators.
The independent \path{verify_seven_plane_certificate.py}
reconstructs the ordered permanent coefficients and exterior
maps directly; it imports no generator or rank/nullspace oracle.
It verifies all $268$ quotient-weight choices, $254$ saturation
minors, and $4275$ linear zero-pattern tests. All annihilator
identities are checked over the integers, and all claimed
nonzero minors are checked modulo the verified prime $1000003$.
The integer products obey the same explicit bounds as in P-09.
A nonzero residue implies a nonzero integer minor, proving
the stated characteristic-zero assertions.
\end{proof}

The geometric dictionary is P-09's proper Grassmannian incidence,
quotient weight bundles and empty projective linear loci, with the
permanent sign convention. It does not require a permanent Borel
compression, and none is asserted.

% P-15. Complete proof; the accompanying exact checker supplies the finite certificate.
\subsection{The support-nine seven-planes}
\begin{finding}{P-15}{a border-apolarity obstruction for three support-nine seven-planes}
Let $k$ be an algebraically closed field of characteristic zero.  Put
$B=C=k^{\{0,1,2\}\times\{0,1,2\}}$, and for row and column pairs
$I=\{r<s\}$, $J=\{c<d\}$ put
\[
 Q^{\epsilon}_{I,J}=
 x_{rc}\otimes x_{sd}+x_{sd}\otimes x_{rc}
 +\epsilon(x_{rd}\otimes x_{sc}+x_{sc}\otimes x_{rd}),
 \qquad \epsilon\in\{-1,1\}.
\]
Here the two displayed factors are $B,C$, rather than multiplication in a
symmetric algebra.  In the following two ordered lists the notation $01$
means the pair $\{0,1\}$:
\begin{align*}
 \mathcal S_A&=((01,01),(01,02),(01,12),(02,01),(02,02),(02,12),(12,01)),\\
 \mathcal S_B&=((01,01),(01,02),(01,12),(02,01),(02,02),(12,01),(12,12)).
\end{align*}
For $A=k^7$ with basis $a_0,\ldots,a_6$, define
$T^{\epsilon}_{\mathcal S}=\sum_{q=0}^6a_q\otimes Q^{\epsilon}_{\mathcal S(q)}$.
Then
\[
 \underline R(T^{-}_{\mathcal S_A})\ge14,\qquad
 \underline R(T^{+}_{\mathcal S_A})\ge14,\qquad
 \underline R(T^{+}_{\mathcal S_B})\ge14.
\]
The first of these is the determinant coordinate type
$[0,1,3,6,7,9,18]$ in the four-by-four rectangle indexing; the other two
are the permanent types $[0,1,3,6,7,9,18]$ and $[0,1,3,6,7,18,21]$.
\end{finding}

\paragraph{Proof and finite certificate.}
We give all reductions, including the meaning and completeness of the
finite linear-arrangement enumeration.  The attached checker
\path{check_support9_sevenplanes.py} reconstructs the matrices below
from the displayed lists and signs.  It reads no project data and imports
no other project code.  Its ranks and kernel computations are over
$\mathbb Q$; there are no floating-point or finite-field rank decisions.
The three resulting verification records are
\path{support9_det_verification.json},
\path{support9_perA_verification.json}, and
\path{support9_perB_verification.json}.

\emph{The necessary extension.}
Let $U=T(C^*)\subset A\otimes B$.  The nine contractions are independent,
so $\dim U=9$.  The proper-incidence argument used in R-04, equivalently
the degree $(2,1,0)$ and $(1,2,0)$ border-apolarity necessary conditions,
says that $\underline R(T)\le13$ would provide a $13$-plane
$E\supset U$ such that
\begin{align}
 \operatorname{rank}\Phi_A(E)&\le78,&
 \Phi_A(E):A\otimes E&\longrightarrow\Lambda^2A\otimes B,\label{eq:P15-A}\\
 \operatorname{rank}\Phi_B(E)&\le104,&
 \Phi_B(E):B\otimes E&\longrightarrow A\otimes\Lambda^2B.\label{eq:P15-B}
\end{align}
On a pure vector $a\otimes b\in E$, the two maps are
$a'\otimes(a\otimes b)\mapsto a\wedge a'\otimes b$ and
$b'\otimes(a\otimes b)\mapsto a\otimes(b\wedge b')$.
The bounds are $(7-1)13$ and $(9-1)13$.
This invocation is a necessary condition, not an assertion that these two
conditions alone characterize border rank.

The exact baseline matrices have dimensions and ranks
\[
\begin{array}{c|c|c|c}
 &\text{matrix on }U&\text{rank on }U&\text{allowed rank increment}\\\hline
 \Phi_A&189\times63&62&16\\
 \Phi_B&252\times81&80&24
\end{array}.
\]
These numbers are the same in the three cases.  The checker constructs
each entry by the displayed wedge formulas and verifies both ranks by
exact Gaussian elimination.

\emph{Reduction to weight subspaces.}
Let $(\mathbb G_m)^6$ act with weights $e_r+f_c$ on both $x_{rc}$ factors,
and weight $-(e_r+e_s+f_c+f_d)$ on the corresponding $a_q$.
It fixes $T$ and preserves $U$.  The closed rank conditions
\eqref{eq:P15-A}--\eqref{eq:P15-B}, together with $U\subset E$, form a
closed torus-invariant subvariety of $\operatorname{Gr}(13,A\otimes B)$.
Choose an integral cocharacter separating all distinct torus characters
on $A\otimes B$.  A limit at zero of its orbit exists by properness of
the Grassmannian.  The limit still satisfies all conditions, and its
eigenspace decomposition for this cocharacter is the torus weight
decomposition.  Consequently an admissible $E$ may be assumed torus
invariant.

Write $V=(A\otimes B)/U=\bigoplus_wV_w$.  For each weight $w$ the
ambient coordinate block consists of the pairs $(q,v)$ with
\[
 w=\operatorname{wt}(x_v)-\operatorname{wt}(Q_{\mathcal S(q)}).
\]
The image $U$ has one dimension in weight $-\operatorname{wt}(x_v)$ for
each of the nine variables, and zero dimensions in the other weights.
In an image-weight block, omit its final coordinate in lexicographic
$(q,v)$ order; the coefficient of the omitted coordinate in the image
vector is nonzero.  The remaining coordinates give a specified basis of
$V_w$.  In a nonimage block take every coordinate.  Thus no choice of
quotient representatives is left unspecified.

For $\mathcal S_A$, there are $33$ nonzero quotient blocks: $15$ of
dimension one, $15$ of dimension two, and three of dimension three.
For $\mathcal S_B$, the corresponding numbers are $34=16+16+2$.
The total quotient dimension is $54$ in both cases.
An admissible extension is therefore specified by subspaces
$F_w\subset V_w$ with $\sum_w\dim F_w=4$.

\emph{The rank formula in quotient weight blocks.}
For either orientation $o\in\{A,B\}$, let $K_o$ be the image of the
baseline map on $U$.  Quotient each codomain weight block by its part of
$K_o$.  In the checker this quotient is represented by an integer-row
basis $H$ for the left nullspace of the baseline block $M_0$.  It checks
$HM_0=0$, $\operatorname{rank}H=\#\operatorname{rows}H$, and
$\#\operatorname{rows}H=\#\operatorname{rows}M_0-\operatorname{rank}M_0$.
Thus $\ker H=\operatorname{im}M_0$, and this is an exact quotient map.
For a factor basis vector with weight $\lambda_j$, write
\[
 L^o_{w,j}:V_w\longrightarrow
       (\operatorname{codomain}\Phi_o/K_o)_{w+\lambda_j}
\]
for the resulting linear map.  Its matrix is $H$ times the explicitly
constructed wedge columns.  For any chosen $F_w$ the exact rank increment
is
\begin{equation}
 \Delta_o(E)=
 \sum_{\gamma}\dim\left(
       \sum_{w+\lambda_j=\gamma}L^o_{w,j}(F_w)\right).
 \label{eq:P15-increment}
\end{equation}
Different target weights form a direct sum.  Formula
\eqref{eq:P15-increment} proves both the weight-count bounds used below
and the validity of subsequently quotienting by already fixed components.

\emph{Initial saturation.}
A support is a set $S$ of at most four quotient weights whose capacities
sum to at least four.  There are $46\,119$ eligible supports for
$\mathcal S_A$ and $51\,985$ for $\mathcal S_B$.
Replace all $F_w$, for $w\in S$, by the whole $V_w$, obtaining a maximal
space $E_S$.  If an admissible $E$ has support $S$, then the kernels of
its wedge maps inject into those for $E_S$.  Equations
\eqref{eq:P15-A}--\eqref{eq:P15-B} require both kernels for $E$ to have
dimension at least $13$, so a kernel of dimension below $13$ for $E_S$
excludes $S$.
The full wedge matrices on $A\otimes B$ have sizes $189\times441$ and
$252\times567$.  Decomposing them into weight blocks and precomputing
the ranks for every local support of at most four weights makes this
finite saturation calculation small.  Every support of size at most two
is excluded.  The surviving supports, before further tests, are
\[
\begin{array}{c|r|r|r}
 &\text{eligible supports}&\text{three weights}&\text{four weights}\\\hline
 T^{-}_{\mathcal S_A}&46\,119&236&13\,019\\
 T^{+}_{\mathcal S_A}&46\,119&244&13\,518\\
 T^{+}_{\mathcal S_B}&51\,985&200&11\,783
\end{array}.
\]
There are only two possible remaining multiplicity patterns:
$(1,1,1,1)$ and $(2,1,1)$.

\emph{Completeness of the linear-arrangement enumeration.}
For linear maps $C_1,\ldots,C_N$ on $k^d$, start with the zero equation
space and repeatedly add the row space of one $C_j$, retaining each
result of rank below $d$ and identifying equal row spaces by exact
row reduction.  This enumerates every possible intersection
$K=\bigcap_{j\in Z}\ker C_j\ne0$.
To such a $K$ assign the indices for which $C_j|_K=0$.
Its actual stratum consists of the nonzero vectors of $K$ on which every
other $C_j$ is nonzero.  It is nonempty: each prohibited additional zero
locus is a proper linear subspace of $K$, and a finite union of proper
linear subspaces cannot cover a nonzero vector space over the infinite
field $k$.
Conversely, the exact zero pattern of any nonzero vector occurs in this
list, by intersecting all the kernels that contain that vector.  Thus
these are locally closed strata; they are not being replaced silently
by their linear spans.

If a subsequent quotient produces new kernel conditions inside $K$, an
intersection $K'\subset K$ is compatible with the original stratum
exactly when none of the originally required nonzero maps vanishes
identically on $K'$.  The same finite-union argument proves both directions.
The checker explicitly makes this compatibility test whenever fixed
components have been quotiented out.

\emph{Four weight lines.}
For a line $F_w=kx$ the characters $w+\lambda_j$, as $j$ varies in one
orientation, are distinct.  Hence its contribution has one dimension
precisely for those $j$ with $L^o_{w,j}x\ne0$.
The combined $7+9$ column-vanishing conditions are linear equations in
at most three coordinates.  The arrangement procedure above therefore
lists every exact joint zero pattern in both orientations.
For several lines, the number of distinct active output weights is a
lower bound for \eqref{eq:P15-increment}.
A preliminary bound sums each line's minimum increment and subtracts
the largest possible loss from repeated output weights; this is valid
because each column has rank at most one.

If an arrangement stratum has a one-dimensional kernel, its line is
fixed.  Compute the actual sum of the images of all fixed lines, and
exclude the case if its ranks already exceed $16$ or $24$.
Otherwise quotient by those fixed images.  In every case that survives
this step there are one or two parameter lines.
For one line, apply the arrangement procedure to its columns in the
new quotient, with the compatibility test just proved.  The rank
increment equals the number of its nonzero columns.
For two lines, do this for each line and use the number of distinct
active weights in their union.  Every resulting case is excluded.
Here are the complete counts; all ranks in the fixed-image step are
computed over $\mathbb Q$:
\[
\begin{array}{c|r|r|r}
 &T^-_{\mathcal S_A}&T^+_{\mathcal S_A}&T^+_{\mathcal S_B}\\\hline
 \text{supports excluded by minimum/overlap}&12\,307&12\,798&11\,519\\
 \text{further supports excluded by joint weight unions}&176&180&36\\
 \text{joint four-line states remaining}&7\,225&7\,697&2\,182\\
 \text{states excluded by fixed-image ranks}&6\,619&7\,057&1\,972\\
 \text{one-parameter-line states, all excluded}&592&626&210\\
 \text{two-parameter-line states, all excluded}&14&14&0\\
 \text{refined two-line state pairs, all excluded}&32&22&0
\end{array}.
\]

\emph{A double component and two lines.}
For a two-dimensional quotient block, its double component is fixed.
For a three-dimensional block, a two-plane has the unique form
$\ker n$ for a nonzero dual normal $n$, modulo scaling.
Let $L:k^3\to D$ be any one of its column maps.
The complete rank-drop rule is
\[
\begin{array}{c|c|l}
 \operatorname{rank}L&\operatorname{rank}(L|_{\ker n})&\text{condition}\\\hline
 0&0&\text{always}\\
 1&0&n\text{ spans the row space of }L\\
 1&1&\text{otherwise}\\
 2&1&n\text{ annihilates }\ker L\\
 2&2&\text{otherwise}\\
 3&2&\text{always}.
\end{array}
\]
Indeed, in rank one the restriction is zero exactly when
$\ker n=\ker L$; in rank two the one-dimensional kernel of $L$ lies in
$\ker n$ exactly in the stated case; the other two rows are immediate.
Both nontrivial drop conditions are linear equations in $n$.
For rank one, equations cutting out the line $\operatorname{row}L$ are
obtained from a basis of its annihilator.  For rank two, the equation is
evaluation on a generator of $\ker L$.

The same arrangement algorithm therefore lists every possible rank
pattern of the double component, jointly in the two orientations.
Combine it with the two line patterns.  At each output weight, the
dimension of the sum in \eqref{eq:P15-increment} is at least the largest
of the dimensions contributed by the three components.  Summing these
maxima excludes every case:
\[
\begin{array}{c|r|r}
 &\text{arrangement states tested}&\text{states excluded}\\\hline
 T^-_{\mathcal S_A}&42\,274&42\,274\\
 T^+_{\mathcal S_A}&45\,489&45\,489\\
 T^+_{\mathcal S_B}&31\,796&31\,796.
\end{array}
\]
This completes the enumeration of every torus-fixed four-dimensional
extension quotient.  There is no admissible $E$, contradicting the
necessary condition for border rank at most $13$.  The three asserted
lower bounds follow.\qed

\paragraph{Algebraic-geometric dictionary for P-15.}
The tensor and both wedge maps are maps between finite free
$\mathbb Z$-modules.  Their rank conditions are determinantal closed
subschemes of the relevant relative Grassmannian.  The incidence
condition $U\subset E$ is a closed condition on the tautological bundle;
the Grassmannian's projectivity gives the proper projection used in R-04
and the cocharacter specialization used here.  The acting split torus is
the group scheme $(\mathbb G_m)^6$, and its character decomposition is
an integral representation-ring decomposition.  Exterior powers in the
two wedge maps are the usual $\lambda$-operations on locally free sheaves.
The locally closed arrangement strata are obtained by intersections of
linear closed subschemes and complements of further such subschemes.
These are precisely the scheme, group-scheme, and $K_0$ constructions
identified in the paper's common Grothendieck locator table.
The certificate uses characteristic-zero rational linear algebra;
no statement about bad primes, purity, an $\ell$-adic weight filtration,
or a motivic realization is inferred from it.

\section{The determinant seventh minimum and the bound forty-three}

\begin{finding}{R-08}{the exact determinant seventh minimum}
Over every algebraically closed field of characteristic zero,
\[
 \mu_7(U^{\det}_{4,2})=14,\qquad
 \BR(\mathcal D_4)\geq43.
\]
In particular the usual complex symmetric border rank of the
four-by-four determinant polynomial is at least forty-three.
This is a lower bound for the full polynomial tensor, not an
assertion that its full border rank equals forty-three.
\end{finding}
\begin{proof}
By P-01 and P-07, any seven-plane in the determinant residual
space degenerates, without increasing border rank, to a
coordinate order ideal in the product of the two componentwise
posets of two-subsets of $\{0,1,2,3\}$. There are exactly
$25$ such seven-element ideals. The independent enumeration
starts with the empty ideal and, at each of seven stages,
adjoins an element only when all its strict predecessors
already occur. Every ideal is reached, since any nonempty
finite ideal has a maximal element whose removal is an ideal.
The final set is compared exactly with the certificate list;
no sampling of Grassmannian points or of coordinate subsets
is used.

All $25$ cases have border rank at least fourteen, as follows.
\begin{center}
\begin{tabular}{p{0.39\linewidth}|r|p{0.40\linewidth}}
Certificate type&Ideals&Lower certificate\\\hline
Ordinary support&2&Fourteen independent support columns\\
Koszul map on the seven-dimensional factor&7&Rank $>13\binom6p$\\
Koszul map on the variable-support factor&2&Rank $6501>13\cdot495$\\
P-09&4&Two prolongation types and their transposes\\
P-12&4&Two further types and their transposes\\
P-14&4&The support-eleven types and their transposes\\
P-15&2&The support-nine type and its transpose
\end{tabular}
\end{center}
The two support-factor Koszul maps have $m=13$, exterior
degree $p=4$ and pure-tensor cost $\binom{12}4=495$.
They are decomposed into row/column torus weight blocks.
Within each block a specified nonzero minor gives a rank
lower bound; distinct source and target weights ensure
that these bounds add. The seven first-factor Koszul
certificates each specify a minor of order
$13\binom6p+1$, with $p=2$ or $3$.
The pure-tensor cost and closed determinantal-rank argument
are the same as in the earlier Koszul findings.

The complete list, all minor indices, all weights and the
finding assigned to each prolongation case are in
\path{det7_global_certificate.json}.
The independent \path{verify_det7_global.py} reconstructs
all $25$ ideals, their determinant tensors, all selected
Koszul entries and the support minors directly. It verifies
$2997$ nonzero minors modulo the verified prime $1000003$,
as well as disjointness and correct weights of the support-factor
blocks. These certify nonzero integer minors, hence the stated
characteristic-zero ranks. Its fourteen prolongation entries
are matched with the exact subsets of the separately checked
P-09, P-12, P-14 and P-15 certificates. Thus every ideal has
the required lower bound, proving $\mu_7\geq14$.

For the upper bound, take seven edges of a $3\times3$
anchored shape meeting all three row and column vertices,
for example all six edges in its first two rows and one
edge in its third row. P-06 gives border rank
$7+3+3+1=14$. This shape embeds in the four-by-four residual
space, so $\mu_7\leq14$.

The complementary-minor image has dimension $36$.
Applying R-01 while retaining seven dimensions gives
\[
 \BR(\mathcal D_4)\geq36-7+14=43.
\]
The symmetric atom variety is contained in the CP atom
variety; consequently this CP lower bound also bounds
symmetric border rank. In characteristic zero the integral
coefficient tensor differs from the normalized polynomial
tensor by a nonzero scalar, which does not change either
border rank.
\end{proof}

\begin{finding}{R-09}{the determinant eighth and ninth minima}
Over every algebraically closed field of characteristic zero,
\[
 \mu_8(U^{\det}_{4,2})=15,\qquad
 \mu_9(U^{\det}_{4,2})=16.
\]
Both yield the same full-tensor lower bound forty-three.
\end{finding}
\begin{proof}
R-07 and R-08 give $\mu_j\geq j-7+14=j+7$ for $j\geq7$.
Anchored $3\times3$ shapes with eight and nine edges meet
all their vertices. P-06 gives matching upper bounds
$j+3+3+1=j+7$. Substitution gives $36-j+(j+7)=43$.
\end{proof}

\paragraph{Geometric dictionary.}
The relative Grassmannian and its torus and root-subgroup
specializations are those of P-01/P-07. The finite order-ideal
cover identifies all possible fixed subbundles relevant to
this reduction. Exterior morphisms of the tautological
bundles supply the Koszul minors, and the proper prolongation
incidences supply the other cases. R-08 retains these maps
and their determinantal ideals, not only their representation
classes in $K_0$. R-09 is the flag-incidence substitution
of R-01 followed by the integral Segre curves of P-06.
The precise EGA~II, SGA~3, SGA~6 and FGA foundations and
archive limitations are recorded in the separate dictionary.

% P-16: exact certificates for three permanent support-eleven orbits.
\begin{finding}{P-16}{three permanent seven-plane orbits with eleven variables}
Over every algebraically closed field of characteristic zero, the
permanent residual tensors for the following subsets have border rank
at least fourteen:
\[
 S_1=\{0,1,2,3,6,7,9\},\qquad
 S_2=\{0,1,2,3,6,7,18\},\qquad
 S_3=\{0,1,3,6,7,15,21\}.
\]
The assertion also holds after row or column relabeling or transposition.
It concerns these three coordinate orbits.
\end{finding}
\begin{proof}
Each tensor has format $7\times11\times11$ and last-factor image
$U\subset A\otimes B$ of dimension eleven. The proper Grassmannian
incidence and torus specialization proved in P-09 imply that border
rank at most thirteen would give a torus-invariant thirteen-plane
$L\supset U$ satisfying
\[
 \rank\Psi_A(L)\le78,\qquad \rank\Psi_B(L)\le130.
 \tag{P16.1}
\]
The quotient $L/U$ is a two-plane. Its weight support therefore consists
of two different weight lines or a two-plane in one quotient weight
block of dimension at least two. The complete finite enumeration is:
\[
\begin{array}{c|r|r|r|r|r|r}
 &\text{weights}&\text{choices}&\text{saturation}
 &\text{base }\Psi_A&\text{pairs}&\text{increment}\\\hline
 S_1&49&1189&1176&74&13&6\\
 S_2&48&1143&1075&71&67&8\\
 S_3&49&1190&1173&72&17&8
\end{array}
\]
The numbers of choices are, respectively,
$\binom{49}{2}+13$, $\binom{48}{2}+15$, and
$\binom{49}{2}+14$. For every saturation exclusion, the certificate
adjoins all coordinate directions in the chosen quotient weight block
or blocks. It exhibits an exterior-map minor forcing the kernel of
this larger space to have dimension at most twelve. The kernel for
$L$ injects into that larger kernel, contradicting the necessary
kernel dimension at least thirteen. This is precisely the saturation
criterion proved in P-09, with permanent signs.

All surviving choices are pairs of different weight lines. For each
weight line, integer annihilators of the baseline exterior image give
linear conditions for a projected column to vanish. The certificate
checks all required stacks of these conditions, obtaining a uniform
lower bound for the number of nonzero projected columns. Within a line
their target weights are distinct, so these columns are independent.
For a pair of lines, subtracting the number of common target weights
from the sum of the two bounds gives a valid joint increment. This
excludes the $13$, $67$, and $17$ pairs displayed in the table: the
resulting $A$-map rank is always at least $79$, contrary to (P16.1).
It finishes $S_1$ and $S_3$ and leaves one pair for $S_2$.

For that last pair the quotient blocks have indices $11,14$ in the
specified certificate ordering, and dimensions one and two. The first
line is forced, with coordinate representative
$X=a_1\otimes b_0$, where $b_0=x_{00}$. Thus every possible extension
contains the fixed twelve-plane
\[
 U_0=U+kX.
\]
The second quotient line has two parameters. Its weight differs from
that of $X$, so adjoining $X$ does not change its chosen parameter basis.
An exact minor gives $\rank\Psi_B(U_0)\ge125$. The certificate supplies
integer annihilators of this enlarged image. Every stack for six of
the eleven new $B$-columns has full column rank two: all
$\binom{11}{6}=462$ stacks are checked. Consequently a nonzero parameter
vector makes at most five columns vanish. The eleven target weights
are distinct, so the rank increment is at least six, giving
\[
 \rank\Psi_B(L)\ge125+6=131>130.
\]
This excludes the last pair and completes all $3522$ choices.

The data are in \path{per7_m11_certificate.json}. The independent
\path{verify_seven_plane_certificate.py} reconstructs the permanent
coefficients, quotient weight blocks, exterior maps, saturation spaces,
and every selected matrix entry. It verifies the $3424$ saturation
minors, all baseline minors, all integer annihilator identities, all
$1939$ linear zero-pattern tests, and the complete coverage of the choices.
It also checks that the conditional fixed block has dimension one and
that its new annihilators vanish on $\Psi_B(U_0)$ itself. All claimed
nonzero minors are checked modulo the verified prime $1000003$,
certifying nonzero integer minors and hence the stated characteristic-zero
ranks. The integer and modular overflow bounds are those of P-09.
Row and column permutations and transposition are invertible product
maps, giving the assertion for the three entire coordinate orbits.
\end{proof}

\paragraph{Geometric dictionary for P-16.}
The proof uses the projective Grassmannian incidence, split-torus
specialization, quotient weight bundles, and exterior morphisms from
P-09. The last pair is the closed incidence stratum containing the
fixed intermediate subbundle $U_0$, followed by empty projective linear
loci for the remaining quotient line. The maps and their integer
determinantal ideals are retained, rather than only their classes in
$K_0$. The EGA, SGA and FGA locators are those in the common dictionary
for these same constructions. No unipotent invariance of the permanent
is used.

\subsection{The last permanent seven-plane: a second prolongation}
\begin{finding}{P-19}{the final permanent seven-plane has border rank at least fourteen}
Over every algebraically closed field of characteristic zero, the
permanent residual tensor for the rectangle subset
\[
 S=\{0,1,2,6,7,18,19\}
\]
has border rank at least $14$.
\end{finding}
\begin{proof}
We give the additional closed necessary condition, the complete reduction
to eight fixed extensions, and their integer certificates.
Use $A=k^7$ and $B=C=k^{11}$, with respective bases
$a_0,\ldots,a_6$ and $x_0,\ldots,x_{10}$.  The variables are the entries
of a three-by-four rectangle, numbered in row order, with $x_{11}$
omitted.  For compactness write $xy$ for $x\otimes y+y\otimes x$ in
$B\otimes C$, throughout this proof.  The tensor is
\[
 T=\sum_{q=0}^6a_q\otimes Q_q,
 \qquad
 \begin{aligned}
 Q_0&=x_0x_5+x_1x_4,& Q_1&=x_0x_6+x_2x_4,\\
 Q_2&=x_0x_7+x_3x_4,& Q_3&=x_0x_9+x_1x_8,\\
 Q_4&=x_0x_{10}+x_2x_8,& Q_5&=x_4x_9+x_5x_8,\\
 Q_6&=x_4x_{10}+x_6x_8.
 \end{aligned}
\]
Thus the displayed notation does not presume a decomposition with equal
second and third factors.  Let $U=T(C^*)\subset A\otimes B$;
its eleven contraction vectors are independent.

\paragraph{A closed higher-prolongation necessary condition.}
For a subspace $E\subset A\otimes B$, use the symmetric-tensor
subspaces to define
\begin{align*}
 \mathcal P_{21}(E)&=(A\otimes E)\cap(\operatorname{Sym}^2A\otimes B),\\
 \mathcal P_{12}(E)&=(E\otimes B)\cap(A\otimes\operatorname{Sym}^2B),\\
 \mathcal P_{22}(E)&=(A\otimes E\otimes B)
               \cap(\operatorname{Sym}^2A\otimes\operatorname{Sym}^2B).
\end{align*}
Factor permutations place these intersections in $A^{\otimes2}\otimes B$,
$A\otimes B^{\otimes2}$, and $A^{\otimes2}\otimes B^{\otimes2}$,
respectively.  In the last expression $E$ occupies the second $A$ and
first $B$ positions.
For each fixed dimension $r$ of $E$, each condition
$\dim\mathcal P_{ij}(E)\ge r$ is closed on
$\operatorname{Gr}(r,A\otimes B)$.  Indeed it is a kernel-dimension
condition for the map from the corresponding tensor product of the
rank-$r$ tautological bundle to the fixed ambient tensor space modulo
the displayed symmetric subspace.  Both source and target have fixed
ranks, so it is determinantal.

If $E$ is spanned by $r$ independent pure vectors
$u_\ell\otimes v_\ell$, its prolongations contain, respectively,
\[
 u_\ell^{\otimes2}\otimes v_\ell,\qquad
 u_\ell\otimes v_\ell^{\otimes2},\qquad
 u_\ell^{\otimes2}\otimes v_\ell^{\otimes2}
 \quad(1\le\ell\le r).
\]
Each of these $r$-tuples is independent.  Choose linear forms nonzero
on all the finitely many $u_\ell$ and $v_\ell$, and contract the added
factors; the result is the original independent tuple multiplied by
nonzero scalars.  Such forms exist because $k$ is infinite.
The open condition that the $r$ pure $AB$ vectors be independent is
nonempty and dense in the parameter space of $r$ pure $ABC$ summands
when $r\le\dim(A\otimes B)$.
Consequently, impose $T(C^*)\subset E$ and all three closed
prolongation conditions above in the product with the projective
Grassmannian.  Its projection is closed and contains the dense
parametrized rank-$r$ set.  This is the same proper-incidence argument
as in R04, now with the additional $\mathcal P_{22}$ condition.
It proves the needed implication
\begin{equation}
 \underline R(T)\le13\ \Longrightarrow\
 \exists E\supset U,\quad \dim E=13,\quad
 \dim\mathcal P_{21}(E),\dim\mathcal P_{12}(E),
       \dim\mathcal P_{22}(E)\ge13.
 \label{eq:P19-necessary}
\end{equation}
In particular, the third condition is imposed directly on the fixed
ambient intersection.  No assertion of continuity for a varying
first-prolongation kernel is used.

\paragraph{The complete first-stage enumeration.}
Apply the row/column split torus with negative quadratic weights on
$A$ and variable weights on $B,C$, as in P-15.  The closed incidence
in \eqref{eq:P19-necessary} is torus invariant.  A separating
cocharacter limit therefore makes $E$ torus invariant without losing
any of its conditions.  The quotient $(A\otimes B)/U$ has $47$
nonzero weight blocks and dimension $66$.  Fourteen blocks have
capacity at least two.  Since $\dim(E/U)=2$, there are
\[
 14+\binom{47}{2}=1095
\]
possible supports, consisting of a two-plane in one weight or lines
in two different weights.

The checker \texttt{check\_per7\_last\_certificate.py} constructs these
blocks in lexicographic $(q,v)$ order.  As in P-15, in every image-weight
block it omits the last coordinate, whose image coefficient is nonzero,
to specify a basis of the quotient.  The baseline wedge matrices are
\[
\begin{array}{c|c|c|c}
 &\text{matrix on }U&\text{rank}&\text{maximum allowed increment}\\\hline
 \Phi_A&231\times77&70&8\\
 \Phi_B&385\times121&115&15.
\end{array}
\]
The allowed total ranks are $78$ and $130$, respectively.
Saturating each chosen support by its entire weight blocks, a
prolongation kernel of dimension at most twelve excludes $971$ of the
$1095$ supports.  Every surviving support consists of two distinct
weight lines; there are $124$ such pairs.

For each line, the projected columns in one orientation have pairwise
distinct output weights.  The complete linear-arrangement algorithm
proved in P-15 gives the uniform minimum number of nonzero columns.
For a pair, the sum of those minima minus the number of common output
weights is a lower bound for the rank increment.  This excludes another
$116$ pairs.  The eight remaining pairs, in the fully specified block
numbering of the checker, are
\[
 (1,23),(5,23),(19,23),(20,23),(21,23),(22,23),(23,24),(23,25).
\]
All the saturation ranks, quotient annihilators, and line arrangements
in these two stages are recomputed over $\mathbb Q$ by the attached
self-contained checker; no earlier search file or numerical rank is
used.

\paragraph{Joint refinement gives eight fixed extensions.}
Block $23$ is one dimensional and has the coordinate representative
$a_2\otimes x_8$.  Thus every remaining extension contains
\[
 U_0=U+\langle a_2\otimes x_8\rangle,\qquad \dim U_0=12.
\]
Its wedge-map ranks are $72$ and $123$.  The remaining line may
therefore add at most six $A$-columns and seven $B$-columns.
Quotient by the actual images on $U_0$, and apply the joint
$7+11$-column arrangement enumeration.  The required zero counts are
at least one $A$-column and four $B$-columns.

Only blocks $1$ and $5$ have parameters; their quotient dimensions
are four, with the following ordered coordinate representatives:
\[
\begin{array}{c|l|c}
 \text{block}&\text{representatives}&\text{only surviving parameter line}\\\hline
 1&a_0x_1,\ a_1x_2,\ a_2x_3,\ a_5x_9&(0,0,1,0)\\
 5&a_0x_5,\ a_1x_6,\ a_2x_7,\ a_3x_9&(0,0,1,0).
\end{array}
\]
Here $a_qx_v$ means $a_q\otimes x_v$.  In each block the exact joint
linear arrangement has six strata, of which the displayed line alone
satisfies both rank conditions.  This uses the entire projective
parameter space; the other five strata have not been replaced by
sampled parameter values.  The six other blocks are already one
dimensional, with representatives
\[
 a_2x_1,\ a_2x_2,\ a_2x_5,\ a_2x_6,\ a_2x_9,\ a_2x_{10}
 \quad\text{for blocks }19,20,21,22,24,25.
\]
It follows that every torus-invariant $E$ satisfying the two first
prolongation conditions is one of the following eight fixed spaces:
\begin{equation}
 E_v=U+\langle a_2\otimes x_8,\ a_2\otimes x_v\rangle,\qquad
 v\in\{1,2,3,5,6,7,9,10\}.
 \label{eq:P19-eight}
\end{equation}
The joint linear equations and their unique surviving parameter
vectors are included explicitly in \texttt{per7\_last\_certificate.json}.

\paragraph{The eight degree $(2,2)$ obstructions.}
For each space in \eqref{eq:P19-eight}, the map
\[
 \Phi_B:E_v\otimes B\longrightarrow A\otimes\Lambda^2B
\]
is a $385\times143$ integer matrix $M_{12,v}$.
The certificate provides a $130\times130$ unit minor and an integer
matrix $N_v$ of size $143\times13$ satisfying
\[
 M_{12,v}N_v=0,
\]
together with a $13\times13$ unit minor of $N_v$.
Thus $\operatorname{rank}M_{12,v}=130$, and the columns of $N_v$ form
an exact basis of its kernel $\mathcal P_{12}(E_v)$.

We use the symmetric-tensor basis
$x_i\otimes x_i$ and
$x_i\otimes x_j+x_j\otimes x_i$ for $i<j$.
The checker embeds each kernel vector in $A\otimes B\otimes B$,
verifies symmetry under exchanging its two $B$ indices, and reads its coefficients
in this basis.  There are $\binom{12}{2}=66$ symmetric coordinates.
The resulting thirteen matrices in $A\otimes\operatorname{Sym}^2B$
give the next wedge map
\[
 M_{22,v}:A\otimes\mathcal P_{12}(E_v)
       \longrightarrow\Lambda^2A\otimes\operatorname{Sym}^2B,
 \qquad \text{matrix size }1386\times91.
\]
Its kernel is exactly $\mathcal P_{22}(E_v)$: first intersect
$A\otimes E_v\otimes B$ with $A^{\otimes2}\otimes\operatorname{Sym}^2B$,
then impose symmetry in the two $A$ factors.  The two successive
intersections equal the defining simultaneous intersection.

Each $M_{22,v}$ has a certified $79\times79$ unit minor.
For clarity, all twenty-four final integer determinants are listed here,
with the selected minor order in square brackets:
\[
\begin{array}{c|r|r|r}
 v&\det(M_{12,v}[130])&\det(N_v[13])&\det(M_{22,v}[79])\\\hline
 1&1&1&1\\
 2&-1&1&1\\
 3&1&1&1\\
 5&-1&1&-1\\
 6&1&1&1\\
 7&-1&1&1\\
 9&1&1&1\\
 10&-1&1&1
\end{array}.
\]
The certificate contains every selected row and column and every entry
of $N_v$.  The checker evaluates the determinants over $\mathbb Z$;
its auxiliary modular elimination only selects the minors.
Consequently
\[
 \dim\mathcal P_{22}(E_v)
   \le91-79=12<13
\]
for each of the eight spaces.  This contradicts
\eqref{eq:P19-necessary} in every possible case and proves the finding.
\end{proof}

\paragraph{Scope and dictionary.}
The last twenty-four minors are units, but the preceding finite
arrangement reduction was checked over $\mathbb Q$.  P-19 therefore
claims characteristic zero, without an unverified extension to every
prime.  Geometrically it uses determinantal closed subschemes of a
relative Grassmannian, its tautological bundle, and a proper incidence
projection, with a split-torus specialization.  The extra condition is
the next symmetric prolongation of that bundle; it is a kernel rank
condition in a fixed ambient locally free sheaf.  Exterior and symmetric
powers supply the standard $\lambda$-ring and Schur-functor dictionary
at the common SGA/FGA locators of the paper.  No purity, $\ell$-adic
weight, or motivic claim is needed to transfer this proof.

\section{The permanent seventh minimum and the bound forty-three}

\begin{finding}{R-10}{the exact permanent seventh, eighth and ninth minima}
Over every algebraically closed field of characteristic zero,
\[
 \mu_7(U^{\mathrm{per}}_{4,2})=14,\qquad
 \mu_8(U^{\mathrm{per}}_{4,2})=15,\qquad
 \mu_9(U^{\mathrm{per}}_{4,2})=16,
 \qquad \BR(\mathcal P_4)\geq43.
\]
The same lower bound holds for the usual complex symmetric border
rank of the four-by-four permanent polynomial.
\end{finding}
\begin{proof}
P-01 reduces the seventh Grassmannian minimum to coordinate
seven-planes among the $36$ complementary permanent minors.
There are $\binom{36}{7}=8\,347\,680$ such subsets. For a subset
$S$, let $m(S)$ be the number of matrix-entry variables appearing
in its seven rectangular supports. The ordinary support flattening
has rank $m(S)$. Indeed, for each appearing variable $x$, choose
one rectangle $s$ containing it. The column indexed by $s$ and
the diagonally opposite variable has a single unit in row $x$;
these chosen columns are distinct for distinct $x$. They form
a unit minor of order $m(S)$. Thus $m(S)\geq14$ gives the required
lower bound directly.

Row and column permutations, together with transposition, act on
the $36$ rectangles and preserve the permanent tensor up to
invertible changes of bases. The resulting group has order
$24^2\cdot2=1152$. The remaining $1\,543\,608$ subsets, those with
$m(S)<14$, form exactly $1517$ orbits. The certificate proves this
finite cover rather than relying on a classification assertion:
it constructs all $1152$ actions, verifies that the $1517$ listed
orbits are pairwise disjoint and have the recorded sizes, and
independently enumerates all $\binom{36}{7}$ subsets to check that
every subset of support less than fourteen belongs to their union.

The lower certificates for these $1517$ representatives are:
\begin{center}
\begin{tabular}{p{0.39\linewidth}|r|p{0.40\linewidth}}
Certificate type&Orbits&Lower certificate\\\hline
Koszul map on the seven-dimensional factor&1503&Rank $>13\binom6p$\\
Koszul map on the variable-support factor&3&Rank $>13\cdot495$\\
P-13&5&Support-twelve prolongation exclusions\\
P-15&2&Support-nine prolongation exclusions\\
P-16&3&Support-eleven prolongation exclusions\\
P-19&1&The last higher-prolongation exclusion
\end{tabular}
\end{center}
For the three support-factor Koszul maps, $m=13$, the exterior
degree is $p=4$, and the rank of the map on a pure tensor is
$\binom{12}{4}=495$. Their certified ranks are respectively
$6457$, $6457$ and $6515$, all strictly greater than
$13\cdot495=6435$. Each rank is a sum of certified ranks of
distinct torus weight blocks. For the first-factor maps, each
listed minor has order strictly greater than thirteen times
the corresponding pure-tensor cost. Determinantal closedness
therefore gives border rank at least fourteen in all these cases.
The eleven remaining representatives are matched, as exact
sorted subsets of rectangle indices, with the separately proved
P-13, P-15, P-16 and P-19 statements.

The file \path{per7_global_certificate.json} records the full
cover, weights and minor indices. The independent verifier
\path{verify_per7_global.py} reconstructs the tensors and all
selected matrix entries, using only the audited arithmetic and
Koszul-entry routines in \path{verify_det7_global.py}.
It verifies $6041$ nonzero minors modulo the prime $1000003$
(whose primality is checked by trial division), including
disjointness and the asserted weights of the blocks. A nonzero
residue proves that the corresponding integer determinant is
nonzero, hence gives the claimed rank over every field of
characteristic zero. This statement does not infer the same
rank in arbitrary positive characteristic from a modular
calculation. The exhaustive cover and the eleven separately
checked prolongation proofs establish $\mu_7\geq14$.

An anchored $3\times3$ shape with seven edges meeting every
row and column vertex has border rank $7+3+3+1=14$ by P-06.
It is a coordinate seven-plane in this residual space, so
$\mu_7\leq14$. R-07 now gives $\mu_j\geq j+7$ for $j\geq7$.
The eight- and nine-edge anchored shapes give equality for
$j=8,9$, again by P-06. Finally, the residual image has
dimension $36$, and R-01 gives
\[
 \BR(\mathcal P_4)\geq36-j+\mu_j=43\qquad(j=7,8,9).
\]
In characteristic zero the integral coefficient tensor differs
from the normalized polynomial tensor by a nonzero scalar.
The symmetric atom variety is contained in the CP atom variety,
so the CP lower bound also bounds symmetric border rank.
\end{proof}

\paragraph{Geometric dictionary.}
R-10 uses the same relative Grassmannian, torus degeneration,
exterior morphisms of locally free sheaves and proper
prolongation incidences as the preceding findings. For the
permanent, the finite symmetry group and exhaustive coordinate
cover replace the determinant's Borel order-ideal reduction;
no characteristic-zero Borel invariance of the permanent is
assumed. The minors are functions defined over $\Z$, and their
nonzero residues certify nonvanishing over $\Q$ and then over
every characteristic-zero extension. The maps and their
determinantal ideals, rather than their $K_0$ classes alone,
carry the rank information. Exact foundational locators and
the absence of an established archival cote for these modern
calculations are recorded in the separate dictionary.

\section{Uniform seventh, eighth and ninth polynomial minima}

\begin{finding}{R-11}{the seven-plane improvement in every matrix size}
Over every algebraically closed field of characteristic zero,
for every $n\geq4$ and for either polynomial family,
\[
 \mu_7(U_{n,2})=14,\qquad
 \mu_8(U_{n,2})=15,\qquad
 \mu_9(U_{n,2})=16.
\]
Consequently
\[
 \BR(\mathcal D_n),\ \BR(\mathcal P_n)
 \ \geq\ \binom n2^2+7\qquad(n\geq4).
\]
This is the consequence of the degree-two residual calculation;
it is not asserted to be the strongest known bound for general $n$.
\end{finding}
\begin{proof}
P-01 reduces the seventh minimum to seven distinct rectangular
quadratic minors or permanents. Let $G$ be the bipartite support
graph: its vertices are the appearing rows and columns and its
edges are the appearing matrix-entry variables. As in R-10,
the support flattening has a unit minor of order $|E(G)|$.
If $|E(G)|\geq14$, the desired lower bound follows.

Suppose therefore $|E(G)|\leq13$. Every edge belongs to one of
the seven selected rectangles. Every appearing row and column
has degree at least two. There are consequently at most six
row vertices and six column vertices. If there are at most four
of each, the tensor embeds, after relabelling and adjoining
unused coordinates, into the size-four situation. R-08 or R-10
then gives the lower bound fourteen. Such adjoining does not
change border rank: inclusion gives one inequality and a split
coordinate projection gives the reverse inequality.

In the remaining cases, transpose if necessary to put the larger
vertex count first. Write the row and column counts as $(r,c)$;
then $r\in\{5,6\}$ and $2\leq c\leq r$. A row is specified by
its nonempty neighbour subset of the $c$ columns. We enumerate
nondecreasing multisets of $r$ such subsets, each of size at
least two, with total size at most thirteen. We discard a
multiset if a column has degree below two, if fewer than seven
rectangles occur, or if some edge belongs to no rectangle.
Finally we identify column permutations, since the sorted row
multisets already identify row permutations. This covers every
possible support graph; it may retain graphs for which no
seven rectangles use every edge, which does not invalidate
the cover.

There are exactly thirty-four retained graph types. For each,
enumerate all seven-subsets of its rectangles whose union is
the whole graph, up to the graph's row/column automorphisms.
The exact cover has the following sizes:
\begin{center}
\begin{tabular}{c|r|r}
$(r,c)$&Support graphs&Seven-subset orbits\\\hline
$(5,2)$&1&4\\
$(5,3)$&6&853\\
$(5,4)$&19&116\\
$(5,5)$&1&6\\
$(6,2)$&1&20\\
$(6,3)$&5&27\\
$(6,4)$&1&1\\\hline
Total&34&1027
\end{tabular}
\end{center}
There are no retained types with $(r,c)=(6,5)$ or $(6,6)$.
The independent verifier enumerates all neighbour multisets
again, compares the complete canonical graph set, constructs
every vertex permutation preserving each graph, and checks
pairwise disjointness of the listed seven-subset orbits.
It then independently tests every seven-subset of that graph
to check exact coverage of those using every edge. A kernel
in the permutation action on rectangles is harmless: the
verifier uses the set of distinct induced actions.

For each family, first-factor Koszul minors give lower bound
fourteen for $1019$ of the $1027$ representatives. The remaining
eight have the following support-factor Koszul certificates.
Seven have thirteen variables; exterior degree one has pure
cost $12$, and the certified ranks are at least $158>13\cdot12$.
The final representative has twelve variables; exterior degree
three has pure cost $\binom{11}{3}=165$, and the ranks are at
least $2150$ for the determinant and $2147$ for the permanent,
both strictly greater than $13\cdot165=2145$.

The full graph cover, selected rectangles, orbit sizes and
minor indices are recorded in \path{uniform7_certificate.json}.
The independent \path{verify_uniform7.py} reconstructs every
selected Koszul entry from its tensor coefficient and exterior
insertion sign. The support-factor maps are decomposed into
distinct row/column torus weight blocks, whose ranks add.
Every listed minor is evaluated modulo the prime $1000003$,
with primality, coordinate ranges, distinct rows and columns,
block weights and the pure-tensor costs checked. Nonzero
residues certify nonzero integer minors and therefore the
required ranks in characteristic zero. The verifier imports
only the independently audited determinant and Koszul-entry
routines from \path{verify_det7_global.py}, not search code.
These finite certificates prove the lower bound in every
remaining case, hence $\mu_7\geq14$ for every $n\geq4$.

The anchored seven-edge construction of P-06 uses only four
rows and four columns and has border rank fourteen. It embeds
in every $n\geq4$, proving equality. R-07 then gives
$\mu_j\geq j+7$ for $j\geq7$, while the anchored eight- and
nine-edge constructions give equality for $j=8,9$.
The flattening image dimension is $\binom n2^2$, so R-01 gives
the asserted full-tensor bounds.
\end{proof}

\paragraph{Geometric dictionary.}
The finite support graph is an index for a coordinate subbundle
in the integral quadratic residual bundle. Split inclusions and
projections carry the small-support tensors between ambient
matrix sizes. P-01 supplies the torus specialization, and
exterior morphisms between locally free sheaves supply the
determinantal conditions. Graph automorphisms act through
permutation matrices, with determinant signs absorbed into
the first-factor basis. The published EGA, SGA~3 and SGA~6
locators are those of the common dictionary. Rank is certified
by the actual morphisms and their minors; no replacement by
a representation-ring class, purity weight or motive is used.

\subsection{All row permanent quadratic minima and a polynomial ten-plane curve}
\begin{finding}{P-20}{row permanent quadratic minima and a ten-plane degeneration}
Let $k$ be any algebraically closed field, let $n\ge2$, and put
\[
 U^{\mathrm{row,per}}_{n,2}
   =\langle e_i\otimes e_j+e_j\otimes e_i:0\le i<j<n\rangle.
\]
For every $1\le j\le\binom n2$, the exact Grassmannian minimum is
\[
 \mu_j(U^{\mathrm{row,per}}_{n,2})=j+1.
\]
In addition, the polynomial permanent coordinate ten-plane with rectangle subset
\[
 S_{10}=\{0,1,2,3,4,6,7,8,9,10\}
\]
has residual border rank at most $18$.  Consequently
\[
 \mu_{10}(U^{\mathrm{per}}_{4,2})\le18
 \quad\text{in every characteristic},
 \qquad
 17\le\mu_{10}(U^{\mathrm{per}}_{4,2})\le18
 \quad\text{in characteristic zero}.
\]
This does not decide between $17$ and $18$.
\end{finding}
\begin{proof}
For a vector $v$ put $q(v)=v\otimes v$ and write
$s_{ij}=e_i\otimes e_j+e_j\otimes e_i$.
We first prove the general row statement.
Every matrix in $U^{\mathrm{row,per}}_{n,2}$ is symmetric with zero
diagonal.  A nonzero rank-one symmetric matrix has the form
$\lambda vv^{\mathsf t}$, since symmetry makes its two rank-one
factors proportional.  Its diagonal cannot vanish over a field unless
the matrix is zero.  Hence every nonzero matrix in every $j$-plane
$W\subset U^{\mathrm{row,per}}_{n,2}$ has rank at least two.
The residual inclusion tensor has first-factor dimension $j$, so
R-01 with one retained dimension gives
$\underline R(P_W)\ge j-1+2=j+1$.  This argument is valid also in
characteristic two.

For the matching upper bound, choose a graph $G$ of $j$ edges containing
a full star centered at vertex $0$ on its used vertices
$0,\ldots,v-1$.  Such a graph exists for every allowed $j$:
if $j<n-1$, take $v=j+1$ and the star itself; otherwise take $v=n$,
all $n-1$ star edges, and any $j-(n-1)$ of the other edges.
Give its edges independent first-factor labels $\gamma_{ij}$ and let
$P_G=\sum_{ij\in G}\gamma_{ij}\otimes s_{ij}$.
For a star edge use
\[
 F_i(t)=\frac{q(e_0+t e_i)-q(e_0)}t=s_{0i}+tq(e_i),
\]
and for a nonstar edge $ij$ use
\[
 F_{ij}(t)=\frac{q(e_0+t(e_i+e_j))-q(e_0+t e_i)
       -q(e_0+t e_j)+q(e_0)}{t^2}=s_{ij}.
\]
All the star points needed by the second formula are present because
$G$ contains the full star.  After collecting coefficients, the curve
\[
 P_G(t)=\sum_{0i\in G}\gamma_{0i}\otimes F_i(t)
          +\sum_{ij\in G,\ i,j>0}\gamma_{ij}\otimes F_{ij}(t)
       =P_G+t\sum_{i=1}^{v-1}\gamma_{0i}\otimes q(e_i)
\]
uses the base point $q(e_0)$, $v-1$ star points, and one point for
each of the $j-(v-1)$ nonstar edges: exactly $j+1$ pure summands.
The identities hold over $\mathbb Z[t,t^{-1}]$ and use no inverse of
an integer, so the degeneration works in every characteristic.
Its regular tensor-valued limit proves
$\underline R(P_G)\le j+1$ and finishes the exact row minimum.

We now specialize this construction to the five-edge graph
$\{01,02,03,12,13\}$, whose row tensor is
\[
 G=\gamma_{01}\otimes s_{01}
   +\gamma_{02}\otimes s_{02}
   +\gamma_{03}\otimes s_{03}
   +\gamma_{12}\otimes s_{12}
   +\gamma_{13}\otimes s_{13}
 \in k^5\otimes k^4\otimes k^4.
\]
For $t\ne0$ define
\begin{align*}
 F_i(t)&=\frac{q(e_0+t e_i)-q(e_0)}{t}
       =s_{0i}+tq(e_i),\qquad i=1,2,3,\\
 F_{ij}(t)&=\frac{q(e_0+t(e_i+e_j))-q(e_0+t e_i)
                    -q(e_0+t e_j)+q(e_0)}{t^2}
       =s_{ij},\qquad (i,j)=(1,2),(1,3).
\end{align*}
These are identities over $\mathbb Z[t,t^{-1}]$ and involve no division
by an integer.  Put
\[
 G(t)=\sum_{i=1}^3\gamma_{0i}\otimes F_i(t)
       +\gamma_{12}\otimes F_{12}(t)
       +\gamma_{13}\otimes F_{13}(t).
\]
After collecting the first-factor coefficients, it uses exactly the
six displayed square vectors
\[
 q(e_0),\ q(e_0+t e_1),\ q(e_0+t e_2),\ q(e_0+t e_3),
 q(e_0+t(e_1+e_2)),\ q(e_0+t(e_1+e_3)).
\]
Each has one arbitrary first-factor coefficient and therefore gives
one pure three-factor summand.  Thus $\operatorname R(G(t))\le6$, and
\[
 G(t)=G+t\sum_{i=1}^3\gamma_{0i}\otimes q(e_i)
 \longrightarrow G.
\]
The sum has a regular limit even though the individual first-factor
coefficients may have poles.

The general row result just proved shows that this factor has border
rank exactly six.

The two-edge star tensor
\[
 H=\alpha_1\otimes s_{01}+\alpha_2\otimes s_{02}
       \in k^2\otimes k^3\otimes k^3
\]
has the three-term curve
\[
 H(t)=\alpha_1\otimes F_1(t)+\alpha_2\otimes F_2(t)
      =H+t\sum_{i=1}^2\alpha_i\otimes q(e_i).
\]
It uses $q(e_0),q(e_0+t e_1),q(e_0+t e_2)$, so its border rank is at
most three, and its support flattening of rank three gives equality
(as also proved in P-03).

Take the Kronecker product of these two explicit curves, grouping
corresponding factors.  For every $t\ne0$, $H(t)\boxtimes G(t)$ is a
sum of at most $3\cdot6=18$ pure tensors in
$k^{10}\otimes k^{12}\otimes k^{12}$, and its regular limit is
$H\boxtimes G$.  Identify $e_r\otimes e_c$ in the last two factors
with the matrix-entry variable $x_{rc}$.  Expanding
$s_{0i}\boxtimes s_{cd}$ gives exactly the four ordered coefficients
of the permanent rectangle with row pair $(0,i)$ and column pair
$(c,d)$.  Here $i=1,2$ and
$(c,d)\in\{01,02,03,12,13\}$.  Their rectangle indices are precisely
$S_{10}$, and the ten first-factor labels are independent.
Thus $H\boxtimes G$ is the required residual inclusion tensor.
This proves the upper bound over every algebraically closed field.

In characteristic zero, R-10 and the retained-dimension propagation
R-07 give $\mu_{10}\ge10+7=17$, proving the stated interval.
The exact script \texttt{check\_per10\_upper18.py} verifies both
integral Laurent identities and every coefficient of the Kronecker
identification; its output is
\texttt{per10\_upper18\_verification.json}.
\end{proof}

\paragraph{Dictionary and remaining gate.}
The constructions are morphisms from the punctured affine line over
$\mathbb Z$ to the respective secant parameterizations whose tensor-valued
sum extends over zero.  Its specialization lies in the secant closure
of the Segre product.  The first and mixed second differences are
integral identities in tensor powers of free modules, so base change
preserves the construction in every characteristic.  The only
characteristic-zero ingredient in the interval is the lower bound
already proved in R-10.  In the common scheme and locally-free-sheaf
dictionary this gives a concrete degeneration, not merely an equality
of $K_0$ classes.  The unresolved ten-plane gate is now whether the
uniform minimum is $17$ or $18$; the curve does not provide a
seventeen-term decomposition.

\subsection{Borel flags for determinant extension spaces}
\begin{finding}{P-17}{a finite Borel-flag obstruction for eighteen determinant ten-planes}
Over an algebraically closed field of characteristic zero, each
determinant residual tensor for a subset in the following table, or
its transpose, has border rank at least eighteen. There are eighteen
distinct subsets in this list and its transposes.
\[
\begin{array}{c|l|r|r|r}
 &S&\dim B&\text{quotient weights}&\text{rank exclusions}\\\hline
 1&0,1,2,3,6,7,8,12,13,14&16&101&3\\
 2&0,1,2,3,4,6,7,8,12,13&15&97&18\\
 3&0,1,2,3,6,7,8,9,12,13&15&95&19\\
 4&0,1,2,3,6,7,8,12,13,18&15&96&25\\
 5&0,1,2,3,4,6,7,8,9,12&14&90&84\\
 6&0,1,2,3,4,6,7,8,12,18&14&90&97\\
 7&0,1,2,3,4,6,7,9,12,13&14&91&61\\
 8&0,1,2,3,4,6,7,12,13,18&14&92&74\\
 9&0,1,2,3,6,7,9,12,13,15&14&88&41\\
 10&0,1,2,3,4,6,7,12,18,24&13&88&263
\end{array}
\]
This is a statement about the listed ten-planes. It does not assert
the uniform minimum $\mu_{10}\ge18$.
\end{finding}
\begin{proof}
We first prove the reduction behind the finite calculation.
Let $S$ be a coordinate order ideal in the two row/column two-subset
posets, let $W$ be the corresponding determinant quadratic space, and
let $A=W^*$. Write $B=C$ for its used-variable space, of dimension $m$,
and $T\in A\otimes B\otimes C$ for the residual inclusion tensor.
The product $H$ of the two upper triangular groups on the row and
column index spaces preserves $W$. It also preserves the used-variable
space: lowering a row or column of a used variable either stays in
its rectangle or belongs to the corresponding lower rectangle of $S$.
The contragredient action on $A$ therefore fixes $T$ and preserves
$U=T(C^*)\subset A\otimes B$.

If $\BR(T)\le r$, the proper-incidence necessary condition of R-04
gives an $r$-plane $E\supset U$ with
\[
 \rank\Psi_A(E)\le(\dim A-1)r,\qquad
 \rank\Psi_B(E)\le(m-1)r.
 \tag{P17.1}
\]
This incidence fiber is projective and $H$-invariant. The triangular
Segre--Pl\"ucker fixed-point argument given in P-11 consequently lets
us take $E$ to be $H$-invariant. Put $V=(A\otimes B)/U$ and
$F=E/U$. An $H$-stable complete flag on the triangular ambient
representation descends, after removing repeated terms, to one on
$V$. Intersect it with $F$. Since each ambient step has dimension one,
the intersection dimensions increase by at most one, and deletion of
repeated terms gives an $H$-stable complete flag on $F$.

Thus every possible $F$ is obtainable by successive invariant lines
in quotients $V/F_0$, starting with $F_0=0$. Each such line has one
torus character. The six adjacent row/column root operators must kill
it in the quotient, because a unipotent group acts trivially on a
one-dimensional representation. In characteristic zero these six
operators generate the relevant unipotent action. For the necessity
of the enumeration it already suffices that they kill every line in
an invariant flag.

For each torus weight $w$, compute the common kernel of the six
induced operators on $(V/F_0)_w$. A zero-dimensional kernel offers no
next line. A one-dimensional kernel offers exactly one, with its full
rational coordinate vector retained. A kernel of dimension at least
two would leave a projective family and would require further work;
it is never replaced by sampled lines. In each of the ten cases in
the table, the complete enumeration below encounters no such family
before all branches are excluded.

The exclusion criterion is monotone. If a partial extension
$U+F_0\subset E$ already has an exterior-map rank larger than a bound
in (P17.1), every completion has at least that rank and is impossible.
We may therefore discard that branch immediately. Otherwise we
enumerate all its next invariant lines. This proves by induction that
every possible $E$ remains in the enumeration until an actual rank
contradiction excludes it.

Here is an integral specification of every matrix used. A basis
vector of $A$ is indexed by $Q\in S$ and a variable basis vector by
$v$. Its weight in $A\otimes B$ is
$\operatorname{wt}(v)-\operatorname{wt}(Q)$. In each weight containing
one image vector from $U$, eliminate the last coordinate in
lexicographic $(Q,v)$ order. Its coefficient is $1$ or $-1$. The
remaining coordinate vectors form the fixed quotient basis, and the
projection and inclusion matrices satisfy
\[
 QU=0,\qquad QX=1.
\]
For an adjacent transvection on variables let $N$ be its matrix on
$B$, and let $R$ be its induced matrix on the quadratic space $W$.
Replacing an index of a two-subset gives the ordinary exterior sign;
the order-ideal property ensures that the replacement stays in $S$.
The induced matrix on $A\otimes B$ is
\[
 G=-R^{\mathsf t}\otimes1+1\otimes N.
\]
The checker verifies the integer equivariance identity
$GU=-UN^{\mathsf t}$ and forms the quotient operator $QGX$.
It then forms each successive quotient by exact rational row
reduction, computes the common kernels in every weight, and retains
all distinct resulting subspaces. These are computations over
$\mathbb Q$, valid after every characteristic-zero field extension.

For the listed cases set $r=17$ and $\dim A=10$. The successive
quotient dimensions required are $17-m$. The numbers of distinct
spaces generated at each depth, before rank pruning, are
\[
\begin{array}{c|l|l}
 \text{case}&\text{generated at successive depths}
             &\text{discarded at successive depths}\\\hline
 1&3&3\\
 2&5,18&0,18\\
 3&5,19&0,19\\
 4&6,25&0,25\\
 5&6,25,84&0,0,84\\
 6&6,27,97&0,0,97\\
 7&5,19,57&0,4,57\\
 8&5,21,73&0,1,73\\
 9&4,15,36&0,5,36\\
 10&6,25,85,249&0,0,14,249
\end{array}
\]
Every last layer is empty after pruning. The self-contained
\path{verify_det10_borel.py} reconstructs all tensors, quotient
matrices, root operators, and flag successors directly from the ten
displayed subsets. It imports no project data, search routine, or
symbolic rank oracle. Its exact rational row reduction checks
$22120$ weight kernels; none requiring a projective choice is
encountered. For each of the $685$ rank exclusions it selects and
independently evaluates a nonzero exterior-map minor modulo the
verified prime $1000003$, of order
\[
 9\cdot17+1=154\quad\hbox{or}\quad(m-1)17+1.
\]
It writes the complete flag spaces, minor indices, and nonzero residues
to \path{det10_borel_certificate.json}. The modular products are less
than $2^{40}$ in magnitude and are held in signed 64-bit integers;
all quotient arithmetic uses arbitrary-precision rational numbers.
The residues certify nonzero integer minors, hence the stated
characteristic-zero rank contradictions. There is no admissible
seventeen-plane $E$ in any case. Transposition gives the other eight
subsets, and the finding follows.
\end{proof}

\paragraph{Dictionary and remaining gate.}
The argument takes place on a proper Grassmannian incidence with
determinantal rank conditions. Its fixed locus is refined by relative
flag incidences for the quotient of the tautological bundle. The
split torus gives the character summands, and the adjacent root
subgroups supply the explicit equivariant endomorphisms used to form
their common kernels. The resulting kernels and their projective
spaces are genuine bundle loci; their dimensions are not inferred
from representation classes in $K_0$. These are the same EGA~II,
SGA~3 and FGA constructions located in the common dictionary.
Together with the earlier six Koszul cases, P-17 reduces the
determinant ten-plane test for a bound forty-four to the remaining
$35$ Borel ideals. The additional exploratory files record smaller
extension-space gates for them. Those exploratory survivors do not
constitute a border-rank upper bound or a proof of the uniform lower
bound eighteen.

\part{All-field conventions and analytic findings}
% Requires amsmath, amssymb, amsthm; no private macros.
\section{F-23 over every algebraically closed field}
\label{sec:fields}

Throughout this section $k$ is algebraically closed, with arbitrary
characteristic. Border rank means CP border rank in the displayed
tensor factors, defined by Zariski closure of sums of pure tensors.
The integral convention below is essential: the normalized
characteristic-zero coefficient tensor cannot simply be reduced
modulo a prime dividing its normalization denominator.

\paragraph{Finding FL-01 (the integral models).}
Let $A_n=\mathbb Z^n$ have basis $e_1,\ldots,e_n$, and let
$V_n=\mathbb Z^{n^2}$ have basis $x_{ij}$.
Define
\begin{align*}
 D_n&=\sum_{\sigma\in\mathfrak S_n}\operatorname{sgn}(\sigma)
             e_{\sigma(1)}\otimes\cdots\otimes e_{\sigma(n)},&
 P_n&=\sum_{\sigma\in\mathfrak S_n}
             e_{\sigma(1)}\otimes\cdots\otimes e_{\sigma(n)},\\
 S_{\det,n}&=\sum_{\sigma,\tau\in\mathfrak S_n}
       \operatorname{sgn}(\sigma)\operatorname{sgn}(\tau)
       x_{\sigma(1),\tau(1)}\otimes\cdots\otimes x_{\sigma(n),\tau(n)},\\
 S_{\mathrm{per},n}&=\sum_{\sigma,\tau\in\mathfrak S_n}
       x_{\sigma(1),\tau(1)}\otimes\cdots\otimes x_{\sigma(n),\tau(n)}.
\end{align*}
The latter two are invariant under permuting the tensor slots.
For free modules they therefore belong to the integral divided-power
model $\Gamma^n(V_n)\simeq(V_n^{\otimes n})^{\mathfrak S_n}$.
In characteristic zero they are $n!$ times the normalized polynomial
coefficient tensors of predecessor GCT\_GINEQ F-23. Their images in
the ordinary symmetric quotient are $n!\det_n$ and $n!\mathrm{per}_n$.

\begin{proof}
An ordered sequence of matrix positions with every row and every
column used once uniquely determines its row permutation $\sigma$
and column permutation $\tau$. Thus each tensor-basis word occurs
once, with coefficient $\pm1$ or $1$. Permuting all slots composes
both permutations by the same permutation, so their sign product
is unchanged. A basis of the invariant tensors consists of the sums
over distinct reorderings of each monomial; this is the usual
divided-power basis. In the ordinary symmetric quotient, each
squarefree matching monomial has $n!$ such reorderings. The claimed
formulas follow. In characteristic two signs coalesce, so the
determinant and permanent models agree. No integral coefficient
vanishes by normalization: each surviving ordered word has a unit
coefficient. If $p\mid n!$, the symmetric quotient of either
invariant tensor is zero while the invariant tensor is nonzero.
\end{proof}

These definitions fix what is meant by the polynomial families in
an all-field statement. They do not identify the ordinary symmetric
quotient with invariant tensors by division by $n!$.

\paragraph{Finding FL-02 (proper-incidence border substitution).}
Let $T\in A\otimes V_1\otimes\cdots\otimes V_b$, let
$W=V_1\otimes\cdots\otimes V_b$, and let $U\subset W$ be the
image of the $A$-flattening. If $q=\dim U\geq1$, then
\begin{equation}\label{eq:fields-substitution}
 \underline{\operatorname R}_k(T)
 \geq q-1+\min_{0\ne u\in U}\underline{\operatorname R}_k(u).
\end{equation}
This is the needed all-field, arbitrary-residual-order form of
border substitution. Landsberg--Micha\l ek
\cite[Proposition~2.3, p.~3]{LM2018} already state the three-factor
principle for every retained dimension, not just dimension one.
The all-field, arbitrary-residual-order form follows from R-01.

\begin{proof}
Apply R-01 with retained dimension $j=1$. A line
$L=\langle u\rangle$ has inclusion tensor $a\otimes u$,
whose border rank equals that of $u$: adjoining a fixed nonzero
factor and contracting against its dual are inverse linear maps
on this tensor. Consequently the one-dimensional Grassmannian
minimum in R-01 is precisely the minimum in
\eqref{eq:fields-substitution}.
\end{proof}

\paragraph{Finding FL-03 (exact residual minima and recursion).}
Let $T_n$ denote any one of the four integral families in FL-01.
For $2\leq b<n$, group the first $n-b$ slots and denote the
remaining flattening image by $U_{n,b}$. Put $s=1$ for the row
families and $s=2$ for the polynomial divided-power families.
Then
\begin{align}
 \min_{0\ne u\in U_{n,b}}\underline{\operatorname R}_k(u)
     &=\underline{\operatorname R}_k(T_b),\label{eq:fields-minimum}\\
 \underline{\operatorname R}_k(T_n)
     &\geq\binom nb^s+\underline{\operatorname R}_k(T_b)-1.
       \label{eq:fields-recursion}
\end{align}

\begin{proof}
In a row family, contract the first $n-b$ slots against an ordered
tuple of coordinate covectors. If an index repeats, the result is
zero. Otherwise the result is, up to a sign, the embedded copy
$Q_I$ of $T_b$ on the complementary $b$-subset $I$. Every $I$
occurs. Distinct $Q_I$ have disjoint tensor-basis supports, each
with a unit coefficient. They therefore form a free integral basis
of the image, remaining independent after every scalar extension.
Its dimension is $\binom nb$.

For the polynomial families, a nonzero contraction is exactly a
partial row-and-column matching. Its residual is, up to a sign,
the coefficient tensor $Q_{I,J}$ of the complementary minor or
permanent. Every pair of $b$-subsets $I,J$ occurs. The residual
supports are disjoint and have unit coefficients, giving image
dimension $\binom nb^2$ over every field. This is the tensor version
of the complementary-minor calculation in GCT\_GINEQ F-22.
In particular, no factor $b!$ or $(n-b)!$ occurs here.

Choose a nonzero coefficient $c$ of $u$ in this basis. In a row
family, project every residual factor onto the coordinates $I$
of that coefficient. Each other $b$-subset has a required index
outside $I$ and is killed, whereas $Q_I$ maps to $T_b$.
In a polynomial family, project onto the matrix-entry block
$I\times J$ in every factor. Each other pair of $b$-subsets has
a required row or column outside this block, and is killed.
The image is $cT_b$.

Products of linear maps preserve pure tensors or send them to
zero; continuity gives the same border-rank monotonicity.
Consequently every nonzero $u$ has border rank at least that of
$T_b$. Conversely, a single $Q_I$ or $Q_{I,J}$ is an embedding of
$T_b$ with the displayed projection as a left inverse, so its
border rank is exactly that of $T_b$. This proves
\eqref{eq:fields-minimum}. Grouping the first $n-b$ factors
cannot increase CP border rank. Apply FL-02 to that grouped
tensor and the image dimension just computed to obtain
\eqref{eq:fields-recursion} for the original tensor.
\end{proof}

\paragraph{Finding FL-04 (unit-minor seeds).}
Over every algebraically closed field,
\[
 \underline{\operatorname R}_k(D_2)=
 \underline{\operatorname R}_k(P_2)=2,\qquad
 \underline{\operatorname R}_k(S_{\det,2})=
 \underline{\operatorname R}_k(S_{\mathrm{per},2})=4.
\]
Hence, for $n\geq2$,
\[
 \underline{\operatorname R}_k(D_n),
 \underline{\operatorname R}_k(P_n)\geq\binom n2+1,
 \qquad
 \underline{\operatorname R}_k(S_{\det,n}),
 \underline{\operatorname R}_k(S_{\mathrm{per},n})
       \geq\binom n2^2+3.
\]

\begin{proof}
The two row matrices are
$\left(\begin{smallmatrix}0&1\\-1&0\end{smallmatrix}\right)$
and $\left(\begin{smallmatrix}0&1\\1&0\end{smallmatrix}\right)$.
The polynomial matrices have two off-diagonal coordinate pairs,
$(x_{11},x_{22})$ and $(x_{12},x_{21})$, with unit entries.
They are signed permutation matrices of size four. Their ranks
are respectively two and four over every field. Matrix rank
equals matrix border rank. For $n>2$ use FL-03 with $b=2$;
for $n=2$ use the seeds themselves.
\end{proof}

Thus the predecessor's bound $39$ at $n=4$ has an all-field
proof in the explicit integral model. Stronger multidimensional
certificates elsewhere in the paper improve this seed recursion
over every algebraically closed field in the integral convention.

\paragraph{Finding FL-05 (ordinary polynomial powers in small characteristic).}
If $\operatorname{char}k=p\leq n$, neither ordinary polynomial
$\det_n$ nor $\mathrm{per}_n$ lies in any secant of the image
$\ell\mapsto\ell^n$ in $\operatorname{Sym}^n(V_n)$.
Their border rank relative to this ordinary-power image is
therefore infinite. This statement concerns the ordinary symmetric
quotient, not the complete-linear-system Veronese in divided powers.

\begin{proof}
The coefficient of a squarefree degree-$n$ monomial in $\ell^n$
is $n!$ times a product of coefficients of $\ell$. It is zero
when $p\leq n$. All sums of ordinary powers, and all their limits,
therefore lie in the closed linear subspace on which every such
coefficient vanishes. Determinant and permanent each have nonzero
squarefree matching coefficients and lie outside that subspace.
\end{proof}

When $n!$ is invertible, the usual symmetrization identification
is valid. In that range, every CP border-rank lower bound here
also bounds ordinary polynomial symmetric border rank. In small
characteristic the finite meaningful assertion uses FL-01's
divided-power model and its complete Veronese; FL-05 explains why
the alternative ordinary-power convention has a different answer.

% BR_PENCIL SDP/orbit sections. Uses standard amsmath/amssymb and theorem/proof environments.
% Root may rename local finding labels S-01,S-02,O-01,O-02 to package identifiers.
\section{An exact DPS certificate outside the four families}
\label{sec:dps-example}

All spaces in the analytic sections are finite-dimensional complex
Hermitian spaces. With $P_U$ the orthogonal projector, we use
\[
 \alpha(U)=\max_{\|z_i\|=1}
 \left\langle\bigotimes_i z_i,
 P_U\bigotimes_i z_i\right\rangle
 =\max_{u\in U,\,\|u\|=1}\varepsilon(u)^2,
\]
where $\varepsilon(u)=\max_{\|z_i\|=1}
|\langle\bigotimes_i z_i,u\rangle|$ is the injective norm.
The equality follows by maximizing the inner product over the unit
sphere of $U$, which gives the norm of the orthogonal projection,
and interchanging two maxima over compact sets. This is the
projector formulation used in predecessor GCT\_GINEQ F-13.
The DPS hierarchy is used with the support convention stated below;
see \cite[Section II and Theorem 1, printed p.~7]{DPS}.

\paragraph{Finding S-01 (exact finite-level certification).}
In $A\otimes B=\mathbb C^2\otimes\mathbb C^3$ let
\[
 u_0=(|00\rangle+|11\rangle)/\sqrt2,\qquad
 u_1=(|01\rangle+|12\rangle)/\sqrt2,\qquad U=\langle u_0,u_1\rangle.
\]
Then $\alpha(U)=3/4$. The first PPT level of the
Doherty--Parrilo--Spedalieri hierarchy certifies this value exactly,
and every higher PPT Bose-symmetric-extension level has zero error.

\begin{proof}
Write $P=|u_0\rangle\langle u_0|+|u_1\rangle\langle u_1|$,
and let $\Gamma$ denote partial transpose on $B$. Define
\begin{align*}
 X={1\over16}\bigl(&|2\cdot00-11\rangle\langle2\cdot00-11|
       +|01-2\cdot12\rangle\langle01-2\cdot12|\bigr),\\
 Y={3\over16}\bigl(&|01-2\cdot10\rangle\langle01-2\cdot10|
       +|2\cdot02-11\rangle\langle2\cdot02-11|\bigr).
\end{align*}
The notation in each ket denotes the indicated linear combination of
basis vectors. Both matrices are positive semidefinite by their Gram
factorizations, and direct expansion gives
\begin{equation}\label{eq:dps-exact-witness}
 {3\over4}I_6-P=X+Y^\Gamma.
\end{equation}
For every trace-one $\rho\succeq0$ with $\rho^\Gamma\succeq0$, this yields
\[
 {3\over4}-\operatorname{tr}(P\rho)
 =\operatorname{tr}(X\rho)+\operatorname{tr}(Y\rho^\Gamma)\geq0.
\]
The product vector
\[
 z={a_0+a_1\over\sqrt2}\otimes{b_0+2b_1+b_2\over\sqrt6}
\]
has $\langle z,Pz\rangle=3/4$; its density matrix is feasible. Thus
this rational primal/dual pair has zero gap. Since product states
are PPT, it also proves $\alpha(U)=3/4$.

For level $k\geq1$, optimize the same objective over marginals of
trace-one positive operators on $A\otimes\operatorname{Sym}^k B$,
with positive partial transposes for every cut. Such marginals are
PPT, so the same dual certificate bounds their objective by $3/4$.
The product extension of $|z\rangle\langle z|$ is feasible for
every $k$ and attains $3/4$. Hence every level has exactly this
value and error zero.
\end{proof}

For an independent direct check, a unit vector $cu_0+du_1$,
$|c|^2+|d|^2=1$, has matrix
\[
 {1\over\sqrt2}\begin{pmatrix}c&d&0\\0&c&d\end{pmatrix}.
\]
Its largest squared Schmidt coefficient is $(1+|cd|)/2$, whose
maximum is $3/4$. The script \path{check_certificates.py} checks
\eqref{eq:dps-exact-witness}, the Gram eigenvalues, the primal
objective and complementary slackness using exact rational arithmetic.
No numerical optimization or tolerance is used.

The subspace is a residual image of
\[
 T=c_0\otimes(a_0b_0+a_1b_1)+c_1\otimes(a_0b_1+a_1b_2)
       \in\mathbb C^2\otimes\mathbb C^2\otimes\mathbb C^3.
\]
This concise format is outside the four determinant/permanent
families. The integral identity
\begin{equation}\label{eq:karatsuba}
 T=c_0a_0(b_0-b_1)+c_1a_1(b_2-b_1)
       +(c_0+c_1)(a_0+a_1)b_1
\end{equation}
gives rank at most three. The $B$-flattening has independent
images $c_0a_0$, $c_0a_1+c_1a_0$, $c_1a_1$, and therefore rank
three. Thus $\operatorname R(T)=\underline{\operatorname R}(T)=3$
over every field. Over $\mathbb C$, the DPS certificate and
GCT\_GINEQ F-13 give the same lower bound
$2+\lceil4/3\rceil-1=3$. This finishes the requested example;
it does not provide a generic efficient solution of predecessor Q-30.

\paragraph{Finding S-02 (a quantitative projector-error bound).}
Let $V\subseteq A\otimes B$ and $d=\dim B$. Write $\beta_k(V)$
for the PPT Bose-symmetric-extension upper bound just defined. Then
\begin{equation}\label{eq:dps-error}
 0\leq\beta_k(V)-\alpha(V)\leq {d-1\over k+d-1}.
\end{equation}
For S-01 the exact bound is zero; the generic dimension-three
bound is $2/(k+2)$.

\begin{proof}
The known input is Navascu\'es--Owari--Plenio,
\emph{The power of symmetric extensions for entanglement detection},
arXiv:0906.2731v1, Theorem 2, equation (7), printed p.~3:
if $\rho$ admits a $k$-Bose-symmetric extension, then
\[
 \sigma={k\rho+\rho_A\otimes I_B\over k+d}
\]
is separable. The support convention is their Definition 1,
printed p.~2 \cite{NOP}. It applies in particular to every feasible
PPT-extension marginal. Both $\rho$ and $\sigma$ are density
operators with marginal $\rho_A$. Separability implies
$R:=\rho_A\otimes I_B-\sigma\succeq0$: in a product-state
convex decomposition each summand is of the form
$p\tau_A\otimes(I_B-\tau_B)$. Also $\operatorname{tr}R=d-1$.
Rearranging gives
\[
 \rho-\sigma={(d-1)\rho-R\over k+d-1}.
\]
Since $0\preceq P_V\preceq I$,
\[
 \operatorname{tr}(P_V\rho)-\operatorname{tr}(P_V\sigma)
 \leq {d-1\over k+d-1}.
\]
The second objective is at most $\alpha(V)$ by separability.
Take the maximum over feasible $\rho$. Conversely, product
states have feasible product extensions, giving the other
inequality. The maxima exist because the extension feasible
set is closed inside the compact trace-one positive cone.
\end{proof}

\section{Orbit optimization: a nonattained constant supremum}
\label{sec:orbit-alpha}

\paragraph{Finding O-01 (the proposed maximum is trivial or absent).}
For every nonzero $U\subseteq\bigotimes_{i=1}^b H_i$,
\begin{equation}\label{eq:orbit-alpha-sup}
 \sup_{g\in\prod_i\operatorname{GL}(H_i)}\alpha(gU)=1.
\end{equation}
The supremum is attained if and only if $U$ contains a nonzero
product tensor; in that case $\alpha(gU)=1$ throughout the orbit.
The assertions also hold for the product of special linear groups.

\begin{proof}
Choose $0\ne u\in U$ and a nonzero coefficient $c$ in a product
orthonormal basis. Relabel that coefficient as $(0,\ldots,0)$.
Let $g_i(t)$ fix $e_0$ and multiply every other basis vector by
$t$, where $0<t<1$. Then
$g(t)u\longrightarrow c z$ with $z=e_0\otimes\cdots\otimes e_0$.
The normalized vectors $v_t=g(t)u/\|g(t)u\|$ belong to $g(t)U$,
and $|\langle v_t,z\rangle|^2\longrightarrow1$.
Since $\alpha(g(t)U)\geq|\langle v_t,z\rangle|^2$ and
$\alpha\leq1$, \eqref{eq:orbit-alpha-sup} follows.

Compactness of the product of unit spheres implies that
$\alpha(W)=1$ precisely when some unit product vector has
projection of norm one onto $W$. The orthogonal decomposition
then forces that vector to belong to $W$. Invertible product
maps preserve this property, which proves the attainment
criterion. Over $\mathbb C$, determinant roots show that
product-GL and product-SL maps induce the same subspace orbits:
they differ only by an overall scalar, which does not change
a subspace. This proves the last assertion.
\end{proof}

If $0\ne U\subseteq\operatorname{Sym}^b H$ in characteristic
zero, the same result holds for the common-factor action
$g^{\otimes b}$ of $\operatorname{GL}(H)$ or $\operatorname{SL}(H)$.
Indeed, for any nonzero $u\in U$, the polynomial
$\ell\mapsto\ell^{\otimes b}(u)$ is nonzero. Choose a unit vector
whose dual evaluates it nontrivially and apply the common
one-parameter degeneration fixing that vector and shrinking
its orthogonal complement. The limit is its nonzero pure
$b$th power. The proof above applies. Attainment is equivalent
to $U$ containing a pure power. This covers the polynomial
residual images and the row permanent under their diagonal
actions. The group must be specified: the row determinant
image $U=\Lambda^b H$ is fixed under diagonal GL, and its
alpha remains $1/b!$ by predecessor F-14.

Thus the intake's proposed \emph{maximum} over the GL orbit
fails to exist for product-free residual spaces, and its
supremum is the constant one. Substituting that value into
F-13 yields only the image-dimension bound. Moreover,
maximizing alpha weakens that certificate; smaller alpha is
the useful direction. This refutes the particular proposed
quantity, without ruling out other orbit optimizations.

\paragraph{Finding O-02 (the actual Pl\"ucker norm condition).}
Put $W=\bigotimes_iH_i$, $q=\dim U$, choose an orthonormal basis
$u_1,\ldots,u_q$, and let
$w=u_1\wedge\cdots\wedge u_q\in\Lambda^q W$. For the product-SL
action, $w$ minimizes its orbit norm if and only if
\begin{equation}\label{eq:plucker-moment}
 \operatorname{Tr}_{\widehat i}P_U={q\over\dim H_i}I_{H_i}
 \qquad\hbox{for every }i.
\end{equation}
Here $\operatorname{Tr}_{\widehat i}$ denotes partial trace
in all factors except $i$.

\begin{proof}
For traceless Hermitian endomorphisms $H_i$ let
$H_W=\sum_i I\otimes\cdots\otimes H_i\otimes\cdots\otimes I$.
Differentiating the exterior-power action and pairing with the
unit wedge gives
\[
 \left.{d\over dt}\right|_{t=0}\log\|\exp(tH_W)w\|^2
 =2\sum_{j=1}^q\langle u_j,H_Wu_j\rangle
 =2\sum_i\operatorname{tr}\bigl(H_i\operatorname{Tr}_{\widehat i}P_U\bigr).
\]
A Hermitian matrix pairs to zero with every traceless Hermitian
matrix precisely when it is scalar. Its trace fixes that
scalar at $q/\dim H_i$. Thus \eqref{eq:plucker-moment} is
exactly the critical-point condition.

Every element of the product-SL group has polar decomposition
$g=k\exp H$ with $k$ unitary and $H$ in the indicated Hermitian
Lie algebra. Diagonalize the induced Hermitian generator on
$\Lambda^q W$. Then
\[
 f(t)=\log\|\exp(tH)w\|^2
     =\log\sum_j |w_j|^2e^{2t\lambda_j}
\]
is convex: its second derivative is four times the variance
of the real $\lambda_j$ under the normalized positive weights.
At a critical point $f'(0)=0$, so $f(1)\geq f(0)$.
The unitary factor does not change the norm, proving global
minimality. Necessity follows by differentiation at a minimum.
\end{proof}

This is the minimal-vector part of Kempf--Ness in explicit
subspace coordinates. Capacity is a \emph{norm minimization}
with scalar homotheties removed. For the corresponding tensor
statement see B\"urgisser--Garg--Oliveira--Walter--Wigderson,
arXiv:1711.08039v1, equation (1.3) and Theorem 3.2, printed p.~13;
the latter uses the \emph{orbit closure} \cite{BGOWW}. No identity between
capacity and the supremum of alpha is asserted.

For example, the line through
$\Omega_d=d^{-1/2}\sum_i e_i\otimes e_i$ satisfies
\eqref{eq:plucker-moment}, while $\alpha(\langle\Omega_d\rangle)=1/d$.
For $d\geq2$, its orbit-alpha supremum is nevertheless one and
is unattained. For S-01, diagonal scaling in $B$ gives
\[
 U_{\mathrm{bal}}=\left\langle
 {\sqrt2|00\rangle+|11\rangle\over\sqrt3},
 {|01\rangle+\sqrt2|12\rangle\over\sqrt3}\right\rangle.
\]
Its projector marginals are $I_2$ and $(2/3)I_3$, so it satisfies
\eqref{eq:plucker-moment}. Every unit vector in it has Schmidt
squares $2/3$ and $1/3$: its row Gram matrix has trace one and
eigenvalue difference
$\sqrt{(|c|^2-|d|^2)^2+4|cd|^2}/3=1/3$.
Hence its alpha is $2/3$, whereas S-01 has alpha $3/4$ and
the common orbit has supremum one. We do not claim here that
$2/3$ is the orbit-infimum of alpha.


\part{Literature, scope, and remaining gates}
\section{Literature comparison, checked on 21 November 2026}
\label{sec:literature}

The comparison distinguishes the row tensors $D_n,P_n$ in
$(\mathbb C^n)^{\otimes n}$ from the polynomial coefficient tensors
$S_{\det,n},S_{\mathrm{per},n}$ in
$(\mathbb C^{n^2})^{\otimes n}$. It also distinguishes CP border
rank, symmetric border rank, Waring rank, and cactus rank.
For a symmetric tensor, CP border rank is at most symmetric
border rank, which is at most Waring rank. A lower bound on
Waring rank alone does not lower-bound border rank; neither does
a lower bound on cactus rank alone without further hypotheses.

\subsection{The numerical comparisons actually verified}

Landsberg--Teitler \cite[Corollary~1.4, p.~2; \S9, pp.~13--14,
Table~3]{LT2010} give Waring-rank lower bounds $43$ for $\det_4$
and $40$ for $\mathrm{per}_4$. Their table gives border lower
bound $36$ for both. The distinction follows from their rank
definitions and the surrounding argument, even where PDF text
extraction loses the underline. The middle catalecticant has
rank $\binom{n}{\lfloor n/2\rfloor}^2$.

Farnsworth \cite[Theorem~1.6, p.~2]{Farnsworth2016} proves
symmetric border rank at least $38$ for $\det_4$.
Theorem~1.3 gives his general-$n$ determinant estimate for $n\geq5$;
Theorem~1.8 gives symmetric border lower bound $14$ for both
degree-three polynomials. His bound $38$ is not silently relabeled
as a CP lower bound. The CP bounds proved here imply symmetric
bounds and therefore improve the specific symmetric comparison
when their value exceeds $38$.

Conner--Harper--Landsberg
\cite[Theorem~1.2, \S1.2]{CHL2023} prove
\[
 \underline{\operatorname R}_{\mathrm{CP}}(S_{\det,3})=17.
\]
Their tensor is explicitly in $(\mathbb C^9)^{\otimes3}$.
Conner--Huang--Landsberg
\cite[Theorem~1.1 and the following identification, p.~3;
proof in \S5]{CHHL2023} prove
\[
 \underline{\operatorname R}_{\mathrm{CP}}(S_{\mathrm{per},3})=16.
\]
These are polynomial tensors, whereas the later small numbers
for size four in the following paragraph concern row tensors.

Han--Ju--Kim \cite[Theorem~1.1 and \S3]{HJK2026} prove
$\underline{\operatorname R}(D_4)=12$ over $\mathbb C$.
The inspected version is arXiv:2510.11051v2, revised on
18 August 2026; its version history records a repaired gap.
Their other paper \cite[Theorems~5.6--5.7]{HJK2025} gives
$\underline{\operatorname R}(P_4)=\operatorname R(P_4)=8$
in characteristic different from two. Neither $12$ nor $8$
is a bound for a polynomial tensor in sixteen variables.

Two further 2026 papers make this ambient distinction
especially important. Han--Song
\cite[Theorem~A, p.~3; Proposition~4.1, p.~10]{HS2026}
prove ordinary CP rank $\operatorname R(P_5)=16$ in
characteristic zero. Their Proposition~4.1 also states
$\underline{\operatorname R}(P_n)\geq
\binom{n}{\lfloor n/2\rfloor}+1$ over $\mathbb C$
and its subfields for $n\geq3$, using a limiting-spans
argument. Thus this row-family ``plus one'' phenomenon
has an explicit current source. Han
\cite[Theorem~A, p.~2]{Han2026} proves symmetric border
rank $2^{n-1}$ for the squarefree monomial
$x_1\cdots x_n$, equivalently for the symmetric row
tensor $P_n$, in characteristic zero. This theorem
concerns symmetric border rank; it does not identify
CP border rank for every $n$, and it does not concern
the permanent polynomial in $n^2$ variables.

The Waring and cactus improvements reviewed by
Boij--Teitler \cite[background, pp.~5--7; \S4, p.~15]{BT2020},
including the Waring lower bound $50$ for $\det_4$, do not
replace the border-rank comparison. Houston--Goucher--Johnston
\cite{HGJ2024} were also checked for their determinant formulas
and tensor/Waring-rank bounds; no competing polynomial border
lower bound at the target size was located in that text.

\subsection{Which part is standard substitution, and which part is evaluated here}

Landsberg--Micha\l ek \cite[Proposition~2.3, p.~3]{LM2018}
already allow every retained dimension $a'\leq a$ in
\[
 \underline{\operatorname R}(T)
 \geq a-a'+
 \min_{A'\in\operatorname{Gr}(a',A^*)}
       \underline{\operatorname R}(T|_{A'}).
\]
Thus passing to $j>1$ is not claimed to invent a substitution
principle. The contribution to verify is the evaluated
Grassmannian minimum for the present residual image, with
the lower certificates and the explicit degenerations giving
upper witnesses. The short proper-incidence proof is included
for the all-characteristics and higher-order formulation
required here.

Torus and Borel symmetry reductions are likewise established
tools. Findings P-01 and P-07 give direct coordinate proofs
with their characteristic hypotheses explicit; the exact
residual minima and certificates are the deductions of this run.

\subsection{The border-apolarity comparison}

Buczy\'nska--Buczy\'nski
\cite[Theorem~3.15, p.~15; Theorem~4.3, p.~18;
Corollary~4.5, p.~19]{BB2021} formulate border apolarity through
the multigraded Hilbert-scheme component
$\operatorname{Slip}_r(X)$ of limits of ideals of points.
The fixed-ideal theorem allows a connected solvable stabilizer
to fix a witnessing ideal. Conner--Harper--Landsberg
\cite[Theorems~2.3--2.4, \S2.3]{CHL2023} use fixed limiting
planes and fixed ideals to make this framework effective.

The current simultaneous version is
Buczy\'nska--Buczy\'nski
\cite[\S4.3, Theorem~4.7; field convention at the start of
\S1]{BB2026}. Over every algebraically closed field, for a
smooth projective toric $X$, a very ample line bundle $L$,
and a subspace $E$ of the corresponding section-dual space
with $\dim E=i$, it gives
\[
 [E]\in\sigma_{r,i}(X)
 \quad\Longleftrightarrow\quad
 \exists[I]\in\operatorname{Slip}_r(X),\quad I\subseteq\operatorname{Ann}(E).
\]
If a connected solvable group preserves $E$, one may choose
$I$ invariant under that group. This provides the appropriate
current framework for simultaneous residual decompositions.
It does not itself supply the numerical minima evaluated here.

Two hypotheses matter. Correct Hilbert function and apolar
containment alone do not replace membership in
$\operatorname{Slip}_r(X)$. Also, the torus in P-01 preserves
the ambient residual space $U$ but generally moves its chosen
subspace $W$. Consequently P-01 is proved by an explicit
degeneration and closedness, not by applying a fixed-ideal
theorem with a group that does not stabilize $W$.

The later Ma\'ndziuk paper
\cite[Theorems~1.1--1.3, p.~2; Theorems~3.5, 4.4, 5.2;
\S6]{Mandziuk2026} adds tests beyond weak border apolarity.
Its examples are a Sch\"onhage sum, a Leitner tensor product,
two skew-Coppersmith--Winograd copies, and a reconsideration
of $3\times3$ matrix multiplication. Dole\v z\'alek--Micha\l ek
\cite[Theorem~1.1; \S5.1; \S6.3]{DM2026} give nonlinear
Kronecker--Koszul equations; their \S6.3 leaves extensive
special-tensor tests, including determinant tensors, for
future work. These current papers were inspected; no
polynomial $4\times4$ border lower bound competing with
the present target was located in their stated results.

\paragraph{Finding FL-06 (a displayed survey bound is false as written).}
Dutta--Lysikov \cite[\S4.3, p.~29, immediately after Corollary~56]{DL2025}
display the universal lower bound $\binom{2n}{n}$ for border
Waring rank of both $\det_n$ and $\mathrm{per}_n$. Taken
literally, this is false at $n=2$.

\begin{proof}
For either sign, $ad\pm bc$ has a nonsingular symmetric
coefficient matrix of size four. The first catalecticant
therefore gives border Waring rank at least four. Conversely,
\[
 ad\pm bc=
 \frac{(a+d)^2-(a-d)^2\pm(b+c)^2\mp(b-c)^2}{4}.
\]
Over $\mathbb C$ the scalar coefficients can be absorbed
into the squared linear forms. Thus its Waring rank and
border Waring rank are four, whereas $\binom42=6$.
\end{proof}

The correct standard comparison verified above is the
middle-catalecticant bound
$\binom{n}{\lfloor n/2\rfloor}^2$. We do not infer which
correction the survey authors intended.

\paragraph{Scope of the search.}
Primary source statements, their ambient definitions, and
relevant citation chains were checked on 21 November 2026.
Queries included determinant/permanent polynomial tensors,
CP and symmetric border rank, the numbers $39$, $40$ and $42$,
Farnsworth, and border apolarity, including 2025--2026 work.
No published polynomial border lower bound of at least $39$
for either size-four polynomial was located in the sources
checked. This is a scoped literature-audit outcome, not a
proof of exhaustive priority. The numerical improvements
are asserted against the specific verified comparisons,
without promoting absence from a search into a historical
novelty theorem.

\part{Unchecked proposals and open numerical gates}
\section{Precisely what remains undecided}

All findings in the earlier parts have complete proofs and recorded
whole-proof checks. The following are questions, not asserted findings.
Their presence does not weaken the checked bounds.

The seven-dimensional question is resolved in characteristic zero:
R-08 and R-10 prove $\mu_7=14$ for both polynomial families, and
R-09/R-10 give $\mu_8=15$, $\mu_9=16$. Each gives the full-tensor
lower bound forty-three. Over every algebraically closed field,
the independently established bound forty-two remains in force.
The characteristic-zero prolongation and modular certificates
are not automatically specialized to all positive characteristics.

The next size-four gate is the tenth minimum. Propagation gives
$\mu_{10}\geq17$. P-20 gives the sharper permanent upper bound
eighteen. For the determinant, an anchored nine-plane has border
rank sixteen; adjoining one further rectangle, of rank four, gives
$\mu_{10}\leq20$ by subadditivity. A uniform lower bound eighteen
would give $36-10+18=44$. The determinant reduces to $59$ Borel
ideals. Direct Koszul tests establish eighteen for six of these;
P-17 excludes eighteen further ideals by exact Borel flags.
The remaining thirty-five ideals are the current checked gate.
Further extension calculations are recorded in the research files
and do not yet constitute a global forty-four theorem. No untested parameter
branch is treated as excluded.

For arbitrary residual order $b>2$, P-02 gives an exact adjacent-pair
reduction of the two-dimensional minimum. P-11 evaluates the row
cubic seed: six for the permanent in every characteristic, eight
for the determinant outside characteristic two, and six in
characteristic two. P-18 bounds the polynomial cubic seed by
$21\leq\BR(M_+)\leq24$ in every characteristic and
$21\leq\BR(M_-)\leq27$ outside characteristic three, with the
common $21$--$24$ interval in characteristic two. The determinant
lower bound twenty-one in characteristic three is not established
by the supplied minors. The polynomial cubic exact values and
the higher-$b$ numerical formulas remain open gates.

R-11 settles the seventh minimum uniformly in matrix size:
$\mu_7=14$ for both polynomial families and every $n\geq4$
in characteristic zero, with eighth and ninth minima fifteen
and sixteen. Its full extra support-graph cover is certified.
P-20 also settles every row permanent quadratic minimum:
$\mu_j=j+1$ for all allowed $n,j$, in every characteristic.

The cubic two-plane route is decisively capped by P-18:
its size-four output is at most forty-one for the determinant
and thirty-eight for the permanent. Its unresolved exact values
therefore cannot improve forty-three by this substitution.

The requested SDP example is finished. A generic efficient
algorithm for predecessor Q-30, and an orbit-infimum version of
alpha, are distinct questions not settled by that example or by
the refutation of the proposed orbit maximum. No such algorithm
or infimum formula is asserted here.

The literature audit establishes the comparison within its listed
primary-source scope, not exhaustive historical priority. Exact
published Grothendieck framework locators are supplied in the
next part; no unpublished manuscript cote for these modern rank
findings has been established.

% This fragment is intended as its own part. It requires only standard LaTeX
% plus amsmath/amssymb and hyperref/url for the bibliography.
\part{The findings in algebraic geometry and Grothendieck's frameworks}
\label{part:dictionary}

\section{Exact published and archival locators}
\label{sec:groth-locators}

The following locators identify the foundations carrying the
proofs across languages. They do not assert that Grothendieck's
texts state the modern tensor-complexity conclusions. PDF leaves
below are one-based counts in the specific facsimiles linked in
the bibliography; printed pagination is listed separately.

\begin{description}
\item[EGA properness.]
EGA~II \cite{EGA2}, Definition~5.4.1 and
Proposition~5.4.2(i)--(iv), printed p.~100, PDF leaf~97;
Theorem~5.5.3(i), printed p.~104, leaf~101.
These give universal closedness, the stability of proper
morphisms, and projective implies proper. Both leaves were
freshly rendered and visually checked in this run.

\item[FGA parameter spaces.]
Grothendieck's Bourbaki expos\'e~221 \cite{FGA},
Theorems~3.1--3.2, printed pp.~259--260, PDF leaves~12--13;
the Hilbert specialization after Proposition~3.8 and
Remarks~3.9, pp.~265--266, leaves~18--19.
These give the Quot/Hilbert parameter spaces and fixed-polynomial
projectivity used for limiting configurations. These four
facsimile leaves were freshly viewed.

\item[Group and module schemes.]
SGA~3 \cite{SGA3}, Expos\'e~I, \S\S4.5--4.6,
including Definition~4.6.1, printed p.~29, PDF leaf~42,
gives the general linear group and module functors. This leaf
was freshly viewed. In particular,
$\operatorname{GL}_n=\operatorname{Spec}
\mathbb Z[T_{ij},\det(T)^{-1}]$.

\item[Exterior powers, representation rings, and the gamma filtration.]
SGA~6 \cite{SGA6}, Expos\'e~V, Definition~2.1,
Examples~2.2(b), and Lemma~2.2.1, printed pp.~314--316,
PDF leaves~321--323; the representation-ring example before
\S2.3, p.~317, leaf~324; Proposition~3.2(ii), p.~325,
leaf~332; \S3.10, equation~(3.10.1), p.~331, leaf~338.
These supply the exterior-power operations on $K_0$, the
representation-ring example, and the gamma filtration.
All six listed leaves were freshly viewed.

\item[Scheme-theoretic images.]
EGA~I \cite{EGA1}, \S9.5, Propositions~9.5.1 and~9.5.4,
Corollary~9.5.2, Definition~9.5.3, printed pp.~176--177;
the predecessor records PDF leaves~175--176. This locator
is carried with its predecessor verification provenance,
not relabeled as a new facsimile check.

\item[Analytic tensor norms.]
The \emph{R\'esum\'e} \cite{Resume} is published-catalogue
item CSG17. The predecessor's inspected USP facsimile has
\S1 no.~1, Definition~1 and Theorem~1 at printed p.~8,
leaf~8; Theorem~2 and its duality corollary at p.~$8'$,
leaf~9. Its fundamental theorem, \S4 no.~2,
Theorems~1--2, is at pp.~59--60, leaves~60--61;
Corollary~1 after Theorem~3, inequality~(8), at p.~61,
leaf~62; the proof opening in \S4 no.~3 at p.~62,
leaf~63. These inherited locators are pertinent where
norms and duality are used, and are not a source for the
norm-free border-rank proof.

\item[\'Etale and weight terminology.]
SGA~4 \cite{SGA4}, Expos\'e~VII, \S1.2, supplies the
\'etale-site foundation. Deligne \cite[Definition~1.2.2,
p.~153]{WeilII} supplies purity terminology; the predecessor
records IAS PDF leaf~17. These source checks retain their
predecessor provenance. No SGA~5 theorem or Weil~II weight
theorem is used to deduce a tensor-rank bound here.
\end{description}

The primary CSG published-works catalogue \cite{CSGWorks}
was consulted in this run. CSG17 identifies the published
\emph{R\'esum\'e}, with its S\~ao Paulo volume and USP reprint;
CSG44 identifies the published Hilbert-schemes expos\'e and
its erratum. These identifiers are not Montpellier
unpublished-manuscript cotes. The Montpellier archive
\cite{MontpellierArchive} was consulted, but no manuscript
cote and scan carrying the modern findings below was
established. For every such finding the unpublished-manuscript
entry is therefore ``not established''; the exact published
foundations are the locators above. No cote or scan is invented.

\section{Dictionary for the residual-rank findings}

\paragraph{P-01: coordinate minima.}
The tensor factors are locally free sheaves over
$\operatorname{Spec}\mathbb Z$. The row/column rescalings
are characters of a split torus in the SGA~3 group schemes.
The residual image is a free submodule with its integral
coordinate basis. A $j$-plane is a point of the relative
Grassmannian. The echelon frames in P-01 give a rank-$j$
subbundle on $\mathbb A^1_k$: the pivot submatrix remains
the identity. Its tautological tensor specializes under
the product group action. The closed fiberwise Segre
secant cone contains the specialization whenever it
contains the nonzero-parameter fibers. This is exactly
the algebraic content of the proof; no metric is present.
The framework locators are SGA~3 Expos\'e~I \S\S4.5--4.6
and EGA~I \S9.5.

\paragraph{P-02: adjacent pairs.}
The maps identifying two distinct selected coordinate
sets are morphisms of free $\mathbb Z$-modules. Their
restrictions to each selected term are injective, and
the induced tensor maps send the residual generators
to the displayed adjacent generators with unit signs.
The Segre atom parametrization is functorial under these
maps. Hence the map sends every secant cone into the
corresponding target secant cone after each base change.
This is the sheaf/group-scheme restatement of the full
projection proof, with the same EGA~I and SGA~3 locators.
For $b>2$ this dictionary retains the exact adjacent-pair
reduction; it does not assign its unknown evaluated value.

\paragraph{P-03: stars.}
For $k\geq1$, the star certificate is an integral family
over $\mathbb A^1$: its displayed Laurent expression
cancels to a polynomial tensor. Over $\mathbb G_m$ it
is a sum of $k+1$ Segre sections. Its specialization
therefore lies in the fiberwise $(k+1)$-secant cone.
The lower bound is a unit minor of a locally free-sheaf
flattening. Both constructions commute with scalar
extension and survive every characteristic.
EGA~I \S9.5 supplies the image/closure language;
SGA~3 supplies the vector-bundle representation language.

\paragraph{P-04 and P-05: the exact pencil minima and Koszul certificate.}
The $p=1$ pencil map is the bundle morphism
\[
 A\otimes B^*\longrightarrow\Lambda^2 A\otimes C,
 \qquad a\otimes\beta\longmapsto
       \sum_h a_h\wedge a\otimes Q_h(\beta).
\]
Its matrix has integral entries. For the coordinate
terms its rank is the dimension of a union of coordinate
subbundles, certified by unit minors. A rank-one Segre
section has map rank one, so the relevant minors vanish
on the claimed lower secant cones. The upper family is
the sum of the explicit star families. This proves the
same exact values in the bundle language. SGA~6
Expos\'e~V, Examples~2.2(b) and Lemma~2.2.1, supplies
the exterior-power construction; the actual minors,
rather than only its $K_0$ class, supply the bound.

\paragraph{P-06: anchored shapes.}
The $h+r+s+1$ displayed vectors are sections of a free
module on $\mathbb A^1$. Their tensor squares are
sections of the complete Veronese cone in invariant
tensors. The four-term finite difference cancels its
constant and linear terms; after dividing by the
displayed power of the parameter, it extends regularly
to zero with the asserted residual tensor as value.
Grouping identical square sections gives the stated
number of Segre atoms on $\mathbb G_m$. The lower bound
is the common-kernel flattening and its unit minors.
This is an explicit secant degeneration in the EGA~I
image language. It requires no unproved smoothability
claim about an auxiliary finite scheme.

\paragraph{P-07: determinant compression.}
The row determinant image is the integral exterior-power
representation; the polynomial determinant image uses
the tensor product of the row and column exterior powers.
The lowering transvections are root subgroups of the
split Borel group scheme in the corresponding SGA~3
general linear group. Each two-element chain gives a
regular degeneration of Grassmannian frames, and the
integer descent terminates at a Borel-stable coordinate
subbundle. The proof uses the exterior-power cancellation
law, valid even in characteristic two, not semisimplicity.
SGA~6 Expos\'e~V, Examples~2.2(b), gives the relevant
exterior-power operation. The same unipotent stability
is not asserted for the permanent image in general.

\paragraph{P-08: finite Boolean differences.}
The vector $v(u,w,y,z)$ is a section of the row--column
tensor representation over the integral affine parameter
space. Its tensor square is a section of the complete
Veronese cone. Inclusion--exclusion on Boolean evaluation
sets extracts a squarefree coefficient after cancellation
of the indicated parameter powers. Grouping evaluations
gives a sum of eleven or twelve Segre sections over
$\mathbb G_m$ and a regular limiting section over
$\mathbb A^1$. This carries the full upper-bound proof
into EGA~I's image/closure framework and SGA~3's module
representation framework. It needs no division by an
integer and no smoothability assumption on an additional
scheme. Its conjunction with R-04 uses the latter's
properness and determinantal argument.

\paragraph{P-09, P-12, P-13, P-14 and P-16: seven-plane prolongation obstructions.}
The proper Grassmann incidence and exterior bundle maps are those of
R-04, with the exact EGA~II and SGA~6 locators recorded above.
A split torus (SGA~3, Expos\'e~I, \S\S4.5--4.6) acts on the
quotient of the tautological bundle. The finite certificates cover
all possible quotient weight supports, and then prove emptiness
of specified projective common-zero loci of linear maps.
P-14 and P-16 also enlarge the fixed subbundle before testing a
remaining variable line; the relevant annihilator is taken after
this enlargement. The map entries, their annihilators and their
minors carry the proof. Their source and target classes in $K_0$
alone do not. The modular and rational certificates have the
stated characteristic-zero scope. No unpublished-manuscript cote
for these modern computations has been established.

\paragraph{P-10 and P-11: integral curves and the row cubic seed.}
P-10's inclusion--exclusion expression is a sum of Segre sections over
$\mathbb G_m$ whose grouped sum extends over $\mathbb A^1$, with
integer ordered coefficients. P-11's six- and eight-term upper
curves are integral, without division by a nonunit integer.
Its double-Koszul map is an exterior-power morphism of locally
free sheaves. The determinant lower bound uses a proper product
of three Grassmannians, the closed multiplication-rank condition,
and a fixed-point argument for a triangular group. The stable
annihilator subbundles reduce to eight explicit triples whose
actual product-map minors are units times powers of two. This
explains exactly the exclusion of characteristic two in that
lower bound; there the determinant and permanent tensors coincide.
The foundations are EGA~I \S9.5, EGA~II \S5.4, SGA~3 Expos\'e~I
\S\S4.5--4.6, SGA~6 Expos\'e~V Definition~2.1, Examples~2.2(b)
and Lemma~2.2.1, and the FGA Grassmann/Hilbert/Quot framework at
the precise printed and facsimile locators above. These foundations
are not citations for the modern numerical values themselves.

\paragraph{P-15: exact parameter arrangements.}
After torus specialization, a quotient line is a point of a
projective weight space. Specifying the zero columns of an
induced morphism gives a locally closed stratum: a linear kernel
intersection with finitely many proper linear subspaces removed.
Over an infinite field this stratum is nonempty precisely when
none of those removed subspaces is the whole intersection.
Two-dimensional quotient subspaces in a three-dimensional weight
space are described by their annihilator lines in the dual
projective plane. The certificates exhaust these linear
arrangements, including compatibility after one quotient line
has been fixed. Thus the proof is an explicit emptiness proof
on projective Grassmann strata, using the same properness and
locally free morphisms as R-04. It does not replace a parameter
space by numerical samples or by a $K_0$ class.

\paragraph{P-17: Borel-stable extension flags.}
The necessary extension spaces form a closed locus in a proper
Grassmannian fiber. Its solvable-group fixed point is refined by
an invariant complete flag in the quotient of the tautological
bundle. Each flag step is a character line annihilated by the
adjacent root operators modulo the preceding subbundle. The
certificate computes these common kernels exactly and retains
any projective family as an unresolved parameter space. Thus
EGA~II properness, the SGA~3 integral module/group schemes and
FGA flag parameter spaces carry the proof. The eighteen named
exclusions follow from actual exterior-map minors; they are
not a consequence of a $K_0$ class or a claim about every
remaining ten-plane.

\paragraph{P-18: the polynomial cubic seed.}
The linear maps $\Phi$ and $\Psi$ between free tensor modules
transfer upper bounds through preimages of closed secant cones.
A $63$-minor of determinant $-16$ proves that the nine-labelled-square
map of affine schemes is dominant after inverting two: its
invertible linear term admits a formal inverse degree by degree.
The tangent curves then give the determinant upper bound twenty-seven;
characteristic two uses the integral permanent curve. The lower
map combines two exterior operations on locally free sheaves,
with explicit integer minors in its torus weight blocks.
These use the EGA~I closure, SGA~3 module-scheme and SGA~6 exterior
power locators already specified. The field conditions come from
the actual determinants and integral identities, not a categorical
dimension count or an assumed smoothability statement.

\paragraph{P-20: all row permanent quadratic minima.}
A graph containing a full star supplies an integral family of
$j+1$ labelled square tensors. First and mixed second differences
produce its quadratic edge inclusion tensor in the special fiber.
The lower bound is the proper substitution of R-01, using the
absence of a nonzero rank-one symmetric zero-diagonal matrix
over every field. The product of two such integral curves gives
the eighteen-term polynomial ten-plane upper bound. These are
explicit sections over $\mathbb G_m$ whose sum extends over
$\mathbb A^1$; the EGA image/closure and SGA~3 module-scheme
foundations apply after every base change. No additional
archival or motivic rank theorem is asserted.

\paragraph{P-19: a closed higher-prolongation condition.}
On the Grassmannian of candidate subbundles $E\subset A\otimes B$,
the condition
\[
 \dim\bigl((A\otimes E\otimes B)
       \cap(\Sym^2 A\otimes\Sym^2 B)\bigr)\geq r
\]
is a determinantal condition for a morphism out of a
constant-rank tautological bundle. The construction is made
in the fixed ambient bundle, so it remains closed even when
an intermediate first-prolongation kernel jumps in dimension.
Proper projection extends necessity from spans of independent
Segre points to their border limits. The successive kernel
calculations are then an exact way to evaluate this fixed
closed condition. This is EGA~II properness and the FGA
parameter-space framework applied to actual bundle morphisms;
the twenty-four final unit minors do not, by themselves,
change the field scope of the earlier rational arrangements.

\paragraph{R-01: multidimensional proper substitution.}
The image map is a regular map to a Grassmannian on
the open locus of constant flattening rank. A retained
$j$-plane is a point of its relative Grassmann bundle.
The low-border-rank condition on its tautological
tensor is closed in local frames and invariant under
changing frames. Proper projection of this closed
incidence gives the closed condition needed to pass
from actual decompositions to border decompositions.
The atom-span dimension argument is unchanged on
every algebraically closed fiber. These are precisely
EGA~II Definition~5.4.1, Proposition~5.4.2 and
Theorem~5.5.3(i), applied to the Grassmannian incidence;
FGA Theorems~3.1--3.2 provide the broader parameter-space
setting if limiting schemes are also retained.

\paragraph{R-02 and R-03: the evaluated higher-dimensional minima.}
The finite coordinate or Borel-ideal reductions parametrize
special points of the same relative Grassmannians.
Their ordinary and Koszul flattenings are explicit
morphisms of locally free sheaves. Nonzero minors give
the lower certificates; the anchored families give
the upper witnesses. This is the exact transfer of
the $j=3,4$ minima.
The SGA~3 actions, SGA~6 exterior powers, and EGA~II
closed-incidence locators above are the foundations.
Each certificate retains its stated characteristic
range. A determinant nonzero modulo one prime proves
nonzero over $\mathbb Q$, but does not by itself prove
nonzero in every characteristic; a unit minor does.

\paragraph{R-04: prolongation and the exceptional permanent subspace.}
The auxiliary limiting ten-plane is retained as a
Grassmannian point, together with the prescribed
subspace it contains. Its prolongation maps are
multiplication/intersection maps of tautological
bundles, so their rank conditions are determinantal
closed conditions. This permits passage to a limiting
span before applying the weight-block exclusions.
The actual exclusion matrices retain their entries
and minors; the recorded unit minors survive every
scalar extension. Properness is the EGA~II input.
FGA's Hilbert/Quot language can retain the limiting
configuration, but the proof does not identify its
actual span with the retained limiting plane when
an evaluation rank drops. The role of the embedding
and the actual multiplication maps is essential.

\paragraph{R-05: the fifth minimum and the bound for both polynomials.}
This deduction combines R-01's proper image condition,
the finite lower certificates including R-04, and
P-06's explicit upper family. In algebraic-geometric
terms it proves avoidance of the relevant lower
Grassmann secant locus, followed by avoidance of the
lower original Segre secant cone. It is asserted in
every characteristic: the permanent minors are units,
the determinant minors are units times powers of two,
and the coefficient tensors coincide in characteristic
two. It does not identify all
special fibers of an integral orbit or secant model
merely by reducing its characteristic-zero equations.

\paragraph{R-06: uniformity in the matrix size.}
Coordinate support submodules and their split inclusions
and projections are locally free-sheaf maps over
$\mathbb Z$. The degree-two support argument reduces
the relevant geometric points to a $4\times4$ coordinate
submodule or to five finite $5\times2$ graph orbits.
The ten further Koszul minors are integral determinantal
certificates. Their characteristic analysis is identical
to R-05, and the anchored upper curves are unchanged
under the split inclusions. The row assertion applies
the same proper Grassmann incidence inside the chosen
tautological residual subbundle and uses an integral
star curve. This is the proof in the EGA~II, SGA~3 and
SGA~6 frameworks with the exact locators already listed.

\paragraph{R-07: propagation along a flag.}
The pairs $V\subset W\subset U$ with dimensions $i,j$
form a relative flag variety, obtained by composing
Grassmann bundles. The proper-incidence argument for
R-01 applies on each fiber over $W$ and compares the
two tautological residual tensors. Taking minima gives
the monotonicity inequality; the complete anchored
$2\times3$ curve supplies the sixth minimum. Properness
is exactly EGA~II Proposition~5.4.2 and
Theorem~5.5.3(i). No new categorical invariant or
cohomological implication is needed for this deduction.

\paragraph{R-08, R-09 and R-10: complete covers and the bound forty-three.}
For the determinant the torus/root-group specialization of
P-07 reduces to a finite order-ideal cover. For the permanent,
row and column permutations and transposition supply a finite
symmetry cover after P-01; no permanent Borel invariance is
assumed. The independent certificates prove that every required
coordinate point belongs to a checked case. Exterior bundle
morphisms give the Koszul determinantal conditions, and proper
Grassmann incidences give the remaining prolongation conditions.
The anchored integral curves attain the minima. R-01 and its
flag version carry these minima back to the original tensor.
The foundational locators are the EGA~II, SGA~3, SGA~6 and FGA
locators already specified. The finite covers and their rank
conclusions are findings of this run, not theorems attributed
to Grothendieck's texts or to an unverified archive cote.

\paragraph{R-11: the uniform support-graph cover.}
Coordinate inclusions and split projections carry a residual
subbundle between ambient matrix sizes without changing border
rank. The support graph records only which free-module basis
sections occur. P-01 reduces to these coordinate subbundles;
the finite graph-and-subset cover exhausts them, and exterior
bundle morphisms with checked minors give the new rank bounds.
The underlying schemes and equivariant maps are those already
located in EGA~II, SGA~3 and SGA~6. The number of graph types is
a finite verified combinatorial calculation, not an invariant
recoverable from an unembedded motive or weight-zero stalk.

\paragraph{FL-01--FL-05: the all-field F-23 theorem.}
FL-01 defines integral tensor sections and makes explicit
the distinction between invariant tensors/divided powers
and the ordinary symmetric quotient. FL-02 is the
proper Grassmannian incidence above with $j=1$.
FL-03 identifies free residual image submodules and
split coordinate projections; unit coefficients
make the image bases survive every base change.
FL-04 is a unit-minor seed calculation. FL-05 is
the factorial obstruction in the ordinary symmetric
quotient. Together these are the complete dictionary
carrying the proofs, not an assumption that a
characteristic-zero normalization exists in every
fiber. Their precise published foundations are
EGA~II's properness statements, SGA~3's module and
group schemes, and SGA~6's locally free operations.

\paragraph{FL-06: the survey counterexample.}
The degree-two Veronese over $\mathbb C$ maps a vector
to its symmetric square. Its secant lower bound is
the rank of the first flattening, a determinantal
closed condition; the four-square expression is an
explicit point of the fourth secant parametrization.
The same proof is therefore an EGA~I image/closure
calculation on a locally free representation. No
additional archive or motivic input is used.

\section{The analytic findings and their precise boundary}

\paragraph{S-01: the exact DPS example.}
The underlying tensor and its three-term decomposition
are maps of free $\mathbb Z$-modules. Its residual
module is a rank-two direct summand of a free rank-six
module, hence a $\mathbb Z$-point of
$\operatorname{Gr}(2,6)$. The lower flattening bound
is an integral minor. These are the SGA~3 and EGA~I
parts. The projector and dual-witness matrices have
entries in $\mathbb Z[1/2]$, but their positive
semidefiniteness is a Hermitian statement on a chosen
complex fiber, certified by the displayed Gram
factorizations. The real inequality is not an
identity of algebraic $K_0$ classes.

\paragraph{S-02: the error estimate.}
The extension space uses the Schur functor
$\operatorname{Sym}^k B$ on a locally free sheaf.
Its rank is $\binom{k+\dim B-1}{\dim B-1}$, and its
$K_0$ generating operation is
$\sigma_t([B])=\lambda_{-t}([B])^{-1}$.
SGA~6 Expos\'e~V, Definition~2.1, Examples~2.2(b)
and Lemma~2.2.1, give this exterior/symmetric-operation
setting. The partial-transpose and positive-cone
constraints require a metric and a complex
archimedean fiber. The proved error bound is carried
by the analytic separability argument, not by a
gamma-filtration or a weight argument.

\paragraph{O-01: the orbit supremum.}
The product-$\operatorname{GL}$ group scheme acts on
the relative Grassmannian and on the Segre scheme.
The closed incidence
$\{U:\mathbb P(U)\cap\operatorname{Segre}\ne\varnothing\}$
is GCT\_GINEQ F-18's $\mathcal I_{q,1}$, obtained by
the EGA~II proper projection. The explicit
one-parameter family carries an accompanying line
to a Segre point; properness of the Grassmannian
and closed incidence retain that line in a limiting
subspace. The supremum-one conclusion is its
archimedean metric interpretation. For a product-free
space, the actual orbit misses this incidence while
its closure meets it. The diagonal action on an
alternating image is a different group action and
has the separate behavior stated in O-01.

\paragraph{O-02: the Pl\"ucker moment equation.}
The Pl\"ucker vector generates the determinant line
$\det(U)=\Lambda^q U$ of the tautological bundle,
inside $\Lambda^q(\bigotimes_i H_i)$. This exterior
power and group action are integral constructions
in SGA~3 and SGA~6. The Hermitian norm, partial
traces, and convexity along Hermitian one-parameter
subgroups are additional analytic structures.
The displayed scalar-marginal equation is the
critical-point condition for that norm, and its
proof establishes minimality. It is not an equality
of $K_0$ classes and is not the orbit-alpha supremum.

\section{What the categorical invariants do and do not retain}

For every Koszul map in this paper, $K_0$ records
the source and target classes, for example
$[B^*]\lambda^p[A]$ and $[C]\lambda^{p+1}[A]$.
For representations one uses the exact tensor
category of finite locally free representations;
SGA~6's representation-ring example is the relevant
locator. Proposition~3.2(ii) and equation~(3.10.1)
then describe the gamma operations and filtration.
Integral Schur functors can be used over
$\mathbb Z$; characteristic-zero Young projectors
requiring division by factorials or hook lengths
cannot be transferred without localization.

These constructions track modules and characters.
The rank of a particular morphism is supplied by
its minors, not by its source and target classes.
GCT\_GINEQ F-20 gives the precise counterexample:
the zero and identity maps have the same such
classes and different ranks. Hilbert and Quot
schemes retain additional geometric data, but
one must keep the actual embedding and limiting
span when those enter the certificate.

No new $\ell$-adic sheaf, purity theorem, or motive
is constructed by the proofs. GCT\_GINEQ F-21
already shows that the constant weight-zero stalk
and absolute point motive cannot distinguish a
zero tensor point from a nonzero pure tensor point.
This does not rule out embedding-sensitive sheaves
or relative motives as future invariants. It
identifies exactly why the present rank certificates
do not follow merely from attaching these names.
There is consequently no fabricated SGA~5 or
unpublished-manuscript theorem attached to such an
inference.


\begin{thebibliography}{99}
% Insert these items inside the final thebibliography environment.
% Keys EGA1,EGA2,SGA3,SGA6,FGA,SGA4,WeilII follow predecessor usage.
% All other keys are new for these fragments; merge duplicates if needed.

\bibitem{LT2010}
J.~M. Landsberg and Z. Teitler,
\emph{On the ranks and border ranks of symmetric tensors},
Found. Comput. Math. 10 (2010), 339--366.
Version inspected: arXiv:0901.0487v3 (28 September 2009),
Corollary~1.4, p.~2; \S9 and Table~3, pp.~13--14.
\url{https://arxiv.org/pdf/0901.0487}.

\bibitem{Farnsworth2016}
C. Farnsworth,
\emph{Koszul--Young flattenings and symmetric border rank of the determinant},
J. Algebra 447 (2016), 664--676.
Version inspected: arXiv:1505.05079v2 (11 November 2015),
Definition~1.2 and Theorems~1.3, 1.6, 1.8, pp.~1--2.
\url{https://arxiv.org/pdf/1505.05079}.

\bibitem{LM2018}
J.~M. Landsberg and M. Micha\l ek,
\emph{A $2n^2-\log_2(n)-1$ lower bound for the border rank of matrix multiplication},
Int. Math. Res. Not. 2018(15), 4722--4733.
Version inspected: arXiv:1608.07486v1 (26 August 2016),
Proposition~2.3, p.~3 (all retained dimensions).
\url{https://arxiv.org/pdf/1608.07486}.

\bibitem{CHL2023}
A. Conner, A. Harper and J.~M. Landsberg,
\emph{New lower bounds for matrix multiplication and $\operatorname{det}_3$},
Forum Math. Pi 11 (2023), e17.
Theorem~1.2, \S1.2; Theorems~2.3--2.4 and \S2.3.
\url{https://doi.org/10.1017/fmp.2023.14}.
Preprint: arXiv:1911.07981.

\bibitem{CHHL2023}
A. Conner, H. Huang and J.~M. Landsberg,
\emph{Bad and good news for Strassen's laser method:
Border rank of $\mathrm{Perm}_3$ and strict submultiplicativity},
Found. Comput. Math. 23 (2023), 2049--2087.
Version inspected: arXiv:2009.11391v1 (23 September 2020),
Theorem~1.1 and following identification, p.~3; proof in \S5.
\url{https://arxiv.org/pdf/2009.11391}.

\bibitem{BB2021}
W. Buczy\'nska and J. Buczy\'nski,
\emph{Apolarity, border rank, and multigraded Hilbert scheme},
Duke Math. J. 170(16) (2021), 3659--3702.
Version inspected: arXiv:1910.01944v2 (7 November 2020),
Theorem~3.15, p.~15; Theorem~4.3, p.~18; Corollary~4.5,
p.~19; \S7.2, p.~32.
\url{https://arxiv.org/pdf/1910.01944}.

\bibitem{BB2026}
W. Buczy\'nska and J. Buczy\'nski,
\emph{Apolarity for border cactus decompositions},
arXiv:2601.19558v1 (27 January 2026),
field convention at the start of \S1; \S4.3, Theorem~4.7.
\url{https://arxiv.org/html/2601.19558v1}.

\bibitem{Mandziuk2026}
T. Ma\'ndziuk,
\emph{Border rank lower bounds beyond weak border apolarity},
arXiv:2609.12121v1 (10 September 2026),
Theorems~1.1--1.3, p.~2; Theorems~3.5, 4.4, 5.2; \S6.
\url{https://arxiv.org/pdf/2609.12121}.

\bibitem{DM2026}
M. Dole\v z\'alek and M. Micha\l ek,
\emph{Nonlinear methods for tensors: determinantal equations
for secant varieties beyond cactus},
arXiv:2602.12762v1 (13 February 2026),
Theorem~1.1; \S5.1; \S6.3.
\url{https://arxiv.org/html/2602.12762v1}.

\bibitem{HJK2026}
Jong In Han, Jeong-Hoon Ju and Yeongrak Kim,
\emph{The border rank of the $4\times4$ determinant tensor is twelve},
arXiv:2510.11051v2 (18 August 2026),
Theorem~1.1 and \S3. The tensor is the row determinant.
\url{https://arxiv.org/html/2510.11051v2}.

\bibitem{HJK2025}
Jong In Han, Jeong-Hoon Ju and Yeongrak Kim,
\emph{Recursive Koszul flattenings of determinant and permanent tensors},
arXiv:2503.12032v1 (15 March 2025),
Theorems~4.7, 5.6--5.7. These are row tensors.
\url{https://arxiv.org/html/2503.12032v1}.

\bibitem{HS2026}
Jong In Han and Jeyoung Song,
\emph{The rank of the $5\times5$ permanent tensor is sixteen},
arXiv:2608.09708v2 (27 August 2026),
Theorem~A, p.~3; Proposition~4.1, p.~10.
These are row tensors; Theorem~A concerns ordinary CP rank.
\url{https://arxiv.org/pdf/2608.09708v2}.

\bibitem{Han2026}
Jong In Han,
\emph{The border Waring rank of $x_1\cdots x_n$ is $2^{n-1}$},
arXiv:2609.09141v1 (8 September 2026),
Theorem~A, p.~2, and the preceding ambient-space identification.
\url{https://arxiv.org/pdf/2609.09141v1}.

\bibitem{BT2020}
M. Boij and Z. Teitler,
\emph{A bound for the Waring rank of the determinant via syzygies},
arXiv:1908.08896v2 (8 November 2019),
background comparisons, pp.~5--7; \S4, p.~15.
\url{https://arxiv.org/pdf/1908.08896}.

\bibitem{HGJ2024}
R. Houston, A.~P. Goucher and N. Johnston,
\emph{A new formula for the determinant and bounds on its tensor and Waring ranks},
Combin. Probab. Comput. 33(6) (2024).
Primary published text checked for the scope of the rank results:
\url{https://www.cambridge.org/core/journals/combinatorics-probability-and-computing/article/new-formula-for-the-determinant-and-bounds-on-its-tensor-and-waring-ranks/E33F5E0726A0691B250C3BB6CE3816F5}.

\bibitem{DL2025}
P. Dutta and V. Lysikov,
\emph{Recent Advances in Debordering Methods},
arXiv:2510.13049v1 (15 October 2025),
\S4.3, p.~29, immediately after Corollary~56.
The displayed universal lower bound is checked and refuted as
written in finding FL-06.
\url{https://arxiv.org/html/2510.13049v1}.

\bibitem{EGA1}
A. Grothendieck and J. Dieudonn\'e,
\emph{\'El\'ements de g\'eom\'etrie alg\'ebrique I},
Publ. Math. IH\'ES 4 (1960), \S9.5, pp.~176--177.
Propositions~9.5.1 and~9.5.4, Corollary~9.5.2, Definition~9.5.3.
NUMDAM \texttt{PMIHES\_1960\_\_4\_\_5\_0}.
\url{https://www.numdam.org/article/PMIHES_1960__4__5_0.pdf}.

\bibitem{EGA2}
A. Grothendieck, with J. Dieudonn\'e,
\emph{\'El\'ements de g\'eom\'etrie alg\'ebrique II},
Publ. Math. IH\'ES 8 (1961), 5--222.
Definition~5.4.1 and Proposition~5.4.2, p.~100,
PDF leaf~97; Theorem~5.5.3(i), p.~104, leaf~101.
\url{https://www.numdam.org/article/PMIHES_1961__8__5_0.pdf}.

\bibitem{SGA3}
M. Demazure and A. Grothendieck et al.,
\emph{Sch\'emas en groupes}, SGA~3,
Expos\'e~I, \S\S4.5--4.6, p.~29,
PDF leaf~42 in the linked facsimile.
\url{https://library.slmath.org/nonmsri/sga/sga/pdf/sga3-1.pdf}.

\bibitem{SGA6}
P. Berthelot, A. Grothendieck, L. Illusie et al.,
\emph{Th\'eorie des intersections et th\'eor\`eme de Riemann--Roch},
SGA~6, Expos\'e~V. Definition~2.1, Examples~2.2(b),
Lemma~2.2.1, pp.~314--316, PDF leaves~321--323;
representation-ring example, p.~317, leaf~324;
Proposition~3.2(ii), p.~325, leaf~332;
\S3.10, equation~(3.10.1), p.~331, leaf~338.
\url{https://library.slmath.org/nonmsri/sga/sga/pdf/sga6.pdf}.

\bibitem{FGA}
A. Grothendieck,
\emph{Techniques de construction et th\'eor\`emes d'existence
en g\'eom\'etrie alg\'ebrique IV: les sch\'emas de Hilbert},
S\'eminaire Bourbaki, expos\'e~221 (1960--1961), 249--276.
Theorems~3.1--3.2, pp.~259--260, PDF leaves~12--13;
after Proposition~3.8 and Remarks~3.9, pp.~265--266,
leaves~18--19. Published-catalogue identifier CSG44.
\url{https://www.numdam.org/article/SB_1960-1961__6__249_0.pdf}.

\bibitem{Resume}
A. Grothendieck,
\emph{R\'esum\'e de la th\'eorie m\'etrique des produits tensoriels topologiques},
Boletim Soc. Mat. S\~ao Paulo 8, 1--79 (dated 1953
in the CSG published-works catalogue, item CSG17).
Facsimile reprint: Resenhas IME-USP 2(4) (1996), 401--480.
\url{https://doi.org/10.11606/resimeusp.v2i4.74836}.
The specific inherited section and leaf locators are given in
\S\ref{sec:groth-locators}; CSG17 is not a manuscript cote.

\bibitem{SGA4}
M. Artin, A. Grothendieck, J.-L. Verdier et al.,
\emph{Th\'eorie des topos et cohomologie \'etale des sch\'emas},
SGA~4, Expos\'e~VII, \S1.2 (the \'etale site).
No tensor-rank implication is attributed to this construction.

\bibitem{WeilII}
P. Deligne,
\emph{La conjecture de Weil II},
Publ. Math. IH\'ES 52 (1980), 137--252.
Definition~1.2.2, p.~153; predecessor's IAS facsimile leaf~17.
\url{https://publications.ias.edu/sites/default/files/Number40.pdf}.

\bibitem{CSGWorks}
CSG,
\emph{Published Works}, entries CSG17 and CSG44,
consulted 21 November 2026.
\url{https://csg.igrothendieck.org/published-works/}.
These are published-work catalogue identifiers.

\bibitem{MontpellierArchive}
Universit\'e de Montpellier,
\emph{Archives Grothendieck}, catalogue consulted
21 November 2026.
\url{https://grothendieck.umontpellier.fr/archives-grothendieck/}.
No unpublished manuscript cote/scan supporting the new findings
was established; the paper uses exact published foundations instead.

% Primary versions and locators verified by /root/sdp_orbits;
% metadata cross-checked by /root/fields_literature.
% Insert inside thebibliography; alias citation keys if parent uses others.

\bibitem{DPS}
A.~C. Doherty, P.~A. Parrilo and F.~M. Spedalieri,
\emph{A complete family of separability criteria},
Phys. Rev. A 69 (2004), 022308.
Version inspected: arXiv:quant-ph/0308032v3 (5 November 2003),
Theorem~1, p.~7.
\url{https://arxiv.org/pdf/quant-ph/0308032v3}.
\url{https://doi.org/10.1103/PhysRevA.69.022308}.

\bibitem{NOP}
M. Navascu\'es, M. Owari and M.~B. Plenio,
\emph{The power of symmetric extensions for entanglement detection},
Phys. Rev. A 80 (2009), 052306.
Version inspected: arXiv:0906.2731v1 (15 June 2009),
Definition~1, p.~2; Theorem~2, equation~(7), p.~3;
Corollary~5, equation~(16), p.~4.
\url{https://arxiv.org/pdf/0906.2731v1}.
\url{https://doi.org/10.1103/PhysRevA.80.052306}.

\bibitem{BGOWW}
P. B\"urgisser, A. Garg, R. Oliveira, M. Walter and A. Wigderson,
\emph{Alternating minimization, scaling algorithms, and the null-cone
problem from invariant theory},
ITCS 2018, LIPIcs 94, article~24, pp.~24:1--24:20.
Full version inspected: arXiv:1711.08039v1 (21 November 2017),
equation~(1.3), p.~4; Theorem~3.2, p.~13.
These locators are for the full arXiv version.
\url{https://arxiv.org/pdf/1711.08039v1}.
\url{https://doi.org/10.4230/LIPIcs.ITCS.2018.24}.

\bibitem{JKK2025}
Jeong-Hoon Ju, Taehyeong Kim and Yeongrak Kim,
\emph{Finding tensor decompositions with sparse optimization},
J. Korean Math. Soc. 62(1) (2025), 33--49.
Theorem~4.1, equation~(15), printed p.~46.
\url{https://doi.org/10.4134/JKMS.j230579}.
Primary published full text:
\url{https://pdf.medrang.co.kr/JKMS/2025/062/JKMS-62-1-33-49.html}.
Also inspected: arXiv:2305.13964v1 (23 May 2023),
Theorem~4.1, equation~(4.8), p.~16;
the author order in that preprint is T.~Kim, J.-H.~Ju, Y.~Kim.
\url{https://arxiv.org/pdf/2305.13964v1}.

\end{thebibliography}

\end{document}
